\documentclass[12pt,a4paper]{article}
\usepackage{fontspec}
\usepackage{geometry}
\usepackage{enumitem}
\usepackage{fancyhdr}
\usepackage{lastpage}
\usepackage{hyperref}

% Configure IBM Plex Sans font family
\setmainfont{IBMPlexSansDevanagari}[
  Path=/home/asi0/asi0-repos/bitcoin-educational-content/scripts/quizz_to_pdfs/fonts/TrueType/IBM-Plex-Sans-Devanagari/,
  Extension=.ttf,
  UprightFont=*-Regular,
  BoldFont=*-Bold,
  Scale=1.0
]

\geometry{margin=2.5cm}
\pagestyle{fancy}
\fancyhf{}
\rhead{Page \thepage\ of \pageref{LastPage}}
\lhead{BTC101 - Quiz (hi)}
\renewcommand{\headrulewidth}{0.4pt}

\title{Quiz: BTC101}
\author{Language: hi}
\date{\today}

\begin{document}

\maketitle
\thispagestyle{fancy}

\section*{Instructions}
Answer all questions by selecting the best option (A, B, C, or D). The answer key is provided at the end of this document.

\bigskip


\subsection*{Question 1}
block size war कब हुआ था?

\begin{enumerate}[label=\Alph*., leftmargin=*]
\item २०१९ में।
\item 2008 में।
\item 2017 में।
\item २०१० में।
\end{enumerate}

\bigskip


\subsection*{Question 2}
Bitcoin में हम फीस क्यों देते हैं?

\begin{enumerate}[label=\Alph*., leftmargin=*]
\item अमीर लोगों को ज्यादा भुगतान करने देना।
\item Bitcoin को माइन कैसे करें, इसका परिचय देने के लिए।
\item ब्लॉक में लेन-देन को शामिल करने के लिए एक मुक्त बाजार बनाना।
\item सरकारों को नेटवर्क से बाहर रखना।
\end{enumerate}

\bigskip


\subsection*{Question 3}
Bitcoin की दो मुख्य विशेषताएं क्या हैं?

\begin{enumerate}[label=\Alph*., leftmargin=*]
\item जीडब्ल्यू-1 एक कंप्यूटर प्रोटोकॉल और साथ ही एक मुद्रा इकाई भी है।
\item जीडब्ल्यू-5 एक केंद्रीकृत और विकेंद्रीकृत दोनों तरह की मुद्रा है।
\item जीडब्ल्यू-4 एक तरह का स्टॉक भी है और एक तरह का बॉन्ड भी।
\item जीडब्ल्यू-2 एक फिजिकल करेंसी है और डिजिटल जीडब्ल्यू-3 भी।
\end{enumerate}

\bigskip


\subsection*{Question 4}
नोबेल पुरस्कार विजेता एफ. ए. हायेक के अनुसार, दोबारा अच्छा पैसा पाने के लिए क्या करना होगा?

\begin{enumerate}[label=\Alph*., leftmargin=*]
\item हम यह चीज़ एक तानाशाही सरकार को दे देते हैं।
\item हम यह चीज़ एक बेहतर सरकार को देते हैं।
\item हम इस चीज़ को सरकार के हाथ से ले लेते हैं।
\item हम लोगों के हाथ से चीज़ें छीन लेते हैं।
\end{enumerate}

\bigskip


\subsection*{Question 5}
Bitcoin को एक न्यूट्रल करेंसी कहने का क्या मतलब है?

\begin{enumerate}[label=\Alph*., leftmargin=*]
\item यह सरकार द्वारा नियंत्रित किया जाता है।
\item यह किसी भी संस्था या संगठन द्वारा नियंत्रित नहीं है।
\item यह एक ही संस्था या संगठन द्वारा नियंत्रित होता है।
\item यह सिर्फ बड़े लेन-देन के लिए इस्तेमाल होता है।
\end{enumerate}

\bigskip


\subsection*{Question 6}
साइफरपंक्स कौन हैं?

\begin{enumerate}[label=\Alph*., leftmargin=*]
\item Bitcoin और दूसरी क्रिप्टोकरेंसी बनाने वाले लोगों का एक समूह।
\item एक ऐसा गुट जो टेक्नोलॉजी के इस्तेमाल के खिलाफ था।
\item डिजिटल युग में प्राइवेसी और आज़ादी की भूमिका पर गहराई से सवाल उठाने वाले लोगों का एक समूह।
\item एक गैंग हैकर्स का जिसने सरकारों से जानकारी चुराई।
\end{enumerate}

\bigskip


\subsection*{Question 7}
Cypherpunk मेनिफेस्टो क्या है?

\begin{enumerate}[label=\Alph*., leftmargin=*]
\item एरिक ह्यूजेस द्वारा लिखा गया एक पाठ जो यह दावा करता है कि गोपनीयता एक मौलिक अधिकार है।
\item जूलियन असांजे द्वारा लिखी गई एक किताब जो टेक्नोलॉजी के भविष्य के बारे में है।
\item डेविड चॉम द्वारा लिखा गया एक घोषणापत्र जो डिजिटल मुद्रा के महत्व के बारे में है।
\item Bitcoin के सिद्धांतों को बताने वाला एक दस्तावेज़।
\end{enumerate}

\bigskip


\subsection*{Question 8}
डिजीकैश प्रोजेक्ट क्या था?

\begin{enumerate}[label=\Alph*., leftmargin=*]
\item डेविड चौम द्वारा गुमनाम इलेक्ट्रॉनिक पैसा बनाने का एक प्रयास।
\item फेडरल रिजर्व द्वारा सभी भौतिक नकदी को डिजिटाइज़ करने का एक प्रयास।
\item एक डिजिटल Wallet बनाने का प्रोजेक्ट Bitcoin के लिए।
\item Bitcoin पर आधारित एक क्रिप्टोकरेंसी बनाने का प्रोजेक्ट।
\end{enumerate}

\bigskip


\subsection*{Question 9}
साइबरस्पेस की स्वतंत्रता की घोषणा क्या है?

\begin{enumerate}[label=\Alph*., leftmargin=*]
\item एक पाठ जिसमें जोर दिया गया है कि साइबरस्पेस को सरकारी विनियमन द्वारा नियंत्रित किया जाना चाहिए।
\item एक पाठ जिसमें यह कहा गया है कि साइबरस्पेस भौतिक क्षेत्र से अलग क्षेत्र है और उस पर समान कानून लागू नहीं होने चाहिए।
\item एक घोषणा जिसमें पारंपरिक वित्तीय प्रणालियों से Bitcoin की स्वतंत्रता का उल्लेख किया गया है।
\item एक घोषणा जो समुद्री डाकू आंदोलन की किसी भी विनियमन से स्वतंत्रता की गारंटी देती है।
\end{enumerate}

\bigskip


\subsection*{Question 10}
वी डाई का बी-मनी क्या है?

\begin{enumerate}[label=\Alph*., leftmargin=*]
\item वेई दाई द्वारा बनाई गई एक डिजिटल करेंसी जो Bitcoin से ज्यादा स्केलेबल होने के लिए उसकी प्रतिस्पर्धी के तौर पर बनाई गई थी।
\item एक तरह का डिजिटल पैसा जिसे CIA ने बनाया था और सेना द्वारा इस्तेमाल किया जाता था।
\item एक डिजिटल करेंसी जो Bitcoin से पहले बहुत प्रचलित थी और साइफरपंक्स के एक समूह द्वारा नियंत्रित की जाती थी।
\item एक गुमनाम डिजिटल करेंसी का आइडिया जहां धोखाधड़ी का पता लगाने का काम किसी केंद्रीय अथॉरिटी की बजाय इवैल्यूएटर्स के एक कम्युनिटी द्वारा किया जाता हो।
\end{enumerate}

\bigskip


\subsection*{Question 11}
Bitcoin सिर्फ़ टेक्नोलॉजी से आगे क्या दर्शाता है?

\begin{enumerate}[label=\Alph*., leftmargin=*]
\item अपराधियों के लिए अवैध लेन-देन करने का एक ज़रिया।
\item पारंपरिक वित्तीय प्रणालियों के खिलाफ एक क्रांति और पारदर्शिता, विकेंद्रीकरण, तथा व्यक्तिगत संप्रभुता पर आधारित एक विकल्प।
\item एक नया डिजिटल करेंसी का रूप जो पारंपरिक मुद्राओं से ज्यादा कारगर है क्योंकि यह तेज़ है।
\item Mining के जरिए लोगों के लिए पैसे कमाने का एक तरीका।
\end{enumerate}

\bigskip


\subsection*{Question 12}
पहले-पहले किस तरह के करेंसी चलते थे?

\begin{enumerate}[label=\Alph*., leftmargin=*]
\item कागज़ के नोट।
\item अनाज और पशुधन जैसे सामान।
\item सोना और चाँदी।
\item डिजिटल संपत्ति।
\end{enumerate}

\bigskip


\subsection*{Question 13}
धन के तीन विहित कार्य क्या हैं?

\begin{enumerate}[label=\Alph*., leftmargin=*]
\item मूल्य का भण्डार, Exchange का माध्यम, लेखा की इकाई।
\item मूल्य का भण्डार, जी.डब्लू.-1 का माध्यम, उत्पादन के साधन।
\item मूल्य का भण्डार, उत्पादन के साधन, लेखांकन की इकाई।
\item जी.डब्लू.-2 का माध्यम, उत्पादन के साधन, लेखांकन की इकाई।
\end{enumerate}

\bigskip


\subsection*{Question 14}
एक प्रभावी करेंसी की मुख्य खासियतें क्या होती हैं?

\begin{enumerate}[label=\Alph*., leftmargin=*]
\item फंजिबल (बदलने योग्य), इंडिविजिबल (अविभाज्य), जीडब्ल्यू-1.
\item फंजिबल, डिविजिबल, पेरिशेबल।
\item फंजिबल (बदलने योग्य), डिविजिबल (विभाज्य), जीडब्ल्यू-0
\item गैर-फंजिबल, विभाज्य, जीडब्ल्यू-2।
\end{enumerate}

\bigskip


\subsection*{Question 15}
मुद्रा के रूप में सोने का एक नुकसान क्या है?

\begin{enumerate}[label=\Alph*., leftmargin=*]
\item यह आसानी से विभाज्य नहीं है.
\item इसका उत्पादन करना आसान है.
\item इसका मूल्य अंतरिक्ष में एक समान नहीं रहता।
\item यह समय के साथ नष्ट हो जाता है।
\end{enumerate}

\bigskip


\subsection*{Question 16}
सोना Exchange का मानक माध्यम क्यों बना?

\begin{enumerate}[label=\Alph*., leftmargin=*]
\item इसकी भरपूर मात्रा, जल्दी खराब होने की प्रवृत्ति, और अविभाज्यता के कारण।
\item इसकी दुर्लभता, टिकाऊपन और विभाज्यता के कारण।
\item इसकी दुर्लभता, नाशवानता और विभाज्यता के कारण।
\item इसकी भरपूर मात्रा, मजबूती, और अविभाज्यता के कारण।
\end{enumerate}

\bigskip


\subsection*{Question 17}
पैसा एक कम्युनिकेशन टूल के रूप में क्या भूमिका निभाता है?

\begin{enumerate}[label=\Alph*., leftmargin=*]
\item यह समय के साथ मूल्य संचार की अनुमति देता है और बिना किसी दूसरे पक्ष पर भरोसा किए।
\item यह समय और स्थान की सीमाओं के पार मूल्य संचार की अनुमति देता है और किसी दूसरे पक्ष पर भरोसा करने की जरूरत नहीं होती।
\item यह स्पेस और टाइम के पार वैल्यू कम्युनिकेशन को मुमकिन बनाता है और किसी काउंटरपार्टी पर भरोसा करने की जरूरत को खत्म करता है।
\item यह सिर्फ जगह के पार वैल्यू कम्युनिकेशन करने देता है और किसी दूसरे पार्टी पर भरोसा करने की जरूरत नहीं होती।
\end{enumerate}

\bigskip


\subsection*{Question 18}
फिएट करेंसी का सबसे बड़ा नुकसान क्या है?

\begin{enumerate}[label=\Alph*., leftmargin=*]
\item वे आसानी से ले जाने योग्य नहीं हैं।
\item उनकी कीमत हमेशा मुद्रास्फीति से बढ़ती रहती है।
\item वे विभाज्य नहीं हैं।
\item उनकी कीमत लगातार मुद्रास्फीति से कम होती जाती है।
\end{enumerate}

\bigskip


\subsection*{Question 19}
Bitcoin का दूसरे करेंसी फॉर्म्स पर मुख्य फायदा क्या है?

\begin{enumerate}[label=\Alph*., leftmargin=*]
\item यह मुद्रास्फीति से जूझ रही मुद्रा का एक बढ़िया उदाहरण है।
\item यह दुनिया के सभी व्यापारियों द्वारा व्यापार में आम तौर पर स्वीकार किया जाता है।
\item यह अपनी सख्ती से सीमित Supply और तटस्थता की वजह से मूल्य का एक उत्कृष्ट भंडार प्रदान करता है।
\item यह अविभाज्य और असंचरणीय है।
\end{enumerate}

\bigskip


\subsection*{Question 20}
सोना ग्लोबलाइज्ड और डिजिटल दुनिया में अच्छा परफॉर्म क्यों नहीं कर पाता?

\begin{enumerate}[label=\Alph*., leftmargin=*]
\item क्योंकि यह इतना कीमती नहीं है।
\item क्योंकि यह इतना दुर्लभ नहीं है।
\item क्योंकि यह समय के साथ घिस जाता है।
\item क्योंकि इसे विभाज्य और आसानी से विनिमेय बनाने के लिए एक केंद्रीय संस्था की जरूरत होती है।
\end{enumerate}

\bigskip


\subsection*{Question 21}
आजकल कॉमर्स में Bitcoin को इतना एक्सेप्ट क्यों नहीं किया जाता?

\begin{enumerate}[label=\Alph*., leftmargin=*]
\item क्योंकि यह छोटी इकाइयों में बाँटा नहीं जा सकता।
\item क्योंकि इसे ज्यादातर लोगों ने अपनाया नहीं है।
\item क्योंकि यह ले जाने योग्य नहीं है।
\item क्योंकि यह Liquid नहीं है।
\end{enumerate}

\bigskip


\subsection*{Question 22}
आज के करेंसी (मुद्रा) का विकास पथ क्या है?

\begin{enumerate}[label=\Alph*., leftmargin=*]
\item कच्चा पत्थर -> सिक्का -> नोट -> Blockchain -> Lightning Network
\item कच्चा पत्थर -> सिक्का -> नोट -> बैंक कार्ड -> जीडब्ल्यू-4
\item कच्चा पत्थर -> सिक्का -> बैंक कार्ड -> नोट -> Blockchain -> Lightning Network
\item कच्चा पत्थर -> सिक्का -> नोट -> बैंक कार्ड -> Blockchain -> Lightning Network
\end{enumerate}

\bigskip


\subsection*{Question 23}
सोना पोर्टेबल मटेरियल क्यों नहीं है?

\begin{enumerate}[label=\Alph*., leftmargin=*]
\item क्योंकि इसे आसानी से छोटे रूपों में ढाला नहीं जा सकता।
\item क्योंकि यह इतना कीमती है कि इसे हिलाया नहीं जा सकता।
\item क्योंकि यह घना और भारी होता है, जिससे बड़ी मात्रा में ले जाना मुश्किल हो जाता है।
\item क्योंकि यह कमरे के तापमान पर एक Liquid है।
\end{enumerate}

\bigskip


\subsection*{Question 24}
फिएट करेंसी का सबसे बड़ा फायदा क्या है?

\begin{enumerate}[label=\Alph*., leftmargin=*]
\item वे समय के साथ मूल्य बनाए रखने में कारगर हैं।
\item वे आसानी से ले जाने योग्य नहीं हैं।
\item वे एक केंद्रीय पार्टी द्वारा नियंत्रित होते हैं।
\item वे आसानी से बाँटे जा सकते हैं।
\end{enumerate}

\bigskip


\subsection*{Question 25}
निम्नलिखित में से कौन सा Bitcoin के मुद्रा के रूप में एक फायदा नहीं है?

\begin{enumerate}[label=\Alph*., leftmargin=*]
\item यह विभाज्य है।
\item यह बिना इजाज़त के है।
\item यह ले जाने योग्य है।
\item इसका इस्तेमाल टैक्स भरने के लिए किया जाता है।
\end{enumerate}

\bigskip


\subsection*{Question 26}
ध्वनि पैसा कैसे बनता है?

\begin{enumerate}[label=\Alph*., leftmargin=*]
\item साउंड मनी को कंपनियों द्वारा हिंसा के इस्तेमाल से थोपना पड़ता है।
\item सही पैसा बनाया नहीं जा सकता।
\item साउंड मनी (मजबूत मुद्रा) स्वाभाविक रूप से बाजार और लोगों की सहमति से ही उभरनी चाहिए।
\item सरकार को मुद्रा संबंधी कानून बनाकर साउंड मनी (मजबूत मुद्रा) को लागू करना चाहिए।
\end{enumerate}

\bigskip


\subsection*{Question 27}
Bitcoin को पैसे का नया रूप क्यों माना जाता है?

\begin{enumerate}[label=\Alph*., leftmargin=*]
\item क्योंकि यह एक डिजिटल पैसा है जिसका कोई केंद्रीय नियंत्रण नहीं है।
\item क्योंकि यह सोने से समर्थित एक डिजिटल मुद्रा है।
\item क्योंकि फेडरल रिजर्व ऐसा कहता है।
\item क्योंकि यह एक ऐसी करेंसी है जिसके लिए Blockchain बैंक की जरूरत पड़ती है।
\end{enumerate}

\bigskip


\subsection*{Question 28}
सोने के मुख्य फायदे क्या हैं मुद्रा के रूप में?

\begin{enumerate}[label=\Alph*., leftmargin=*]
\item यह आसानी से सत्यापित किया जा सकता है।
\item ये बहुत ही मजबूत है।
\item यह आसानी से बाँटा जा सकता है या लंबी दूरी तक ले जाया जा सकता है।
\item यह सरकार द्वारा नियंत्रित नहीं है।
\end{enumerate}

\bigskip


\subsection*{Question 29}
फिएट मनी का क्या अर्थ है?

\begin{enumerate}[label=\Alph*., leftmargin=*]
\item फिएट मनी एक प्रकार की क्रिप्टोकरेंसी है जो विकेन्द्रीकृत नेटवर्क पर संचालित होती है, जिससे केंद्रीय प्राधिकरण की आवश्यकता के बिना पीयर-टू-पीयर लेनदेन की अनुमति मिलती है।
\item फिएट मुद्रा वह मुद्रा है जिसका कोई आंतरिक मूल्य नहीं होता है और इसका उपयोग लेनदेन के लिए नहीं किया जा सकता है, जिससे यह अर्थव्यवस्था में अनिवार्यतः बेकार हो जाती है।
\item फिएट मुद्रा वह मुद्रा है जिसका मूल्य इसलिए होता है क्योंकि सरकार इसका रखरखाव करती है और लोगों को इसके मूल्य पर विश्वास होता है, न कि यह सोने जैसी भौतिक वस्तु द्वारा समर्थित होती है।
\item फिएट मुद्रा वह मुद्रा होती है जो किसी भौतिक वस्तु, जैसे सोना या चांदी, द्वारा समर्थित होती है, जिसका अर्थ है कि इसका मूल्य सीधे उस वस्तु के मूल्य से जुड़ा होता है।
\end{enumerate}

\bigskip


\subsection*{Question 30}
मुद्राओं के संदर्भ में "फिड्यूशियरी" शब्द का क्या मतलब है?

\begin{enumerate}[label=\Alph*., leftmargin=*]
\item यह उन मुद्राओं को संदर्भित करता है जो केवल किसी विशेष संस्थान के भीतर ही इस्तेमाल की जाती हैं।
\item यह उन मुद्राओं की बात करता है जिनकी कीमत सिर्फ इसलिए होती है क्योंकि लोग उस केंद्रीय संस्था पर भरोसा करते हैं जो इन्हें नियंत्रित करती है।
\item यह उन मुद्राओं को संदर्भित करता है जो केवल विश्वास-आधारित लेन-देन (फिड्यूशियरी ट्रांजैक्शन) के लिए मूल्यवान होती हैं।
\item यह उन मुद्राओं को संदर्भित करता है जो सोने द्वारा समर्थित हैं और छोटे वाणिज्यिक बैंकों द्वारा जारी किए जाते हैं।
\end{enumerate}

\bigskip


\subsection*{Question 31}
फ़िड्यूशियरी करेंसी जारी करने की ज़िम्मेदारी किन संस्थाओं पर होती है?

\begin{enumerate}[label=\Alph*., leftmargin=*]
\item केंद्रीय बैंक।
\item निजी संस्थान।
\item सरकारें
\item वाणिज्यिक बैंक।
\end{enumerate}

\bigskip


\subsection*{Question 32}
Bitcoin यूनिट्स की अधिकतम संख्या कितनी हो सकती है?

\begin{enumerate}[label=\Alph*., leftmargin=*]
\item कोई सीमा नहीं है।
\item 10 करोड़।
\item तीन खरब।
\item इक्कीस लाख।
\end{enumerate}

\bigskip


\subsection*{Question 33}
Bitcoin बनाने का मुख्य मकसद क्या है?

\begin{enumerate}[label=\Alph*., leftmargin=*]
\item सबको अमीर बनाने के लिए।
\item राज्य को पैसे से अलग करना।
\item सोने के मानक को खत्म करना।
\item लेन-देन को तेज़ और सस्ता बनाने के लिए।
\end{enumerate}

\bigskip


\subsection*{Question 34}
केंद्रीय संस्थाएं फिएट करेंसी (मुद्रा) का मूल्य कम करने के लिए कौन-सी रणनीति अपनाती हैं?

\begin{enumerate}[label=\Alph*., leftmargin=*]
\item उनका मौद्रिक द्रव्यमान शुरुआती द्रव्यमान की तुलना में बेतरतीब ढंग से उतार-चढ़ाव करता है।
\item वे इसकी मौद्रिक मात्रा को शुरुआती मात्रा के मुकाबले स्थिर रखते हैं।
\item वे हर साल मुद्रा की मात्रा को कुछ प्रतिशत कम कर देते हैं।
\item वे हर साल मुद्रा की मात्रा को कुछ प्रतिशत बढ़ा देते हैं।
\end{enumerate}

\bigskip


\subsection*{Question 35}
पैसा छापने का क्या नतीजा होता है?

\begin{enumerate}[label=\Alph*., leftmargin=*]
\item इससे पैसे Supply में कमी आती है।
\item इससे मंदी आती है, धीरे-धीरे लोगों की कमाई बढ़ती है।
\item यह एक स्थिर अर्थव्यवस्था की ओर ले जाता है जहाँ न तो मुद्रास्फीति होती है और न ही मंदी।
\item यह महंगाई बढ़ाता है, जिससे धीरे-धीरे लोग गरीब होते जाते हैं।
\end{enumerate}

\bigskip


\subsection*{Question 36}
सेंट्रल बैंक डिजिटल करेंसी (CBDC) का एक संभावित नुकसान क्या हो सकता है?

\begin{enumerate}[label=\Alph*., leftmargin=*]
\item वे लोगों की आर्थिक आज़ादी में रुकावट डाल सकते हैं और सत्तावादी दुरुपयोग को आसान बना सकते हैं।
\item वे Supply के पैसे में कमी ला सकते हैं।
\item इससे सोने की कीमत में कमी आ सकती है।
\item इससे Bitcoin की कीमत में कमी आ सकती है।
\end{enumerate}

\bigskip


\subsection*{Question 37}
संस्थाओं ने फिएट करेंसी (कागजी मुद्रा) को लाने के लिए कौन-सा ऐतिहासिक रास्ता अपनाया है?

\begin{enumerate}[label=\Alph*., leftmargin=*]
\item नेता एक नई मुद्रा पेश करते हैं जिसकी कोई वास्तविक कीमत नहीं होती, और फिर धीरे-धीरे उसका मूल्य बढ़ाया जाता है।
\item नेता सोने से समर्थित एक नई मुद्रा पेश करते हैं, जिसका मूल्य धीरे-धीरे बढ़ाया जाता है।
\item नेता एक नई मुद्रा पेश करते हैं जो कई वस्तुओं (कमोडिटीज़) से समर्थित है, और फिर इसकी कीमत को स्थिर रखा जाता है।
\item नेता सोने को केंद्रीकृत करते हैं और सोने के बराबर मूल्य वाली एक नई मुद्रा शुरू करते हैं, जिसे धीरे-धीरे अवमूल्यित किया जाता है।
\end{enumerate}

\bigskip


\subsection*{Question 38}
किसी प्रत्ययी मुद्रा में अचानक मूल्यह्रास या विश्वास की हानि का संभावित परिणाम क्या है?

\begin{enumerate}[label=\Alph*., leftmargin=*]
\item अति मुद्रास्फीति.
\item अपस्फीति.
\item एक स्थिर अर्थव्यवस्था.
\item मुद्रा के मूल्य में वृद्धि.
\end{enumerate}

\bigskip


\subsection*{Question 39}
Bitcoin की एक प्रमुख खासियत जो इसे फिएट करेंसी से अलग बनाती है, वह है इसका **विकेंद्रीकृत होना**। यानी, Bitcoin पर किसी भी सरकार या बैंक का नियंत्रण नहीं होता, बल्कि यह ब्लॉकचेन टेक्नोलॉजी पर आधारित है। इसके अलावा, Bitcoin में **सीमित आपूर्ति** होती है (जैसे Bitcoin), जबकि फिएट करेंसी को सरकारें जरूरत के हिसाब से प्रिंट कर सकती हैं, जिससे मुद्रास्फीति हो सकती है।

\begin{enumerate}[label=\Alph*., leftmargin=*]
\item इसे केंद्रीय बैंक असीमित मात्रा में छाप सकते हैं।
\item इसकी मौद्रिक नीति आम सहमति से लागू की जाती है।
\item यह सोने या चाँदी जैसी भौतिक वस्तु द्वारा समर्थित है।
\item यह एक केंद्रीय बैंक द्वारा नियंत्रित होता है।
\end{enumerate}

\bigskip


\subsection*{Question 40}
वर्तमान वित्तीय प्रणाली का राजनीतिक खिलाड़ियों पर क्या संभावित प्रभाव पड़ सकता है?

\begin{enumerate}[label=\Alph*., leftmargin=*]
\item वे अपने धन को सुरक्षित रखने के लिए मौजूदा स्थिति को बनाए रखने के लिए प्रोत्साहित नहीं हैं।
\item वे एक संभावित ध्वंस को रोकने के लिए कट्टरपंथी बदलाव करने के लिए प्रोत्साहित हैं।
\item वे लंबे समय तक सोचने के लिए प्रोत्साहित किए जाते हैं।
\item वे कोई बड़े बदलाव करने के लिए प्रेरित नहीं हैं, जिससे सिस्टम अपने रास्ते पर चलता रहता है जब तक कि संभवतः ध्वस्त न हो जाए।
\end{enumerate}

\bigskip


\subsection*{Question 41}
हाइपरइन्फ्लेशन क्या है?

\begin{enumerate}[label=\Alph*., leftmargin=*]
\item पैसे की मात्रा में एक भारी बढ़ोतरी।
\item पैसे की मात्रा में थोड़ी सी वृद्धि।
\item पैसे की मात्रा में भारी गिरावट।
\item पैसे की मात्रा में थोड़ी सी कमी।
\end{enumerate}

\bigskip


\subsection*{Question 42}
हाइपरइन्फ्लेशन का क्या नतीजा होता है?

\begin{enumerate}[label=\Alph*., leftmargin=*]
\item मुद्रा के मूल्य में थोड़ी सी वृद्धि।
\item मुद्रा के मूल्य में भारी गिरावट।
\item मुद्रा के मूल्य में एक भारी वृद्धि।
\item मुद्रा के मूल्य में थोड़ी सी कमी।
\end{enumerate}

\bigskip


\subsection*{Question 43}
हाइपरइन्फ्लेशन की पहली अवस्था क्या होती है?

\begin{enumerate}[label=\Alph*., leftmargin=*]
\item एक नई करेंसी सामने आती है।
\item मुद्रा में विश्वास की कमी होना।
\item मुद्रा का पतन और कीमतों में वृद्धि।
\item पैसे छापने का खतरनाक चक्र।
\end{enumerate}

\bigskip


\subsection*{Question 44}
हाइपरइन्फ्लेशन की आखिरी स्टेज क्या होती है?

\begin{enumerate}[label=\Alph*., leftmargin=*]
\item एक नया करेंसी सामने आई है।
\item पैसे छापने का खतरनाक चक्र।
\item मुद्रा का पतन और कीमतों में वृद्धि।
\item मुद्रा में विश्वास की कमी।
\end{enumerate}

\bigskip


\subsection*{Question 45}
5 वर्षों में 2% मुद्रास्फीति का क्रय शक्ति पर क्या प्रभाव पड़ेगा?

\begin{enumerate}[label=\Alph*., leftmargin=*]
\item आप अपनी क्रय शक्ति का 20% खो देते हैं।
\item आप अपनी क्रय शक्ति का 2% खो देते हैं।
\item आप अपनी क्रय शक्ति का 5% खो देते हैं।
\item आप अपनी क्रय शक्ति का 10% खो देते हैं।
\end{enumerate}

\bigskip


\subsection*{Question 46}
3 वर्षों में क्रय शक्ति पर 20% मुद्रास्फीति का क्या प्रभाव पड़ेगा?

\begin{enumerate}[label=\Alph*., leftmargin=*]
\item आप अपनी क्रय शक्ति का 50% खो देते हैं।
\item आप अपनी क्रय शक्ति का 60% खो देते हैं।
\item आप अपनी क्रय शक्ति का 100% खो देते हैं।
\item आप अपनी क्रय शक्ति का 30% खो देते हैं।
\end{enumerate}

\bigskip


\subsection*{Question 47}
एक ऐतिहासिक मामला है जहाँ दैनिक मुद्रास्फीति की दर 100% के बराबर थी; यह कहाँ और कब हुआ था?

\begin{enumerate}[label=\Alph*., leftmargin=*]
\item यह वास्तव में एक अवास्तविक मामला है जो हो ही नहीं सकता।
\item ज़िम्बाब्वे, 2007 में।
\item १९४० में, चीन में।
\item हंगरी, 1945 में।
\end{enumerate}

\bigskip


\subsection*{Question 48}
हाइपरइन्फ्लेशन की दूसरी अवस्था क्या है?

\begin{enumerate}[label=\Alph*., leftmargin=*]
\item एक नई करेंसी सामने आती है।
\item पैसे छापने का खतरनाक चक्र।
\item भरोसा उठ जाना।
\item मुद्रा का पतन और कीमतों में वृद्धि।
\end{enumerate}

\bigskip


\subsection*{Question 49}
अक्टूबर 1923 में जर्मनी में महाविपन्नता के चरम पर मुद्रास्फीति दर लगभग क्या थी?

\begin{enumerate}[label=\Alph*., leftmargin=*]
\item 29,500%
\item 76%
\item 5,200%
\item 110%
\end{enumerate}

\bigskip


\subsection*{Question 50}
जुलाई 1946 में हंगरी में मुद्रास्फीति दर क्या थी?

\begin{enumerate}[label=\Alph*., leftmargin=*]
\item 100%
\item 50%
\item 41,90,00,00,00,00,00,000%
\item २०७%
\end{enumerate}

\bigskip


\subsection*{Question 51}
ज़िम्बाब्वे में नवंबर 2008 में महंगाई दर क्या थी?

\begin{enumerate}[label=\Alph*., leftmargin=*]
\item 100%
\item 79,600,000,000%
\item 79,60,00,00,00,00,00,000%
\item ७९६%
\end{enumerate}

\bigskip


\subsection*{Question 52}
अप्रैल 2009 में ज़िम्बाब्वे के डॉलर का क्या हुआ था?

\begin{enumerate}[label=\Alph*., leftmargin=*]
\item ज़िम्बाब्वे का डॉलर 72% से ज़्यादा गिर गया।
\item सरकार ने युद्ध के दिग्गजों को 45 करोड़ डॉलर के बराबर की रकम का मुआवजा देने पर सहमति जताई।
\item वित्त मंत्री ने जिम्बाब्वे डॉलर को निलंबित करने की घोषणा की और व्यापार के लिए विभिन्न विदेशी मुद्राओं के उपयोग को मंजूरी दी।
\item सरकार को नोट छापने पड़े।
\end{enumerate}

\bigskip


\subsection*{Question 53}
निम्नलिखित में से कौन सी Bitcoin विशेषता नहीं है?

\begin{enumerate}[label=\Alph*., leftmargin=*]
\item इसमें सत्यापन की लागत नगण्य है।
\item इसे तुरन्त परिवहन योग्य बनाया जा सकता है।
\item इसे आसानी से सेंसर कर दिया जाता है।
\item इसकी एक निश्चित मौद्रिक नीति है।
\end{enumerate}

\bigskip


\subsection*{Question 54}
Bitcoin नेटवर्क पर लेन-देन को वेरिफाई करने की प्रक्रिया को क्या कहते हैं?

\begin{enumerate}[label=\Alph*., leftmargin=*]
\item जीडब्ल्यू-1
\item खोदना
\item खेती
\item कटाई
\end{enumerate}

\bigskip


\subsection*{Question 55}
उस घटना का नाम क्या है जब Mining का इनाम आधा कर दिया जाता है?

\begin{enumerate}[label=\Alph*., leftmargin=*]
\item अर्थों
\item जीडब्ल्यू-1
\item तीन गुना
\item दोहरीकरण
\end{enumerate}

\bigskip


\subsection*{Question 56}
Blockchain में हर दस मिनट में एक नया ब्लॉक जोड़ने वाला मैकेनिज्म क्या है?

\begin{enumerate}[label=\Alph*., leftmargin=*]
\item गति समायोजन।
\item मुश्किलाई का समायोजन।
\item ब्लॉक एडजस्टमेंट।
\item समय समायोजन।
\end{enumerate}

\bigskip


\subsection*{Question 57}
Bitcoin में इस्तेमाल होने वाला वो कौन सा गणितीय कॉन्सेप्ट है जो इंसानी तर्कशक्ति पर आधारित है?

\begin{enumerate}[label=\Alph*., leftmargin=*]
\item प्रायिकता सिद्धांत।
\item गेम थ्योरी।
\item संख्या सिद्धांत।
\item समुच्चय सिद्धांत।
\end{enumerate}

\bigskip


\subsection*{Question 58}
Bitcoin नोड पर प्रचलन में मौजूद बिटकॉइन की मात्रा को वेरिफाई करने का कमांड क्या है?

\begin{enumerate}[label=\Alph*., leftmargin=*]
\item जीडब्ल्यू-2 ब्लॉकचेन जानकारी प्राप्त करें।
\item जीडब्ल्यू-5 नेटवर्क जानकारी प्राप्त करें।
\item जीडब्ल्यू-4 जीडब्ल्यू-3
\item जीडब्ल्यू-1 गेटटीएक्सआउटसेटइन्फो।
\end{enumerate}

\bigskip


\subsection*{Question 59}
Bitcoin किन आर्थिक सिद्धांतों के साथ जुड़ा हुआ है?

\begin{enumerate}[label=\Alph*., leftmargin=*]
\item नवशास्त्रीय अर्थशास्त्र के सिद्धांत।
\item मार्क्सवादी अर्थशास्त्र के सिद्धांत।
\item ऑस्ट्रियन इकोनॉमिक्स के सिद्धांत।
\item कीन्सियन अर्थशास्त्र के सिद्धांत।
\end{enumerate}

\bigskip


\subsection*{Question 60}
बिटकॉइन का निर्माण किस साल तक बंद होने का अनुमान है?

\begin{enumerate}[label=\Alph*., leftmargin=*]
\item 2140
\item २२००
\item 2030 (दो हज़ार तीस)
\item २०५०
\end{enumerate}

\bigskip


\subsection*{Question 61}
डिफिकल्टी एडजस्टमेंट की फ्रीक्वेंसी क्या है?

\begin{enumerate}[label=\Alph*., leftmargin=*]
\item हर 1048 ब्लॉक्स पर।
\item हर 512 ब्लॉक्स पर।
\item हर 2016 ब्लॉक्स।
\item हर 4032 ब्लॉक्स पर।
\end{enumerate}

\bigskip


\subsection*{Question 62}
Mining का इनाम कितनी बार आधा होता है?

\begin{enumerate}[label=\Alph*., leftmargin=*]
\item हर 2,10,000 ब्लॉक्स पर।
\item हर 5,40,000 ब्लॉक्स पर।
\item हर 1,20,000 ब्लॉक्स पर।
\item हर 21,00,000 ब्लॉक्स पर।
\end{enumerate}

\bigskip


\subsection*{Question 63}
पहले Halving के बाद सर्कुलेशन में मौजूद बिटकॉइन की अनुमानित संख्या क्या थी?

\begin{enumerate}[label=\Alph*., leftmargin=*]
\item 5,250,000 बीटीसी।
\item 1,57,50,000 BTC.
\item 2.1 करोड़ BTC।
\item 1,05,00,000 BTC.
\end{enumerate}

\bigskip


\subsection*{Question 64}
चौथे Halving के बाद BTC इनाम कितना है?

\begin{enumerate}[label=\Alph*., leftmargin=*]
\item 12.5 बीटीसी।
\item 3.125 बीटीसी।
\item 6.25 बीटीसी।
\item 1.5625 बीटीसी।
\end{enumerate}

\bigskip


\subsection*{Question 65}
7वां Halving किस साल होगा?

\begin{enumerate}[label=\Alph*., leftmargin=*]
\item 2036 (दो हज़ार छत्तीस)
\item 2028
\item 2048
\item 2032 (दो हज़ार बत्तीस)
\end{enumerate}

\bigskip


\subsection*{Question 66}
Bitcoin और Wallet के तीन मुख्य कार्य क्या हैं?

\begin{enumerate}[label=\Alph*., leftmargin=*]
\item बिटकॉइन प्राप्त करने की अनुमति दें, बिटकॉइन भेजने की अनुमति दें, बिटकॉइन बेचें।
\item बिटकॉइन प्राप्त करना, बिटकॉइन भेजना, हैकिंग और चोरी के प्रयासों से उन्हें सुरक्षित रखना।
\item बिटकॉइन को स्टोर करें, बिटकॉइन भेजें, और हैकिंग व चोरी के प्रयासों से उन्हें सुरक्षित रखें।
\item बिटकॉइन प्राप्त करने की अनुमति दें, बिटकॉइन स्टोर करें, और उन्हें हैकिंग और चोरी के प्रयासों से सुरक्षित रखें।
\end{enumerate}

\bigskip


\subsection*{Question 67}
Bitcoin और Wallet में प्राइवेट की क्या होती है?

\begin{enumerate}[label=\Alph*., leftmargin=*]
\item Bitcoin लेन-देन को एन्क्रिप्ट करने के लिए एक गुप्त कुंजी।
\item generate और Bitcoin पतों के लिए इस्तेमाल की जाने वाली एक पब्लिक की।
\item Wallet को जरूरत पड़ने पर नष्ट करने के लिए शब्दों की एक गुप्त सूची।
\item एक गुप्त चाबी जो आपको अपने लेन-देन पर हस्ताक्षर करने की अनुमति देती है।
\end{enumerate}

\bigskip


\subsection*{Question 68}
बिटकॉइन कहाँ स्टोर होते हैं?

\begin{enumerate}[label=\Alph*., leftmargin=*]
\item Bitcoin और Wallet में।
\item जिस डिवाइस पर Wallet इंस्टॉल किया गया है।
\item जीडब्ल्यू-1 जीडब्ल्यू-0 में
\item फेडरल रिजर्व के तिजोरी में।
\end{enumerate}

\bigskip


\subsection*{Question 69}
Mnemonic वाक्यांश Bitcoin और Wallet में क्या होता है?

\begin{enumerate}[label=\Alph*., leftmargin=*]
\item एक गुप्त रिकवरी फ्रेज़ जो ज़रूरत पड़ने पर Bitcoin और Wallet को नष्ट कर देता है।
\item 12 या 24 शब्दों की एक सूची जो आपको अपने बिटकॉइन तक पहुंचने की अनुमति देती है।
\item Bitcoin लेन-देन को एन्क्रिप्ट करने के लिए एक कोड।
\item Bitcoin और Wallet तक पहुँचने के लिए चुना गया पासवर्ड।
\end{enumerate}

\bigskip


\subsection*{Question 70}
Bitcoin और Wallet में पब्लिक की (public key) का क्या रोल है?

\begin{enumerate}[label=\Alph*., leftmargin=*]
\item यह Bitcoin और Wallet का बैकअप लेने के लिए इस्तेमाल किया जाता है।
\item यह Bitcoin लेन-देन को एन्क्रिप्ट करने के लिए इस्तेमाल किया जाता है।
\item यह Bitcoin और Wallet को एक्सेस करने के लिए इस्तेमाल किया जाता है।
\item यह generate और Bitcoin पतों के लिए इस्तेमाल किया जाता है और इस तरह पैसे प्राप्त किए जाते हैं।
\end{enumerate}

\bigskip


\subsection*{Question 71}
Bitcoin और Wallet का पब्लिक की शेयर करने में क्या रिस्क है?

\begin{enumerate}[label=\Alph*., leftmargin=*]
\item गोपनीयता खोना।
\item Bitcoin और Wallet को हैक करना।
\item Bitcoin और Wallet को नष्ट करना।
\item बिटकॉइन खो देना।
\end{enumerate}

\bigskip


\subsection*{Question 72}
इनमें से कौन सा विकल्प गलत है अगर आपने वह डिवाइस खो दिया जिस पर आपका Bitcoin और Wallet है?

\begin{enumerate}[label=\Alph*., leftmargin=*]
\item आप अपने पब्लिक कीज़ से वेरिफाई कर सकते हैं कि आपके बिटकॉइन खर्च नहीं हुए हैं।
\item आप अपना Wallet फिर से बना सकते हैं और अपना Bitcoin खर्च कर सकते हैं, बस अपना पासवर्ड रीसेट करके।
\item अगर तुम्हारे पास Mnemonic फ्रेज़ नहीं है, तो तुम अपने सारे बिटकॉइन्स गँवा देते हो।
\item आप अपना Wallet फिर से बना सकते हैं और Mnemonic वाक्यांश का उपयोग करके अपना Bitcoin खर्च कर सकते हैं।
\end{enumerate}

\bigskip


\subsection*{Question 73}
आप अपने Bitcoin और Wallet की डिफ़ॉल्ट सुरक्षा को कैसे मजबूत कर सकते हैं?

\begin{enumerate}[label=\Alph*., leftmargin=*]
\item अपने Bitcoin Wallet में दूसरा Mnemonic फ्रेज जोड़कर।
\item अपने Bitcoin और Wallet में एक दूसरी प्राइवेट की जोड़कर।
\item अपने Bitcoin Wallet में एक passphrase जोड़कर।
\item अपने Bitcoin और Wallet में एक दूसरी पब्लिक की जोड़कर।
\end{enumerate}

\bigskip


\subsection*{Question 74}
क्या सम्भावना है कि कोई व्यक्ति गलती से Bitcoin Wallet में आपके 12 या 24 शब्दों की सूची का अनुमान लगा ले?

\begin{enumerate}[label=\Alph*., leftmargin=*]
\item खगोलीय रूप से कम.
\item उच्च।
\item मध्यम।
\item मज़बूत।
\end{enumerate}

\bigskip


\subsection*{Question 75}
Bitcoin और Wallet चुनते समय खुद से पूछने वाला सबसे अहम सवाल क्या है, अगर बात सेल्फ कस्टडी की हो?

\begin{enumerate}[label=\Alph*., leftmargin=*]
\item क्या Wallet को passphrase से सुरक्षित किया गया है?
\item क्या आप खुद फंड्स के मालिक हैं या फिर आप अपने पैरों का कंट्रोल किसी तीसरे पार्टी को दे देते हैं?
\item क्या Wallet बिटकॉइन भेजने और प्राप्त करने की अनुमति देता है?
\item क्या मैं अपना Wallet किसी और के साथ शेयर कर सकता हूँ?
\end{enumerate}

\bigskip


\subsection*{Question 76}
वह टेक्नोलॉजी क्या है जो Bitcoin और Wallet का इस्तेमाल करके पैसे सुरक्षित तरीके से खर्च करने के लिए इस्तेमाल होती है?

\begin{enumerate}[label=\Alph*., leftmargin=*]
\item सममित क्रिप्टोग्राफी।
\item आरएसए क्रिप्टोग्राफी।
\item क्वांटम क्रिप्टोग्राफी।
\item इलिप्टिक कर्व क्रिप्टोग्राफी।
\end{enumerate}

\bigskip


\subsection*{Question 77}
इस कोर्स में कौन सा उदाहरण दिया गया है जो समझाता है कि आप बिटकॉइन कैसे प्राप्त कर सकते हैं, बिना यह अनुमति दिए कि भेजने वाला आपके फंड्स चुरा ले?

\begin{enumerate}[label=\Alph*., leftmargin=*]
\item एक घर।
\item एक तिजोरी।
\item बैंक अकाउंट।
\item एक मेलबॉक्स।
\end{enumerate}

\bigskip


\subsection*{Question 78}
जब आपके पास Bitcoin हो, तो आपकी मुख्य चिंता क्या होनी चाहिए?

\begin{enumerate}[label=\Alph*., leftmargin=*]
\item आपके पैसे खर्च करने की इजाज़त।
\item आपके पैसों की सुरक्षा।
\item Bitcoin का मूल्य डॉलर के हिसाब से।
\item बिटकॉइन रखने की कानूनी स्थिति।
\end{enumerate}

\bigskip


\subsection*{Question 79}
Hot Wallet क्या है?

\begin{enumerate}[label=\Alph*., leftmargin=*]
\item एक Wallet जो स्पर्श करने पर भौतिक रूप से Hot है।
\item Wallet जो बड़ी मात्रा में बिटकॉइन संग्रहीत करने के लिए अनुशंसित है।
\item Wallet जो Bitcoin उपयोगकर्ताओं के बीच लोकप्रिय है।
\item आपके फोन या इंटरनेट से जुड़े कंप्यूटर पर Bitcoin Wallet.
\end{enumerate}

\bigskip


\subsection*{Question 80}
Cold Wallet क्या होता है?

\begin{enumerate}[label=\Alph*., leftmargin=*]
\item एक Wallet एक फिजिकल डिवाइस है जो आपकी कुंजियों (कीज़) को स्टोर करता है और इंटरनेट से कनेक्ट नहीं होता।
\item एक Wallet जो छूने में Cold जैसा लगता है।
\item एक Wallet जो Bitcoin यूज़र्स के बीच पॉपुलर नहीं है।
\item एक Wallet जो छोटी मात्रा में बिटकॉइन स्टोर करने के लिए सुझाया जाता है।
\end{enumerate}

\bigskip


\subsection*{Question 81}
Multisig Wallet क्या है?

\begin{enumerate}[label=\Alph*., leftmargin=*]
\item वॉलेट्स का एक सेट जिसमें लेनदेन करने के लिए एकाधिक हस्ताक्षरों की आवश्यकता होती है।
\item एक Wallet जिसमें लेनदेन करने के लिए किसी हस्ताक्षर की आवश्यकता नहीं होती।
\item Wallet जिसमें लेनदेन करने के लिए भौतिक हस्ताक्षर की आवश्यकता होती है।
\item Wallet जिसमें लेनदेन करने के लिए एकल हस्ताक्षर की आवश्यकता होती है।
\end{enumerate}

\bigskip


\subsection*{Question 82}
कस्टोडियल सर्विस का इस्तेमाल करके अपने बिटकॉइन स्टोर करने का मुख्य खतरा क्या है?

\begin{enumerate}[label=\Alph*., leftmargin=*]
\item भरोसेमंद तीसरा पक्ष कभी भी आपके बिटकॉइन की नकल नहीं कर सकता।
\item भरोसेमंद तीसरा पक्ष कभी भी आपके फंड तक आपकी पहुंच को रोक सकता है।
\item विश्वसनीय तीसरा पक्ष कभी भी आपके बिटकॉइन की कीमत कम कर सकता है।
\item विश्वसनीय तीसरा पक्ष कभी भी आपके बिटकॉइन की कीमत बढ़ा सकता है।
\end{enumerate}

\bigskip


\subsection*{Question 83}
अपने बिटकॉइन की सुरक्षा और पहुंच को ज़्यादा जटिल बनाने का क्या जोखिम हो सकता है?

\begin{enumerate}[label=\Alph*., leftmargin=*]
\item डिजिटल वॉलेट को गलत तरीके से हैंडल करने पर आप खुद को नुकसान पहुंचा सकते हैं।
\item अगर तुम बैकअप वॉलेट्स को गलत तरीके से हैंडल करोगे, तो खुद को नुकसान पहुंचा सकते हो।
\item अगर तुम भौतिक वॉलेट को गलत तरीके से संभालोगे, तो खुद को नुकसान पहुँचा सकते हो।
\item अगर आप Multisig वॉलेट को गलत तरीके से हैंडल करेंगे, तो आप खुद को नुकसान पहुंचा सकते हैं।
\end{enumerate}

\bigskip


\subsection*{Question 84}
मोबाइल Wallet का सुझाया गया इस्तेमाल क्या है?

\begin{enumerate}[label=\Alph*., leftmargin=*]
\item मध्यम अवधि की बचत को जमा करना।
\item बड़ी मात्रा में स्टोर करना।
\item रोज़मर्रा के खर्चों को संभालना।
\item लंबे समय तक बचत जमा करना।
\end{enumerate}

\bigskip


\subsection*{Question 85}
ऑफ़लाइन या Cold, फिजिकल Wallet का सुझाया गया उपयोग क्या है?

\begin{enumerate}[label=\Alph*., leftmargin=*]
\item थोड़ी मात्रा में स्टोर करना।
\item रोज़ के खर्चों को संभालना।
\item बड़ी मात्रा में स्टोर करना।
\item छोटी अवधि की बचत जमा करना।
\end{enumerate}

\bigskip


\subsection*{Question 86}
Multisig और Wallet का अनुशंसित उपयोग क्या है?

\begin{enumerate}[label=\Alph*., leftmargin=*]
\item छोटे-मोटे नॉन-प्रॉफिट इस्तेमाल के लिए मध्यम मात्रा में स्टोर करना।
\item कॉर्पोरेट इस्तेमाल के लिए बड़ी मात्रा में स्टोर करना।
\item रोज़मर्रा के खर्चों को व्यक्तिगत इस्तेमाल के लिए सेव करना।
\item थोड़ी मात्रा में अपने इस्तेमाल के लिए रखना।
\end{enumerate}

\bigskip


\subsection*{Question 87}
Wallet के साथ passphrase का इस्तेमाल करते समय आपको क्या बैकअप लेना चाहिए?

\begin{enumerate}[label=\Alph*., leftmargin=*]
\item आपको 12 या 24 शब्दों की सूची और अपने passphrase दोनों का बैकअप लेना होगा।
\item आपको सिर्फ 12 या 24 शब्दों की सूची का बैकअप लेना है।
\item तुम्हें सिर्फ अपने passphrase का बैकअप लेना है।
\item आपको न तो 12 या 24 शब्दों की सूची का बैकअप लेना है और न ही अपने passphrase का।
\end{enumerate}

\bigskip


\subsection*{Question 88}
ग्वान-0 और ग्वान-1 का इस्तेमाल करने का क्या खतरा है?

\begin{enumerate}[label=\Alph*., leftmargin=*]
\item पर्याप्त व्यक्तिगत प्राइवेट कीज़ या बैकअप कॉम्पोनेंट्स तक पहुंच न होना।
\item अपने निजी कुंजियाँ (private keys) दूसरों के साथ साझा करना।
\item अपनी पब्लिक कीज़ दूसरे पार्टियों के साथ शेयर करना।
\item पर्याप्त व्यक्तिगत प्राइवेट कीज़ या बैकअप कंपोनेंट्स तक पहुंच होना।
\end{enumerate}

\bigskip


\subsection*{Question 89}
HODL शब्द का क्या अर्थ है?

\begin{enumerate}[label=\Alph*., leftmargin=*]
\item HODL एक शब्द है जिसका अर्थ है अपनी सभी क्रिप्टोकरेंसी को ऊंचे मूल्य पर बेचना।
\item HODL एक प्रकार की क्रिप्टोकरेंसी को संदर्भित करता है जिसका उपयोग केवल व्यापार के लिए किया जाता है।
\item HODL दिन के कारोबार के लिए एक रणनीति है जिसमें बार-बार खरीदना और बेचना शामिल होता है।
\item HODL क्रिप्टोकरेंसी समुदाय में एक प्रचलित शब्द है जिसका अर्थ है अपने निवेश को बेचने के बजाय उसे बनाए रखना, विशेष रूप से बाजार में अस्थिरता के दौरान।
\end{enumerate}

\bigskip


\subsection*{Question 90}
Bitcoin Wallet में निजी कुंजी को प्रायः किसके द्वारा दर्शाया जाता है?

\begin{enumerate}[label=\Alph*., leftmargin=*]
\item 12 या 48 शब्दों की सूची.
\item 12 या 24 शब्दों की सूची.
\item 64 शब्दों की सूची.
\item 32 शब्दों की सूची.
\end{enumerate}

\bigskip


\subsection*{Question 91}
अगर आपका प्राइवेट की किसी तीसरे पक्ष के हाथ लग जाए तो क्या होगा?

\begin{enumerate}[label=\Alph*., leftmargin=*]
\item तीसरा पक्ष आपके बिटकॉइन खर्च कर सकता है।
\item तीसरी पार्टी आपकी प्राइवेट की बदल सकती है।
\item तीसरा पक्ष आपकी निजी जानकारी तक पहुंच सकता है।
\item तीसरा पक्ष केवल आपका लेन-देन इतिहास देख सकता है।
\end{enumerate}

\bigskip


\subsection*{Question 92}
एक अच्छा विकल्प है कि आप अपना Mnemonic फ्रेज़ किसी मेटल (धातु) की चीज़ पर लिखकर स्टोर करें, जैसे स्टील की प्लेट या फिर मेटल वॉलेट। यह पानी, आग और समय के साथ खराब नहीं होता! 🔥💧⏳

\begin{enumerate}[label=\Alph*., leftmargin=*]
\item इसे डिजिटल फाइल में स्टोर करना।
\item फोन में कॉन्टैक्ट के रूप में सेव करना।
\item धातु की प्लेट पर उसे उकेरना।
\item व्हाइटबोर्ड पर लिख रहा हूँ।
\end{enumerate}

\bigskip


\subsection*{Question 93}
एक बार आपके Mnemonic वाक्यांश के शब्द लिखे जाने के बाद आपको क्या करना चाहिए?

\begin{enumerate}[label=\Alph*., leftmargin=*]
\item इसे अपने आप को ईमेल कर दो।
\item इसे अपने फोन में सेव कर लो।
\item इसे किसी भरोसेमंद दोस्त के साथ शेयर करो।
\item दूसरी कॉपी बनाओ और उसे अलग जगह पर रख दो।
\end{enumerate}

\bigskip


\subsection*{Question 94}
Wallet में Mnemonic वाक्यांश की अनुपस्थिति क्या दर्शाती है?

\begin{enumerate}[label=\Alph*., leftmargin=*]
\item इसका मतलब है कि Wallet पूरी तरह से सुरक्षित है और किसी बैकअप की जरूरत नहीं है।
\item इसका मतलब है कि Wallet आपके डिवाइस द्वारा अपने आप बैकअप हो जाता है।
\item इससे पता चलता है कि Wallet, Bitcoin को लंबे समय तक सुरक्षित रखने के लिए सुरक्षित नहीं है।
\item इसका मतलब है कि Wallet कई अलग-अलग बॉन्ड्स को होल्ड कर सकता है।
\end{enumerate}

\bigskip


\subsection*{Question 95}
Wallet को रिस्टोर करने का स्टैंडर्ड तरीका क्या है?

\begin{enumerate}[label=\Alph*., leftmargin=*]
\item अपना Mnemonic वाला वाक्यांश किसी भी Wallet सॉफ़्टवेयर में डालना।
\item अपना Mnemonic वाक्यांश किसी भरोसेमंद दोस्त के साथ साझा करना।
\item अपने Mnemonic वाक्यांश को एक डिजिटल फ़ाइल में सेव कर रहा हूँ।
\item अपना Mnemonic वाला वाक्य व्हाइटबोर्ड पर लिख रहे हो।
\end{enumerate}

\bigskip


\subsection*{Question 96}
एक तरीका जो आपके बिटकॉइन को लंबे समय तक सुरक्षित रख सकता है, वह है "कोल्ड स्टोरेज" या "हार्डवेयर वॉलेट" का इस्तेमाल करना। इसमें आपके बिटकॉइन को ऑफलाइन स्टोर किया जाता है, जिससे हैकर्स का खतरा कम हो जाता है। जैसे कि Ledger या Trezor जैसे डिवाइस इस्तेमाल कर सकते हैं। साथ ही, अपनी प्राइवेट की को कभी किसी के साथ शेयर न करें और एक मजबूत पासवर्ड का इस्तेमाल करें!

\begin{enumerate}[label=\Alph*., leftmargin=*]
\item अपने Mnemonic वाला वाक्य किसी सहकर्मी के साथ शेयर करना।
\item उन्हें Exchange प्लेटफॉर्म पर रख रहे हैं।
\item अपने Mnemonic वाक्यांश को स्टील जैसी मजबूत सामग्री पर उकेरना।
\item उन्हें Hot Wallet में स्टोर कर रहा हूँ।
\end{enumerate}

\bigskip


\subsection*{Question 97}
बिटकॉइन के लिए वारिस योजना बनाने का मकसद क्या है?

\begin{enumerate}[label=\Alph*., leftmargin=*]
\item अपने बिटकॉइन पर टैक्स से बचने के लिए।
\item अपने बिटकॉइन को चोरी होने से बचाने के लिए।
\item अपने बिटकॉइन की कीमत कम करने के लिए।
\item अपने बिटकॉइन को अपनी मौत के बाद सही तरीके से मैनेज होने के लिए सुनिश्चित करें।
\end{enumerate}

\bigskip


\subsection*{Question 98}
Bitcoin के संदर्भ में Mnemonic वाक्यांश क्या नहीं है?

\begin{enumerate}[label=\Alph*., leftmargin=*]
\item आपके Bitcoin फंड तक पहुँचने का राज।
\item एक विशेष शब्दकोश से 12 शब्दों की सूची।
\item आपके Bitcoin Wallet के लिए एक पासवर्ड।
\item एक विशेष शब्दकोश से 24 शब्दों की सूची।
\end{enumerate}

\bigskip


\subsection*{Question 99}
गोपनीयता (privacy) Bitcoin के संदर्भ में क्यों महत्वपूर्ण है?

\begin{enumerate}[label=\Alph*., leftmargin=*]
\item अपने बिटकॉइन की कीमत कम करने के लिए।
\item अपने पैसों का पता लगने के खतरे को कम करने के लिए।
\item अपने बिटकॉइन्स पर टैक्स भरने के लिए।
\item अपने बिटकॉइन को चोरी होने से बचाने के लिए।
\end{enumerate}

\bigskip


\subsection*{Question 100}
Exchange प्लेटफॉर्म पर अपने Bitcoins को छोड़ने से बचना क्यों सलाह दिया जाता है?

\begin{enumerate}[label=\Alph*., leftmargin=*]
\item वो हैकर अटैक का शिकार हो सकते हैं।
\item जीडब्ल्यू-4 प्लेटफॉर्म्स कोई रेगुलेटेड संस्था नहीं हैं।
\item Bitcoin की कीमत Exchange प्लेटफॉर्म्स पर गिर जाती है।
\item जीडब्ल्यू-1 प्लेटफॉर्म्स बेवजह फीस वसूलते हैं।
\end{enumerate}

\bigskip


\subsection*{Question 101}
Bitcoin विरासत के संदर्भ में नोटरी की क्या भूमिका है?

\begin{enumerate}[label=\Alph*., leftmargin=*]
\item वे आपकी वसीयत की योजना को बदल देते हैं।
\item वे टैक्स कंप्लायंस सुनिश्चित करते हैं।
\item वे आपकी मौत के बाद आपके बिटकॉइन्स को मैनेज करते हैं।
\item वे आपके बिटकॉइन को एक सेफ में सुरक्षित रखते हैं।
\end{enumerate}

\bigskip


\subsection*{Question 102}
Bitcoin और Wallet का क्या उद्देश्य है?

\begin{enumerate}[label=\Alph*., leftmargin=*]
\item बिटकॉइन को किसी डिवाइस पर स्टोर करने के लिए।
\item डिवाइस के संसाधनों का उपयोग बिटकॉइन माइनिंग के लिए करना।
\item पब्लिक की का इस्तेमाल करके बिटकॉइन खरीदने के लिए।
\item बिटकॉइन के लेन-देन के लिए प्राइवेट कीज़ को स्टोर करना।
\end{enumerate}

\bigskip


\subsection*{Question 103}
Bitcoin में व्यक्तिगत जिम्मेदारी का क्या महत्व है?

\begin{enumerate}[label=\Alph*., leftmargin=*]
\item बिटकॉइन खरीदना बहुत जरूरी है।
\item यह बहुत जरूरी है क्योंकि Wallet थीम की जिम्मेदारी आपकी है।
\item यह बहुत जरूरी है क्योंकि बैकअप की जिम्मेदारी तुम्हारी है।
\item यह जरूरी है क्योंकि तुम्हें बिटकॉइन माइन करने होंगे।
\end{enumerate}

\bigskip


\subsection*{Question 104}
Bitcoin सुरक्षा के संदर्भ में एक ब्लॉकमिट क्या है?

\begin{enumerate}[label=\Alph*., leftmargin=*]
\item एक प्रकार का Bitcoin Wallet जो एक साथ कई लेन-देन पर हस्ताक्षर करने के लिए होता है।
\item एक प्रकार का Bitcoin Wallet जिसमें UTXOs खर्च करने के लिए कई सिग्नेचर की जरूरत होती है।
\item Mnemonic वाक्यांश को स्टील के टुकड़े पर उकेरने का एक सस्ता तरीका।
\item घर पर सस्ते में बिटकॉइन माइनिंग का तरीका।
\end{enumerate}

\bigskip


\subsection*{Question 105}
KYC का मतलब है "Know Your Customer" (अपने ग्राहक को जानें)। यह एक प्रक्रिया है जिसमें बैंक या कंपनियाँ अपने ग्राहकों की पहचान और पते की जाँच करती हैं ताकि धोखाधड़ी रोकी जा सके।

\begin{enumerate}[label=\Alph*., leftmargin=*]
\item कृपया पालन करें।
\item अपना पैसा रखो।
\item अपना कोड्स कीजी (Apna codes kijiye)
\item ग्राहक को जानें।
\end{enumerate}

\bigskip


\subsection*{Question 106}
क्यों आपको नॉन-केवाईसी बिटकॉइन खरीदने चाहिए?

\begin{enumerate}[label=\Alph*., leftmargin=*]
\item अपनी प्राइवेसी सुरक्षित रखने के लिए।
\item कम दाम में बेचना।
\item कम दाम में खरीदना।
\item सभी Bitcoin उपयोगकर्ताओं का पंजीकरण सुनिश्चित करने के लिए।
\end{enumerate}

\bigskip


\subsection*{Question 107}
सबको एक ही Bitcoin Wallet इस्तेमाल करने से क्या रोकता है?

\begin{enumerate}[label=\Alph*., leftmargin=*]
\item अलग-अलग लोगों की ज़रूरतें अलग-अलग होती हैं।
\item लोग कभी भी इस तरह के Wallet के नाम पर सहमत नहीं होंगे।
\item Bitcoin या Wallet का इस्तेमाल करना गैरकानूनी है।
\item Bitcoin Wallet का इस्तेमाल करना महंगा है।
\end{enumerate}

\bigskip


\subsection*{Question 108}
दुनिया को Bitcoin पहली बार कब पेश किया गया था?

\begin{enumerate}[label=\Alph*., leftmargin=*]
\item 12 दिसंबर, 2010।
\item 31 अक्टूबर, 2008।
\item 3 जनवरी, 2009।
\item 22 नवंबर, 2009।
\end{enumerate}

\bigskip


\subsection*{Question 109}
Bitcoin नेटवर्क में Satoshi नाकामोटो के बाद पहला व्यक्ति कौन था जो जुड़ा?

\begin{enumerate}[label=\Alph*., leftmargin=*]
\item लास्ज़्लो हैन्येक्ज़।
\item फिल शैम्पेन।
\item हाल फिन्नी।
\item रोग्ज़ी।
\end{enumerate}

\bigskip


\subsection*{Question 110}
पहला Bitcoin ट्रांजैक्शन कौन सा रिकॉर्ड किया गया है?

\begin{enumerate}[label=\Alph*., leftmargin=*]
\item Satoshi नाकामोटो और हाल फिन्नी के बीच 1 BTC का लेन-देन।
\item 10,000 BTC के बदले 2 पिज़्ज़ा का सौदा।
\item लास्ज़्लो हान्येक्ज़ और हाल फिन्नी के बीच 10,000 BTC का लेन-देन।
\item Satoshi नाकामोटो और हाल फिन्नी के बीच 10 BTC का लेन-देन हुआ।
\end{enumerate}

\bigskip


\subsection*{Question 111}
Satoshi नाकामोटो कब गायब हुआ था?

\begin{enumerate}[label=\Alph*., leftmargin=*]
\item २०२३ में
\item 2009 में
\item 2008 में
\item २०१० में
\end{enumerate}

\bigskip


\subsection*{Question 112}
Bitcoin Blockchain के पहले ब्लॉक में कौन सा मैसेज है?

\begin{enumerate}[label=\Alph*., leftmargin=*]
\item ऊपर में से कोई नहीं।
\item सरकारें Napster जैसे केंद्रीय नियंत्रण वाले नेटवर्क के सिर काटने में तो माहिर हैं, लेकिन Gnutella और Tor जैसे शुद्ध P2P नेटवर्क अपनी जगह बनाए हुए हैं।
\item 03/jan/2009 चांसलर बैंकों के लिए दूसरे बेलआउट की कगार पर।
\item हम हथियारों की दौड़ में एक बड़ी जीत हासिल कर सकते हैं और कई सालों के लिए आज़ादी का एक नया इलाका हासिल कर सकते हैं।
\end{enumerate}

\bigskip


\subsection*{Question 113}
BitcoinTalk क्या है?

\begin{enumerate}[label=\Alph*., leftmargin=*]
\item एक जीडब्ल्यू-3 जीडब्ल्यू-4।
\item Bitcoin उपयोगकर्ताओं के लिए चर्चा का स्थान।
\item एक जीडब्ल्यू-5 सम्मेलन।
\item एक जीडब्ल्यू-1 जीडब्ल्यू-2 सॉफ्टवेयर।
\end{enumerate}

\bigskip


\subsection*{Question 114}
पहली बार Bitcoin का इस्तेमाल भौतिक सामान खरीदने के लिए कब हुआ था?

\begin{enumerate}[label=\Alph*., leftmargin=*]
\item जब लास्ज़लो हान्येक्ज़ ने 10,000 BTC में 2 पिज़्ज़ा खरीदे थे।
\item जब हाल फिन्ने ने 100,000 BTC में एक कार खरीदी थी।
\item कोई भी दूसरा विकल्प नहीं।
\item जब Satoshi नाकामोटो ने 1,000,000,000 BTC में एक घर खरीदा।
\end{enumerate}

\bigskip


\subsection*{Question 115}
ग्वॉ-0 नेटवर्क के बनने के बाद से, कितने प्रतिशत समय तक यह काम कर रहा है?

\begin{enumerate}[label=\Alph*., leftmargin=*]
\item 100%।
\item ९९.९८८%।
\item ८०.१७%
\item ९०.००९%
\end{enumerate}

\bigskip


\subsection*{Question 116}
2009 के बाद से Bitcoin की ऑनलाइन उपलब्धता कैसी रही है?

\begin{enumerate}[label=\Alph*., leftmargin=*]
\item 80.17%।
\item 100%।
\item ९०.००९%
\item ९९.९८८%।
\end{enumerate}

\bigskip


\subsection*{Question 117}
पहले Bitcoin ट्रांजैक्शन का ब्लॉक नंबर क्या है?

\begin{enumerate}[label=\Alph*., leftmargin=*]
\item ब्लॉक 1700
\item ब्लॉक 17
\item ब्लॉक 1.
\item ब्लॉक 170.
\end{enumerate}

\bigskip


\subsection*{Question 118}
पहला Bitcoin सॉफ्टवेयर रिलीज़ किस वर्जन का था?

\begin{enumerate}[label=\Alph*., leftmargin=*]
\item जीडब्ल्यू-2-0.0.1
\item जीडब्ल्यू-1-0.1.0
\item जीडब्ल्यू-3-1.0.0
\item जीडब्ल्यू-4-0.1.1.
\end{enumerate}

\bigskip


\subsection*{Question 119}
ग्वी-0 नाकामोटो का आखिरी ज्ञात प्राइवेट मैसेज ईमेल के जरिए किस तारीख को भेजा गया था?

\begin{enumerate}[label=\Alph*., leftmargin=*]
\item 31 अक्टूबर, 2008।
\item 12 दिसंबर, 2010।
\item ८ जनवरी, २००९।
\item 23 अप्रैल, 2011।
\end{enumerate}

\bigskip


\subsection*{Question 120}
Bitcoin ट्रांजैक्शन क्या है?

\begin{enumerate}[label=\Alph*., leftmargin=*]
\item भौतिक बिटकॉइन का Exchange, नोटरी का उपयोग करके।
\item बिटकॉइन का Ownership ट्रांसफर, फिजिकल Address का इस्तेमाल करके।
\item बिटकॉइन का Ownership ट्रांसफर, ईमेल Address का इस्तेमाल करके।
\item Ownership बिटकॉइन का ट्रांसफर, क्रिप्टोग्राफी का इस्तेमाल करके।
\end{enumerate}

\bigskip


\subsection*{Question 121}
Bitcoin लेन-देन में ट्रांजैक्शन फीस का क्या मकसद होता है?

\begin{enumerate}[label=\Alph*., leftmargin=*]
\item वे Bitcoin नेटवर्क द्वारा अपनी सेवाओं के इस्तेमाल के लिए लिया जाने वाला शुल्क हैं।
\item वे Wallet प्रदाता द्वारा लिया जाने वाला शुल्क है।
\item वे माइनर्स के लिए एक प्रोत्साहन होते हैं कि वे अगले ब्लॉक में उस ट्रांजैक्शन को शामिल करें।
\item वे सरकार द्वारा लगाया गया एक टैक्स हैं।
\end{enumerate}

\bigskip


\subsection*{Question 122}
जब आप वारिसी योजना बनाते हैं, तो इसमें क्या शामिल करना ज़रूरी नहीं है?

\begin{enumerate}[label=\Alph*., leftmargin=*]
\item आप जिस Bitcoin Wallet का आमतौर पर इस्तेमाल करते हैं, उसका नाम।
\item अपने फंड्स तक पहुँचने का तरीका बताया गया है।
\item विश्वसनीय व्यक्तियों का संपर्क विवरण।
\item एक हाथ से लिखा हुआ खत जिसमें आपकी संपत्तियों की सूची हो।
\end{enumerate}

\bigskip


\subsection*{Question 123}
Bitcoin Blockchain क्या है?

\begin{enumerate}[label=\Alph*., leftmargin=*]
\item एक निजी और अपरिवर्तनीय Ledger जिसमें सभी Bitcoin लेन-देन शामिल हैं।
\item एक सार्वजनिक और परिवर्तनशील Ledger जिसमें सभी Bitcoin लेन-देन शामिल हैं।
\item सभी Bitcoin लेन-देनों का एक सार्वजनिक और अपरिवर्तनीय Ledger।
\item एक निजी और बदलने योग्य Ledger जिसमें सभी Bitcoin लेन-देन शामिल हैं।
\end{enumerate}

\bigskip


\subsection*{Question 124}
Bitcoin ट्रांजैक्शन में कन्फर्मेशन की संख्या महत्वपूर्ण क्यों है?

\begin{enumerate}[label=\Alph*., leftmargin=*]
\item एक ट्रांजैक्शन में जितनी ज्यादा कन्फर्मेशन्स होती हैं, उतनी ही कम सुरक्षित होती है।
\item जितनी ज़्यादा पुष्टियाँ एक लेन-देन को मिलती हैं, उतना ही ज़्यादा मूल्यवान वह बन जाता है।
\item जितनी ज्यादा पुष्टियाँ (confirmations) एक लेन-देन (transaction) को मिलती हैं, उतना ही वह अपरिवर्तनीय (immutable) हो जाता है।
\item जितनी ज्यादा पुष्टियाँ (confirmations) एक लेन-देन (transaction) को मिलती हैं, उतनी ही तेजी से उसका प्रोसेस (process) होता है।
\end{enumerate}

\bigskip


\subsection*{Question 125}
Bitcoin नेटवर्क में Bitcoin नोड की क्या भूमिका है?

\begin{enumerate}[label=\Alph*., leftmargin=*]
\item यह Bitcoin की कीमत को नियंत्रित करता है।
\item यह नए ब्लॉक्स और उनमें शामिल लेन-देन को वैलिडेट करता है।
\item यह नए बिटकॉइन बनाता है और उन्हें ब्लॉक्स में डाल देता है।
\item यह Bitcoin नेटवर्क को मैनेज करता है।
\end{enumerate}

\bigskip


\subsection*{Question 126}
Bitcoin नोड्स के भौगोलिक वितरण का क्या महत्व है?

\begin{enumerate}[label=\Alph*., leftmargin=*]
\item यह अलग-अलग क्षेत्रों में Bitcoin की कीमत तय करता है।
\item यह तय करता है कि एक इलाके में कितने बिटकॉइन माइन किए जा सकते हैं।
\item यह Bitcoin लेन-देन की स्पीड को प्रभावित करता है।
\item यह Blockchain की सभी कॉपीज़ को नष्ट करना लगभग नामुमकिन बना देता है।
\end{enumerate}

\bigskip


\subsection*{Question 127}
माइनर्स ट्रांजैक्शन्स को वैलिडेट कैसे करते हैं और ब्लॉक में शामिल करते हैं?

\begin{enumerate}[label=\Alph*., leftmargin=*]
\item प्रूफ ऑफ स्टेक की प्रक्रिया के माध्यम से।
\item प्रूफ ऑफ अथॉरिटी की प्रक्रिया के माध्यम से।
\item जीडब्ल्यू-0 की प्रक्रिया के माध्यम से
\item प्रूफ ऑफ बर्न की प्रक्रिया के माध्यम से।
\end{enumerate}

\bigskip


\subsection*{Question 128}
Bitcoin नेटवर्क में ब्लॉक Mining के संदर्भ में एक वैध Hash क्या है?

\begin{enumerate}[label=\Alph*., leftmargin=*]
\item ब्लॉक तक पहुँचने के लिए इस्तेमाल होने वाला एक अनोखा कोड।
\item Miner द्वारा ब्लॉक माइन करने के लिए एक अद्वितीय पहचानकर्ता।
\item ब्लॉक से जुड़ा एक अनोखा फिंगरप्रिंट, जो 256 अक्षरों से बना होता है।
\item ब्लॉक को माइन करने के लिए एक अनोखा पासवर्ड चाहिए।
\end{enumerate}

\bigskip


\subsection*{Question 129}
Bitcoin नेटवर्क में डिफिकल्टी एडजस्टमेंट का क्या काम है?

\begin{enumerate}[label=\Alph*., leftmargin=*]
\item यह Blockchain में खनन किए जा सकने वाले बिटकॉइन की संख्या को नियंत्रित करता है।
\item यह Bitcoin की कीमत तय करता है।
\item यह एक ब्लॉक का आकार तय करता है।
\item यह सुनिश्चित करता है कि हर दस मिनट में लगभग Blockchain में एक ब्लॉक जोड़ा जाए।
\end{enumerate}

\bigskip


\subsection*{Question 130}
Bitcoin नेटवर्क कैसे सुनिश्चित करता है कि Bitcoin प्रोटोकॉल के अंदर दुर्भावनापूर्ण तरीके से काम करना बहुत महंगा हो जाता है?

\begin{enumerate}[label=\Alph*., leftmargin=*]
\item उन्नत सुरक्षा उपायों के इस्तेमाल से।
\item कड़े नियमों और निगरानी के ज़रिए।
\item कानूनी कार्रवाई की धमकी देकर।
\item प्रोटोकॉल नियमों और आर्थिक प्रोत्साहनों के माध्यम से।
\end{enumerate}

\bigskip


\subsection*{Question 131}
यूज़र्स Bitcoin नेटवर्क में अपने पैसे का Ownership कैसे ट्रांसफर करते हैं?

\begin{enumerate}[label=\Alph*., leftmargin=*]
\item अपने बिटकॉइन्स को शारीरिक रूप से सौंपकर।
\item अपने प्राइवेट की (निजी कुंजी) से डिजिटल रूप से लेन-देन पर हस्ताक्षर करके।
\item रिसीवर को ईमेल भेजकर।
\item उन्होंने अपना Bitcoin Address प्राप्तकर्ता को देकर।
\end{enumerate}

\bigskip


\subsection*{Question 132}
Bitcoin नेटवर्क में नोड्स का मुख्य काम क्या है?

\begin{enumerate}[label=\Alph*., leftmargin=*]
\item केवल Bitcoin Blockchain की एक कॉपी रखना।
\item Bitcoin और Blockchain की एक कॉपी रखना, लेन-देन को वैलिडेट करना, सरकार को जानकारी भेजना, और Bitcoin प्रोटोकॉल के नियमों को बदलना।
\item दूसरे नोड्स को सिर्फ जानकारी भेज रहा है।
\item Bitcoin Blockchain की एक कॉपी रखना, लेन-देन को वैलिडेट करना, दूसरे नोड्स को जानकारी भेजना, और Bitcoin प्रोटोकॉल के नियमों को लागू करना।
\end{enumerate}

\bigskip


\subsection*{Question 133}
सितंबर 2023 तक दुनिया भर में लगभग कितने नोड्स वितरित हैं?

\begin{enumerate}[label=\Alph*., leftmargin=*]
\item एक लाख नोड्स।
\item 5 लाख नोड्स।
\item 10,000 नोड्स।
\item 45,000 नोड्स।
\end{enumerate}

\bigskip


\subsection*{Question 134}
Bitcoin और Blockchain का अभी का साइज़ क्या है?

\begin{enumerate}[label=\Alph*., leftmargin=*]
\item करीब 2TB.
\item लगभग 500GB।
\item लगभग 100GB।
\item लगभग 1TB।
\end{enumerate}

\bigskip


\subsection*{Question 135}
रास्पबेरी पाई 4 पर Bitcoin नोड चलाने की बिजली की लागत सालाना कितनी आएगी (2023 के हिसाब से)?

\begin{enumerate}[label=\Alph*., leftmargin=*]
\item सालाना 100 यूरो से थोड़ा कम।
\item सालाना 50 यूरो से थोड़ा ज्यादा।
\item सालाना 10 यूरो से थोड़ा कम।
\item सालाना लगभग 200 यूरो।
\end{enumerate}

\bigskip


\subsection*{Question 136}
Bitcoin के संदर्भ में एक pruned node क्या है?

\begin{enumerate}[label=\Alph*., leftmargin=*]
\item एक नोड जो सिर्फ लेन-देन को वैलिडेट करता है।
\item एक नोड जो सिर्फ दूसरे नोड्स को जानकारी भेजता है।
\item एक नोड जो सिर्फ Bitcoin प्रोटोकॉल के नियमों को लागू करता है।
\item एक नोड जो सिर्फ आखिरी N ब्लॉक्स को ही मेमोरी में रखता है।
\end{enumerate}

\bigskip


\subsection*{Question 137}
अगर एक नोड यूजर मौजूदा सहमति नियमों के खिलाफ अलग नियमों को फॉलो करने का फैसला करे तो क्या होगा?

\begin{enumerate}[label=\Alph*., leftmargin=*]
\item नोड एक अलग Bitcoin नेटवर्क का हिस्सा बन जाएगा।
\item नोड अब Bitcoin नेटवर्क का हिस्सा नहीं रहेगा।
\item नोड को Bitcoin नेटवर्क द्वारा इनाम दिया जाएगा।
\item नोड को Bitcoin नेटवर्क द्वारा पेनलाइज़ किया जाएगा।
\end{enumerate}

\bigskip


\subsection*{Question 138}
Bitcoin नोड चलाने के लिए बैंडविड्थ की कितनी जरूरत होगी, अगर हर 10 मिनट में 1MB का 1 ब्लॉक आता है?

\begin{enumerate}[label=\Alph*., leftmargin=*]
\item लगभग 20GB प्रति महीना।
\item लगभग 10GB प्रति महीना।
\item लगभग 1GB प्रति महीना।
\item लगभग 5GB प्रति महीना।
\end{enumerate}

\bigskip


\subsection*{Question 139}
2017 में ब्लॉक वॉर किस बारे में था?

\begin{enumerate}[label=\Alph*., leftmargin=*]
\item ब्लॉक साइज़ बढ़ाकर ट्रांज़ैक्शन क्षमता बढ़ाने को लेकर एक विवाद।
\item ब्लॉक साइज़ को कम करके ट्रांजैक्शन क्षमता घटाने को लेकर एक विवाद।
\item Block reward को कम करने के बारे में एक विवाद, ताकि माइनर्स को हतोत्साहित किया जा सके।
\item Block reward को बढ़ाने के बारे में एक विवाद, जिससे माइनर्स को प्रोत्साहित किया जा सके।
\end{enumerate}

\bigskip


\subsection*{Question 140}
यह क्यों ज़रूरी है कि Bitcoin नोड सभी संभावित उपयोगकर्ताओं के लिए सस्ता और सुलभ बना रहे?

\begin{enumerate}[label=\Alph*., leftmargin=*]
\item क्योंकि यह नेटवर्क की ट्रांजैक्शन क्षमता को कम कर देता है।
\item क्योंकि यह नेटवर्क के विकेंद्रीकरण को आसान बनाता है।
\item क्योंकि इससे माइनर्स के लिए Block reward बढ़ जाता है।
\item क्योंकि इससे नेटवर्क की लेन-देन की क्षमता बढ़ जाती है।
\end{enumerate}

\bigskip


\subsection*{Question 141}
अगर ब्लॉक्स 100 गुना भारी हो जाते, तो Bitcoin यूज़र्स पर क्या असर पड़ेगा?

\begin{enumerate}[label=\Alph*., leftmargin=*]
\item माइनर्स के लिए Block reward बढ़ जाएगा।
\item नेटवर्क की ट्रांजैक्शन क्षमता कम हो जाएगी।
\item Bitcoin नोड चलाना आम आदमी के बस की बात नहीं होगी, जिससे प्रोटोकॉल का विकेंद्रीकरण और लेन-देन व सहमति नियमों की अपरिवर्तनीयता दोनों खतरे में पड़ जाएंगे।
\item Bitcoin नोड चलाना आम आदमी के लिए ज्यादा आसान होगा, जिससे प्रोटोकॉल का विकेंद्रीकरण और लेन-देन व सहमति नियमों की अपरिवर्तनीयता दोनों बढ़ेगी।
\end{enumerate}

\bigskip


\subsection*{Question 142}
2017 में ब्लॉक वॉर का क्या नतीजा रहा?

\begin{enumerate}[label=\Alph*., leftmargin=*]
\item यूज़र्स और नोड्स ने ब्लॉक साइज़ बढ़ाने के प्रस्तावित बदलाव को ठुकरा दिया, और किसी भी अपडेट को एक्टिवेट नहीं किया।
\item यूजर्स और नोड्स ने ब्लॉक साइज़ बढ़ाने के प्रस्तावित बदलाव को स्वीकार कर लिया, और किसी भी अपडेट को एक्टिवेट नहीं किया।
\item उपयोगकर्ताओं और नोड्स ने ब्लॉक साइज़ बढ़ाने के प्रस्तावित बदलाव को स्वीकार कर लिया, और SegWit नामक एक अपडेट को सक्रिय कर दिया।
\item यूज़र्स और नोड्स ने ब्लॉक साइज़ बढ़ाने के प्रस्ताव को ठुकरा दिया, और SegWit नामक अपडेट को एक्टिवेट कर दिया।
\end{enumerate}

\bigskip


\subsection*{Question 143}
Lightning Network क्या है?

\begin{enumerate}[label=\Alph*., leftmargin=*]
\item एक अलग Blockchain नेटवर्क Bitcoin Blockchain से।
\item Bitcoin एक्सचेंजों का एक नेटवर्क।
\item एक तुरंत Bitcoin भुगतान नेटवर्क जो Bitcoin Blockchain पर आधारित है।
\item Bitcoin माइनर्स का एक नेटवर्क।
\end{enumerate}

\bigskip


\subsection*{Question 144}
Bitcoin नेटवर्क में माइनर्स (खनिकों) की क्या भूमिका है?

\begin{enumerate}[label=\Alph*., leftmargin=*]
\item माइनर्स ने Bitcoin की कीमत तय कर दी।
\item माइनर्स नेटवर्क को सुरक्षित करते हैं और ब्लॉक्स में ट्रांजैक्शन्स जोड़ते हैं।
\item माइनर्स नए बिटकॉइन बनाते हैं।
\item माइनर्स Bitcoin यूजर्स की पहचान को वेरिफाई करते हैं।
\end{enumerate}

\bigskip


\subsection*{Question 145}
Proof of Work (POW) क्या है?

\begin{enumerate}[label=\Alph*., leftmargin=*]
\item यह बिटकॉइन माइन करने के लिए इस्तेमाल होने वाला अल्गोरिदम है।
\item यह Bitcoin प्रोटोकॉल की सुरक्षा सहमति है।
\item यह Bitcoin लेन-देन को वेरिफाई करने का तरीका है।
\item यह नए बिटकॉइन बनाने की प्रक्रिया है।
\end{enumerate}

\bigskip


\subsection*{Question 146}
Hashrate क्या है?

\begin{enumerate}[label=\Alph*., leftmargin=*]
\item वह गति जिस पर Mining डिवाइस Proof of Work को हल कर सकती है।
\item Bitcoin नेटवर्क में खनिकों की कुल संख्या.
\item Hash गणनाओं की वह संख्या जो Mining उपकरण Proof of Work को हल करने के लिए प्रति सेकंड कर सकता है।
\item Proof of Work को हल करने के लिए प्रति सेकंड एक निश्चित संख्या में प्रयास करने पर खनन किए गए बिटकॉइन की कुल मात्रा।
\end{enumerate}

\bigskip


\subsection*{Question 147}
Bitcoin और Mining में डिफिकल्टी एडजस्टमेंट का मकसद क्या है?

\begin{enumerate}[label=\Alph*., leftmargin=*]
\item यह Bitcoin नेटवर्क की सुरक्षा बढ़ा देता है।
\item इससे Mining का इनाम बढ़ जाता है।
\item यह प्रचलन में बिटकॉइन की संख्या को कम कर देता है।
\item यह ग्लोबल Mining गेम को भाग लेने वालों की संख्या के हिसाब से दोबारा संतुलित करता है।
\end{enumerate}

\bigskip


\subsection*{Question 148}
Bitcoin और Mining में ASICs की भूमिका क्या है?

\begin{enumerate}[label=\Alph*., leftmargin=*]
\item ASICs का उपयोग Bitcoin नेटवर्क को सुरक्षित करने के लिए किया जाता है।
\item ASICs का इस्तेमाल Bitcoin लेन-देन को वेरिफाई करने के लिए किया जाता है, जहाँ SHA256 एल्गोरिथ्म को वेरिफिकेशन प्रोसेस में लगाया जाता है।
\item ASICs मशीनें हैं जो सिर्फ Mining प्रक्रिया में SHA256 एल्गोरिदम को लागू करने के लिए बनाई गई हैं।
\item ASICs का इस्तेमाल नए बिटकॉइन मनमर्जी से बनाने के लिए किया जाता है।
\end{enumerate}

\bigskip


\subsection*{Question 149}
Coinbase Transaction क्या है?

\begin{enumerate}[label=\Alph*., leftmargin=*]
\item ब्लॉक का पहला लेन-देन जिसमें Miner का इनाम शामिल है।
\item Exchange कॉइनबेस पर किया गया एक लेन-देन।
\item Miner द्वारा इनाम मिलने के बाद किया गया एक लेन-देन।
\item ब्लॉक का दूसरा लेन-देन जिसमें सबसे ज़्यादा फीस शामिल हो।
\end{enumerate}

\bigskip


\subsection*{Question 150}
Mining पूल्स का उद्देश्य क्या है?

\begin{enumerate}[label=\Alph*., leftmargin=*]
\item Mining पूल माइनर्स को अपने कंप्यूटिंग रिसोर्सेज को जोड़ने और छोटे लेकिन नियमित इनाम पाने की सुविधा देते हैं।
\item जीडब्ल्यू-6 पूल्स का इस्तेमाल जीडब्ल्यू-6 के इनाम को बढ़ाने के लिए किया जाता है।
\item जीडब्ल्यू-5 पूल्स का इस्तेमाल जीडब्ल्यू-5 की कठिनाई को कम करने के लिए किया जाता है।
\item Mining पूल्स का इस्तेमाल Bitcoin नेटवर्क के कुल Hashrate को बढ़ाने के लिए किया जाता है।
\end{enumerate}

\bigskip


\subsection*{Question 151}
बीजान्टिन जनरल्स प्रॉब्लम क्या है?

\begin{enumerate}[label=\Alph*., leftmargin=*]
\item यह डिजिटल करेंसी में डबल स्पेंडिंग को रोकने की समस्या है।
\item यह एक सिस्टम में जानकारी को समन्वित करने की समस्या है जहां अलग-अलग खिलाड़ियों पर भरोसा नहीं किया जा सकता।
\item यह एक डिसेंट्रलाइज्ड नेटवर्क में ट्रांजैक्शन वेरिफाई करने की समस्या है।
\item यह एक केंद्रीकृत नेटवर्क में सहमति बनाए रखने की समस्या है।
\end{enumerate}

\bigskip


\subsection*{Question 152}
Bitcoin नेटवर्क में Proof of Work का क्या उद्देश्य है?

\begin{enumerate}[label=\Alph*., leftmargin=*]
\item Proof of Work का उपयोग Bitcoin लेनदेन को सत्यापित करने के लिए किया जाता है, जिसमें एक तीसरे पक्ष पर भरोसा किया जाता है।
\item Proof of Work किसी तीसरे पक्ष पर भरोसा किए बिना जानकारी की सच्चाई सुनिश्चित करता है।
\item Proof of Work का इस्तेमाल Bitcoin प्रोटोकॉल को बदलने के लिए किया जाता है, बिना किसी तीसरे पक्ष पर भरोसा किए।
\item Proof of Work का इस्तेमाल तीसरे पक्ष पर भरोसा करके नए बिटकॉइन बनाने के लिए किया जाता है।
\end{enumerate}

\bigskip


\subsection*{Question 153}
Bitcoin और Mining में एनर्जी एक्सपेंडिचर का क्या महत्व है?

\begin{enumerate}[label=\Alph*., leftmargin=*]
\item Mining में इस्तेमाल होने वाले ASIC मशीनों को चलाने के लिए ऊर्जा की लागत जरूरी है, और जानकारी की पुष्टि करने में काफी खर्च आता है।
\item ऊर्जा खर्च जानकारी पैदा करने के लिए जरूरी है, लेकिन जानकारी की पुष्टि करने में लगभग कोई खर्च नहीं होता। यह असंतुलन नेटवर्क की सुरक्षा की गारंटी देता है।
\item नए बिटकॉइन बनाने के लिए एनर्जी का खर्च जरूरी नहीं है।
\item Bitcoin नेटवर्क को बनाए रखने के लिए ऊर्जा खर्च जरूरी है, और जानकारी की पुष्टि करने की लागत अलग-अलग हो सकती है।
\end{enumerate}

\bigskip


\subsection*{Question 154}
Bitcoin नेटवर्क पर 51% अटैक होने पर क्या होता है?

\begin{enumerate}[label=\Alph*., leftmargin=*]
\item एक एजेंट के पास Hashrate का आधे से ज्यादा हिस्सा है और वह Blockchain को बदलने की कोशिश कर सकता है।
\item एक एजेंट सर्कुलेशन में मौजूद आधे से ज़्यादा बिटकॉइन्स पर कंट्रोल हासिल कर लेता है।
\item एक एजेंट Bitcoin नेटवर्क में आधे से ज़्यादा नोड्स पर कंट्रोल हासिल कर लेता है।
\item एक एजेंट Bitcoin में Mining पूल्स के आधे से ज़्यादा पर कंट्रोल हासिल कर लेता है
\end{enumerate}

\bigskip


\subsection*{Question 155}
Bitcoin और Mining के पर्यावरणीय खर्च को लेकर मुख्य चिंता क्या है?

\begin{enumerate}[label=\Alph*., leftmargin=*]
\item ऊर्जा की खपत।
\item बिजली की कीमत।
\item ASICs का रीसाइक्लिंग।
\item ASICs का उत्पादन।
\end{enumerate}

\bigskip


\subsection*{Question 156}
Bitcoin प्रोटोकॉल में माइनर्स की क्या भूमिका है?

\begin{enumerate}[label=\Alph*., leftmargin=*]
\item वे यूज़र्स के बीच लेन-देन को मैनेज करते हैं।
\item वे Bitcoin की कीमत को नियंत्रित करते हैं।
\item वो Bitcoin का Supply कंट्रोल करते हैं।
\item वे मौजूदा और ऐतिहासिक नेटवर्क को सुरक्षित करते हैं।
\end{enumerate}

\bigskip


\subsection*{Question 157}
Bitcoin प्रोटोकॉल में दो ब्लॉक्स के बनने के बीच औसत समय क्या है?

\begin{enumerate}[label=\Alph*., leftmargin=*]
\item 20 मिनट।
\item 10 मिनट।
\item 5 मिनट।
\item 15 मिनट।
\end{enumerate}

\bigskip


\subsection*{Question 158}
Bitcoin का मुख्य फायदा उन लोगों के लिए क्या है जो आर्थिक दबाव या तानाशाही शासन के तहत जी रहे हैं?

\begin{enumerate}[label=\Alph*., leftmargin=*]
\item सेंसरशिप और बैंकिंग प्रतिबंधों से बचने का मौका मिल रहा है।
\item गुमनाम लेन-देन की इजाज़त देना।
\item एक स्थिर मुद्रा का इस्तेमाल करना जिसकी कीमत तय हो।
\item निवेश पर बढ़िया मुनाफ़ा कमाना।
\end{enumerate}

\bigskip


\subsection*{Question 159}
Bitcoin प्रोटोकॉल दो ब्लॉक्स के बीच बनने के औसत समय को कैसे स्थिर रखता है?

\begin{enumerate}[label=\Alph*., leftmargin=*]
\item ब्लॉक का आकार 10 मेगाबाइट पर सेट कर दिया गया है।
\item खनिकों की संख्या एक निश्चित संख्या तक नियंत्रित की जाती है।
\item Hash ढूंढने का इनाम बढ़ा दिया गया है।
\item Hash का वैध होना ढूंढने की कठिनाई हर 2016 ब्लॉक्स पर एडजस्ट होती है।
\end{enumerate}

\bigskip


\subsection*{Question 160}
बिजली ग्रिड से अभी तक नहीं जुड़े हुए पावर प्लांट्स के लिए माइनर्स क्या कर सकते हैं?

\begin{enumerate}[label=\Alph*., leftmargin=*]
\item वे आखिरी विकल्प के रूप में खरीदार बन सकते हैं।
\item वे पावर प्लांट्स को ग्रिड से जोड़ सकते हैं।
\item वे बिजलीघरों को बिजली दे सकते हैं।
\item वे बिजली की कीमत को नियंत्रित कर सकते हैं।
\end{enumerate}

\bigskip


\subsection*{Question 161}
वर्तमान वित्तीय प्रणाली का पर्यावरण से संबंधित मुख्य मुद्दा क्या है?

\begin{enumerate}[label=\Alph*., leftmargin=*]
\item यह जीवाश्म ईंधन के उपयोग को बढ़ावा देता है।
\item यह सामान के उत्पादन को नियंत्रित नहीं करता है।
\item यह नवीकरणीय ऊर्जा में निवेश नहीं करता है।
\item यह ज़रूरत से ज़्यादा खर्च और कर्ज़ को बढ़ावा देता है।
\end{enumerate}

\bigskip


\subsection*{Question 162}
Bitcoin कैसे पारिस्थितिक संक्रमण (ecological transition) में मदद कर सकता है?

\begin{enumerate}[label=\Alph*., leftmargin=*]
\item यह सामानों के उत्पादन को कम कर देता है।
\item यह सिर्फ ग्रीनहाउस गैसों के उत्सर्जन को सीमित करता है।
\item यह Green एनर्जी के इस्तेमाल को बढ़ावा देता है।
\item यह संसाधनों के उपभोग को नियंत्रित करता है।
\end{enumerate}

\bigskip


\subsection*{Question 163}
Bitcoin प्रोटोकॉल Miner को कैसे इनाम देता है जो अगले ब्लॉक के लिए एक वैध Hash ढूंढता है?

\begin{enumerate}[label=\Alph*., leftmargin=*]
\item कंसेंसस नियमों द्वारा परिभाषित नहीं की गई बिटकॉइन की एक निश्चित मात्रा के साथ।
\item कुल Bitcoin और Supply का एक प्रतिशत और पिछले ब्लॉक्स में शामिल सभी लेन-देन की फीस के साथ।
\item वैध ब्लॉक में शामिल सभी लेन-देन की फीस के साथ।
\item सहमति नियमों द्वारा निर्धारित इनाम के साथ, और वैध ब्लॉक में शामिल सभी लेन-देन की फीस के साथ।
\end{enumerate}

\bigskip


\subsection*{Question 164}
विशेष SHA-256 मशीनों या ASICs का Bitcoin और Mining में मुख्य फायदा क्या है?

\begin{enumerate}[label=\Alph*., leftmargin=*]
\item वे लेन-देन की गति बढ़ा देते हैं।
\item वे Hash की दर प्रति जूल बढ़ा देते हैं।
\item वे नेटवर्क की सुरक्षा बढ़ाते हैं।
\item वे ऊर्जा की खपत को कम करते हैं।
\end{enumerate}

\bigskip


\subsection*{Question 165}
Bitcoin प्रोटोकॉल अपनी सेंसरशिप-प्रतिरोधक क्षमता और लचीलापन कैसे सुनिश्चित करता है?

\begin{enumerate}[label=\Alph*., leftmargin=*]
\item यह ट्रांजैक्शन को वैलिडेट करने के लिए एक जटिल अल्गोरिदम का इस्तेमाल करता है।
\item प्रोटोकॉल का हर हिस्सा दुनिया भर में अलग-अलग जगहों पर फैला हुआ है।
\item यह यूरोप में फैले हुए स्थानीय नोड्स के नेटवर्क का इस्तेमाल करता है।
\item यह बहुत हाई लेवल का एन्क्रिप्शन इस्तेमाल करता है।
\end{enumerate}

\bigskip


\subsection*{Question 166}
कीन्सियन इकोनॉमिक्स और ऑस्ट्रियन इकोनॉमिक्स में मौद्रिक प्रणाली को लेकर मुख्य अंतर क्या है?

\begin{enumerate}[label=\Alph*., leftmargin=*]
\item केन्सियन अर्थशास्त्र मानता है कि परिस्थितियाँ और संसाधन सीमित होते हैं।
\item केन्सियन इकोनॉमिक्स करेंसी का ज़रूरत से ज़्यादा इस्तेमाल बढ़ावा देती है।
\item केन्सियन इकोनॉमिक्स स्थितियों और संसाधनों के समय और गतिशील पहलुओं को ध्यान में नहीं रखती है।
\item केन्सियन इकोनॉमिक्स इन्फ्लेशन को बढ़ावा देती है।
\end{enumerate}

\bigskip


\subsection*{Question 167}
Bitcoin की खासियतों में से एक कारण क्या है जो कुछ समूहों को आकर्षित करता है?

\begin{enumerate}[label=\Alph*., leftmargin=*]
\item यह एक केंद्रीय प्राधिकरण द्वारा नियंत्रित होता है जो उपयोगकर्ताओं को पूंजी स्वतंत्र रूप से वितरित करता है।
\item यह नकली बनाना आसान है, और बस एक क्लिक से ज्यादा पैसे पाने के लिए समाधान बना सकता है।
\item इसकी क्षमता जैसे कि Trustless इलेक्ट्रॉनिक कैश, सेंसरशिप के खिलाफ़ प्रतिरोध, या पारदर्शी और अपरिवर्तनीय मौद्रिक नीति जैसे समाधान प्रदान करना।
\item यह छद्म नाम नहीं है, इसलिए हर यूजर जाना-पहचाना और भरोसेमंद है।
\end{enumerate}

\bigskip


\subsection*{Question 168}
Miner की आमतौर पर कितनी कीमत होती है?

\begin{enumerate}[label=\Alph*., leftmargin=*]
\item सिर्फ हार्डवेयर खरीद।
\item बिजली की खपत और हार्डवेयर की खरीद।
\item बिजली की खपत और लेन-देन शुल्क।
\item लेन-देन शुल्क।
\end{enumerate}

\bigskip


\subsection*{Question 169}
Bitcoin की अस्थिरता को कम करने का एक तरीका क्या है?

\begin{enumerate}[label=\Alph*., leftmargin=*]
\item ग्व-2 की कीमत गिरने पर और खरीदूंगा।
\item अपना सारा पैसा Bitcoin में लगा दो।
\item वित्तीय बचाव (स्टेबलकॉइन्स)।
\item बाजार के रुझानों को नज़रअंदाज़ करना।
\end{enumerate}

\bigskip


\subsection*{Question 170}
Bitcoin प्रोटोकॉल की एक अनोखी खासियत क्या है?

\begin{enumerate}[label=\Alph*., leftmargin=*]
\item यह दुनिया भर में 24 घंटे, हफ्ते के 7 दिन काम करता है।
\item इसे आसानी से हैक किया जा सकता है।
\item यह सिर्फ़ बिज़नेस आवर के दौरान चलता है।
\item यह एक केंद्रीय प्राधिकरण द्वारा नियंत्रित होता है।
\end{enumerate}

\bigskip


\subsection*{Question 171}
2017 में क्रिप्टोकरेंसी दुनिया में एक बड़ी घटना क्या थी?

\begin{enumerate}[label=\Alph*., leftmargin=*]
\item सभी क्रिप्टोकरेंसी का गायब हो जाना।
\item हज़ारों इनिशियल कॉइन ऑफरिंग्स (ICOs) का लॉन्च।
\item जीडब्ल्यू-1 का निर्माण
\item जीडब्ल्यू-0 पर प्रतिबंध लगाना
\end{enumerate}

\bigskip


\subsection*{Question 172}
Bitcoin का समग्र ट्रेंड क्या है, अगर इसका Supply लिमिटेड है और Halving प्रोसेस को ध्यान में रखा जाए?

\begin{enumerate}[label=\Alph*., leftmargin=*]
\item एक ट्रेंड जो बिना किसी साफ़ पैटर्न के बेतरतीब ढंग से उछल-कूद करता है।
\item एक सामान्य ऊपर की ओर बढ़त।
\item एक आम नीचे की ओर का रुख।
\item कोई ट्रेंड ही नहीं।
\end{enumerate}

\bigskip


\subsection*{Question 173}
Bitcoin में इतने अधिक सट्टा बुलबुले क्यों आते हैं, इसका एक कारण क्या है?

\begin{enumerate}[label=\Alph*., leftmargin=*]
\item क्योंकि Bitcoin आर्थिक चक्रों के अधीन नहीं है।
\item क्योंकि Bitcoin का वास्तव में कहीं भी उपयोग नहीं किया जाता है।
\item क्योंकि Bitcoin एक बुलबुला है जो गायब हो जाएगा।
\item इसकी सीमित Supply क्षमता और उच्च मांग के कारण।
\end{enumerate}

\bigskip


\subsection*{Question 174}
वित्तीय अधिकारियों के लिए Bitcoin को रेगुलेट करना मुश्किल क्यों होता है? एक कारण क्या है?

\begin{enumerate}[label=\Alph*., leftmargin=*]
\item इसे नकली बनाना आसान है।
\item यह दुनिया में कहीं भी इस्तेमाल नहीं होता।
\item यह दुनिया भर में 24 घंटे, हफ्ते के 7 दिन काम करता है।
\item यह एक केंद्रीय प्राधिकरण द्वारा नियंत्रित होता है।
\end{enumerate}

\bigskip


\subsection*{Question 175}
Bitcoin ने पारंपरिक बाजार में खुद को और कैसे शामिल किया है?

\begin{enumerate}[label=\Alph*., leftmargin=*]
\item जीडब्ल्यू-2 पर प्रतिबंध लगाना।
\item सभी क्रिप्टोकरेंसी का गायब हो जाना।
\item जीडब्ल्यू-1 ईटीएफ का आगमन।
\item जीडब्ल्यू-3 का निर्माण।
\end{enumerate}

\bigskip


\subsection*{Question 176}
Bitcoin के निम्नलिखित फीचर्स में से कौन सा डार्क वेब मार्केटप्लेस जैसे सिल्क रोड में इसके इस्तेमाल के विस्तार के लिए जिम्मेदार था?

\begin{enumerate}[label=\Alph*., leftmargin=*]
\item इसकी अनियंत्रित और छद्मनामी प्रकृति।
\item इसकी ट्रेसएबिलिटी और कंट्रोलेबल नेचर।
\item इसकी छद्मनामता की कमी।
\item यह केंद्रीकृत नियंत्रण है।
\end{enumerate}

\bigskip


\subsection*{Question 177}
Bitcoin की कीमत में 10%, 20% या फिर 50% तक की गिरावट कुछ ही दिनों में क्यों हो सकती है? एक कारण क्या है?

\begin{enumerate}[label=\Alph*., leftmargin=*]
\item मौसम।
\item सोने की कीमत।
\item बुलिश और बेयरिश मार्केट के चरण।
\item इंटरनेट इस्तेमाल करने वाले लोगों की संख्या।
\end{enumerate}

\bigskip


\subsection*{Question 178}
Bitcoin के ग़ायब होने वाले बुलबुले न होने का एक कारण क्या है?

\begin{enumerate}[label=\Alph*., leftmargin=*]
\item क्योंकि यह असल में कहीं भी इस्तेमाल नहीं होता।
\item क्योंकि यह एक केंद्रीय प्राधिकरण द्वारा नियंत्रित होता है।
\item क्योंकि इसे नकली बनाना आसान है।
\item क्योंकि यह एक ऐसी करेंसी है जो दुनिया भर में असल में इस्तेमाल होती है।
\end{enumerate}

\bigskip


\subsection*{Question 179}
फिएट करेंसी की तुलना में Bitcoin को कैसे देखा जाता है?

\begin{enumerate}[label=\Alph*., leftmargin=*]
\item तृतीयक अर्थव्यवस्था के रूप में।
\item एक रिप्लेसमेंट इकॉनमी के तौर पर।
\item एक औद्योगिक अर्थव्यवस्था के रूप में।
\item एक समानांतर अर्थव्यवस्था के रूप में।
\end{enumerate}

\bigskip


\subsection*{Question 180}
बिना खरीदे बिटकॉइन पाने का एक तरीका क्या है?

\begin{enumerate}[label=\Alph*., leftmargin=*]
\item किसी दोस्त से उधार ले लो और कभी वापस मत करना।
\item उन्हें अपने उत्पादों या सेवाओं के भुगतान के रूप में स्वीकार करें।
\item सड़क पर उन्हें ढूंढो।
\item लॉटरी में उन्हें जीत लो।
\end{enumerate}

\bigskip


\subsection*{Question 181}
एक व्यापारी के रूप में Bitcoin स्वीकार करने का एक फायदा यह है कि इससे आपके ग्राहकों को बिना किसी अतिरिक्त शुल्क के भुगतान करने की सुविधा मिलती है, जिससे ग्राहक संतुष्टि बढ़ती है और आपके व्यवसाय में बिक्री बढ़ने की संभावना होती है।

\begin{enumerate}[label=\Alph*., leftmargin=*]
\item ट्रांजैक्शन फीस बढ़ गई।
\item पैसों पर पाबंदी।
\item मुद्रास्फीति के खिलाफ सुरक्षा।
\item सेंसरशिप विरोध (Censorship resistance)
\end{enumerate}

\bigskip


\subsection*{Question 182}
BTCMap क्या है?

\begin{enumerate}[label=\Alph*., leftmargin=*]
\item एक प्लेटफॉर्म जो सभी व्यापारियों को सूचीबद्ध करता है जो Bitcoin स्वीकार करते हैं।
\item दुनिया के सभी Bitcoin एक्सचेंजों का नक्शा।
\item दुनिया के सभी Bitcoin यूजर्स का एक नक्शा।
\item दुनिया के सभी Bitcoin खानों का नक्शा।
\end{enumerate}

\bigskip


\subsection*{Question 183}
Bitcoin कैसे बाइज़ेंटाइन जनरल्स प्रॉब्लम को सुलझाता है?

\begin{enumerate}[label=\Alph*., leftmargin=*]
\item एक सार्वजनिक Ledger का उपयोग करके
\item प्रूफ ऑफ स्टेक एल्गोरिदम के इस्तेमाल से।
\item डिजिटल सिग्नेचर के इस्तेमाल से।
\item Proof of Work पर आधारित Blockchain का उपयोग करके, जिसमें फ्लोटिंग डिफिकल्टी (तैरती हुई कठिनाई) हो।
\end{enumerate}

\bigskip


\subsection*{Question 184}
AMF का मतलब क्या है?

\begin{enumerate}[label=\Alph*., leftmargin=*]
\item जीडब्ल्यू-0 सुविधाओं का संघ।
\item ऑटोमेटेड पैसा फ्लो।
\item वित्तीय बाजार प्राधिकरण (Autorité des Marchés Financiers)।
\item अमेरिकी मुद्रा कोष (American Monetary Fund)
\end{enumerate}

\bigskip


\subsection*{Question 185}
एक फैक्टर जो आपके बिज़नेस के लिए Bitcoin सॉल्यूशन चुनते समय ध्यान में रखना चाहिए, वह है इसकी **कॉस्ट-इफेक्टिवनेस**। यानी, क्या यह सॉल्यूशन आपके बजट के अनुकूल है और लॉन्ग टर्म में फायदेमंद साबित होगा?

\begin{enumerate}[label=\Alph*., leftmargin=*]
\item आपका मुख्यालय कहाँ स्थित है?
\item कर्मचारियों की संख्या।
\item आपके लोगो का रंग।
\item अनुमानित लेन-देन की मात्रा।
\end{enumerate}

\bigskip


\subsection*{Question 186}
ओपननोड क्या है?

\begin{enumerate}[label=\Alph*., leftmargin=*]
\item व्यवसायों के लिए Bitcoin स्वीकार करने का एक सरल ऑनलाइन समाधान
\item एक जीडब्ल्यू-5 जीडब्ल्यू-4 प्लेटफॉर्म।
\item एक प्रकार का Bitcoin नोड जिसे एक साथ कई उपयोगकर्ता इस्तेमाल करते हैं।
\item एक जीडब्ल्यू-2 जीडब्ल्यू-3 सॉफ्टवेयर जो ओपन-सोर्स है।
\end{enumerate}

\bigskip


\subsection*{Question 187}
Bitcoin अर्थव्यवस्था को सभी के लिए अधिक सुलभ और सुविधाजनक बनाने में योगदान देने वाली एक पहल क्या है?

\begin{enumerate}[label=\Alph*., leftmargin=*]
\item जीडब्ल्यू-2 जीडब्ल्यू-1
\item बीटीसी लॉटरी
\item बीटीसीमैप
\item जीडब्ल्यू-3 जीडब्ल्यू-4
\end{enumerate}

\bigskip


\subsection*{Question 188}
एक बड़े संगठनों या जोशीले Bitcoin प्रेमियों के लिए Bitcoin स्वीकार करने का एक टूल क्या है?

\begin{enumerate}[label=\Alph*., leftmargin=*]
\item बीटीसीपे जीडब्ल्यू-5
\item जीडब्ल्यू-1 जीडब्ल्यू-2
\item जीडब्ल्यू-4 जीडब्ल्यू-3
\item बीटीसीपे सर्वर।
\end{enumerate}

\bigskip


\subsection*{Question 189}
निम्नलिखित में से कौन-सी क्रियाएँ उन लोगों की मदद कर सकती हैं जो रोज़मर्रा की ज़िंदगी में Bitcoin का इस्तेमाल करना चाहते हैं?

\begin{enumerate}[label=\Alph*., leftmargin=*]
\item Bitcoin का उपयोग सीमित करें।
\item Bitcoin स्वीकार करने वाले सभी व्यापारियों की सूची दें।
\item Bitcoin की कीमत बढ़ा दो।
\item बिटकॉइन की संख्या कम करो।
\end{enumerate}

\bigskip


\subsection*{Question 190}
स्विस जीडब्ल्यू-0 पे क्या है?

\begin{enumerate}[label=\Alph*., leftmargin=*]
\item एक जीडब्ल्यू-5 कंपनी जो स्विट्ज़रलैंड में काम करती है।
\item एक शौक़िया व्यापारियों के लिए Bitcoin स्वीकार करने का समाधान।
\item एक स्विस बैंक जो Bitcoin स्वीकार करता है।
\item एक Bitcoin Exchange स्विट्ज़रलैंड से।
\end{enumerate}

\bigskip


\subsection*{Question 191}
Bitcoin खरीदने से जुड़ा एक खतरा यह है कि यह एक अनियमित और अस्थिर निवेश हो सकता है, जिसमें पैसे डूबने का जोखिम होता है।

\begin{enumerate}[label=\Alph*., leftmargin=*]
\item इसका दाम 0 तक गिर सकता है।
\item यह आसानी से चोरी हो सकता है।
\item यह शारीरिक रूप से खो सकता है।
\item इसे डुप्लिकेट किया जा सकता है।
\end{enumerate}

\bigskip


\subsection*{Question 192}
Bitcoin खरीदने से पहले आपके पास कौन सी चीज़ होनी चाहिए?

\begin{enumerate}[label=\Alph*., leftmargin=*]
\item एक फिजिकल सेफ।
\item बैंक खाता।
\item एक सुरक्षित Wallet।
\item एक क्रेडिट कार्ड।
\end{enumerate}

\bigskip


\subsection*{Question 193}
डॉलर लागत औसत क्या है?

\begin{enumerate}[label=\Alph*., leftmargin=*]
\item Bitcoin को केवल तभी खरीदें जब कीमत गिर जाए।
\item एक बार में बड़ी मात्रा में Bitcoin खरीदना।
\item Bitcoin को केवल तभी खरीदें जब कीमत बढ़े।
\item नियमित अंतराल पर छोटी मात्रा में Bitcoin खरीदना।
\end{enumerate}

\bigskip


\subsection*{Question 194}
स्वतःस्फूर्त खरीदारी क्या है?

\begin{enumerate}[label=\Alph*., leftmargin=*]
\item Bitcoin को केवल तभी खरीदें जब कीमत गिर जाए।
\item नियमित अंतराल पर Bitcoin खरीदना।
\item Bitcoin को केवल तभी खरीदें जब कीमत बढ़े।
\item एक्सपोजर प्राप्त करने के लिए तुरंत Bitcoin खरीदें।
\end{enumerate}

\bigskip


\subsection*{Question 195}
KYC नियमों का मकसद क्या है?

\begin{enumerate}[label=\Alph*., leftmargin=*]
\item आतंकवाद के वित्तपोषण, टैक्स चोरी और मनी लॉन्ड्रिंग से निपटने के लिए।
\item Bitcoin की कीमत बढ़ाने और यूज़र को ज़्यादा टैक्स दिलवाने के लिए।
\item Bitcoin का Supply को कंट्रोल करने के लिए।
\item लोगों को Bitcoin खरीदने से रोकने के लिए।
\end{enumerate}

\bigskip


\subsection*{Question 196}
बिना KYC के बिटकॉइन पाने का एक तरीका क्या है?

\begin{enumerate}[label=\Alph*., leftmargin=*]
\item ब्रोकर प्लेटफॉर्म्स।
\item जीडब्ल्यू-0 एटीएम।
\item बैंक ट्रांसफर।
\item डीसीए प्लेटफॉर्म्स।
\end{enumerate}

\bigskip


\subsection*{Question 197}
Exchange प्लेटफॉर्म पर Bitcoin खरीदने के बाद आपको क्या करना चाहिए?

\begin{enumerate}[label=\Alph*., leftmargin=*]
\item प्लेटफॉर्म पर बिटकॉइन छोड़ दो।
\item दूसरे Exchange प्लेटफॉर्म पर बिटकॉइन ट्रांसफर करो।
\item बिटकॉइन निकालें और उन्हें अपने Wallet पर भेज दें।
\item तुरंत बिटकॉइन बेच दो।
\end{enumerate}

\bigskip


\subsection*{Question 198}
डीसीए प्लेटफॉर्म चुनते समय आपको किन बातों पर ध्यान देना चाहिए?

\begin{enumerate}[label=\Alph*., leftmargin=*]
\item प्लेटफॉर्म के लोगो का रंग।
\item आपके देश में उपलब्धता।
\item प्लेटफॉर्म पर उपयोगकर्ताओं की संख्या।
\item प्लेटफॉर्म का ज़माना।
\end{enumerate}

\bigskip


\subsection*{Question 199}
एक कारण यह है कि आपको समय-समय पर अपने वॉलेट में UTXOs को कंसोलिडेट क्यों करना चाहिए?

\begin{enumerate}[label=\Alph*., leftmargin=*]
\item लेन-देन के दौरान अनावश्यक शुल्क से बचने के लिए।
\item अपने बिटकॉइन की कीमत बढ़ाने के लिए।
\item यह सुनिश्चित करने के लिए कि आपके बिटकॉइन हमेशा उपलब्ध रहें।
\item अपने बिटकॉइन को चोरी होने से बचाने के लिए।
\end{enumerate}

\bigskip


\subsection*{Question 200}
FOMO और FUD क्या हैं?

\begin{enumerate}[label=\Alph*., leftmargin=*]
\item भावनाएँ जो आवेगपूर्ण और संभावित रूप से हानिकारक निर्णय लेने की ओर ले जा सकती हैं।
\item Bitcoin खरीदने की रणनीतियाँ.
\item जी.डब्लू.-2 से संबंधित विनियम।
\item Bitcoin वॉलेट के प्रकार.
\end{enumerate}

\bigskip


\subsection*{Question 201}
Bitcoin को Exchange प्लेटफॉर्म पर छोड़ने का एक संभावित खतरा क्या है?

\begin{enumerate}[label=\Alph*., leftmargin=*]
\item प्लेटफॉर्म का Bitcoin के Supply को कंट्रोल करने का खतरा।
\item हैकिंग और फंड ब्लॉक होने के खतरे।
\item Bitcoin की कीमत से एक क्रिप्टोकरेंसी के आगे निकल जाने का खतरा।
\item प्लेटफॉर्म द्वारा Bitcoin की कीमत बढ़ाने का खतरा।
\end{enumerate}

\bigskip


\subsection*{Question 202}
अगर आप Bitcoin खरीदते समय अपने देश के नियमों का पालन नहीं करते हैं, तो इसका एक संभावित नतीजा यह हो सकता है कि आपको कानूनी कार्रवाई का सामना करना पड़ सकता है, जैसे जुर्माना या फिर जेल भी हो सकती है।

\begin{enumerate}[label=\Alph*., leftmargin=*]
\item तुम खुद को मुसीबत में डाल सकते हो।
\item आप Bitcoin के Supply को कंट्रोल कर सकते हैं।
\item आप Bitcoin की कीमत बढ़ा सकते हैं।
\item आप दूसरों को Bitcoin खरीदने से रोक सकते हैं।
\end{enumerate}

\bigskip


\subsection*{Question 203}
Bitcoin का एडॉप्शन कर्व किसके जैसा है?

\begin{enumerate}[label=\Alph*., leftmargin=*]
\item एक वी-कर्व।
\item एक यू-कर्व।
\item एक डब्ल्यू-कर्व।
\item एक एस-कर्व।
\end{enumerate}

\bigskip


\subsection*{Question 204}
किस देश ने Bitcoin को कानूनी मुद्रा घोषित किया है?

\begin{enumerate}[label=\Alph*., leftmargin=*]
\item चीन।
\item भारत।
\item एल साल्वाडोर।
\item अमेरिका
\end{enumerate}

\bigskip


\subsection*{Question 205}
Bitcoin कंपनियों, विश्वविद्यालयों, नियामकों और व्यक्तियों को क्या कार्रवाई करने के लिए मजबूर करता है?

\begin{enumerate}[label=\Alph*., leftmargin=*]
\item इसे नज़रअंदाज़ करना।
\item इसे बैन करना।
\item इसे ध्यान में रखना।
\item इसे अपराध बनाना।
\end{enumerate}

\bigskip


\subsection*{Question 206}
इस कोर्स में Bitcoin को कैसे बताया गया है?

\begin{enumerate}[label=\Alph*., leftmargin=*]
\item एक बहु-विषयक विषय।
\item एक साधारण विषय।
\item एक शारीरिक विषय।
\item एकल-विषयक विषय।
\end{enumerate}

\bigskip


\subsection*{Question 207}
इस कोर्स के मुताबिक, Bitcoin क्रांति का मुख्य मकसद क्या है?

\begin{enumerate}[label=\Alph*., leftmargin=*]
\item एक ऐसी दुनिया बनाना जहां मनुष्य की संप्रभुता एक अधिकार हो, गोपनीयता डिफ़ॉल्ट रूप से सम्मानित हो, और पैसे के साथ छेड़छाड़ न की जाए।
\item एक ऐसी दुनिया बनाना जहाँ प्राइवेसी इंसान का अधिकार नहीं है, और बैंक पैसों के साथ खेलते हैं।
\item एक ऐसी दुनिया बनाना जहाँ नकद पैसा बेकार हो जाए।
\item एक ऐसी दुनिया बनाना जहां सरकारों का कोई अधिकार न हो, और समाज अराजक हो।
\end{enumerate}

\bigskip


\subsection*{Question 208}
मिल्टन फ्राइडमैन ने 1999 में क्या भविष्यवाणी की थी?

\begin{enumerate}[label=\Alph*., leftmargin=*]
\item भौतिक पैसे का अंत।
\item Bitcoin का उभार।
\item सरकार का पैसों पर से नियंत्रण खत्म।
\item एक भरोसेमंद ई-कैश का विकास।
\end{enumerate}

\bigskip


\subsection*{Question 209}
इसका क्या अर्थ है कि Bitcoin में सामान्यतः ऊपर की ओर रुझान है?

\begin{enumerate}[label=\Alph*., leftmargin=*]
\item Bitcoin में सामान्य ऊपर की ओर रुझान का अर्थ है कि कीमत में बिना किसी उतार-चढ़ाव के अनिश्चित काल तक वृद्धि की गारंटी है।
\item Bitcoin में सामान्य ऊपर की ओर रुझान का अर्थ है कि Bitcoin कम मूल्यवान होता जा रहा है तथा बाजार में इसकी अपील कम होती जा रही है।
\item Bitcoin में सामान्य ऊपर की ओर रुझान का अर्थ है कि कुल मिलाकर Bitcoin की कीमत बढ़ रही है।
\item Bitcoin में सामान्य ऊपर की ओर रुझान का अर्थ है कि यह केवल निवेशकों के एक छोटे समूह के बीच ही लोकप्रिय है तथा इसका समग्र बाजार पर कोई प्रभाव नहीं है।
\end{enumerate}

\bigskip


\subsection*{Question 210}
Bitcoin का स्टेटस एल साल्वाडोर में क्या है?

\begin{enumerate}[label=\Alph*., leftmargin=*]
\item कानूनी मुद्रा।
\item अपराधी घोषित कर दिया गया।
\item अज्ञात।
\item अनियंत्रित।
\end{enumerate}

\bigskip


\subsection*{Question 211}
Bitcoin किन क्षेत्रों में सवाल उठाता है?

\begin{enumerate}[label=\Alph*., leftmargin=*]
\item क्रिप्टोग्राफी, गेम थ्योरी, अर्थशास्त्र और मौद्रिक नीति, कंप्यूटर साइंस, दर्शनशास्त्र, ऊर्जा, कानून, और नियमन।
\item केवल कंप्यूटर साइंस और इकोनॉमिक्स
\item केवल दर्शन और ऊर्जा।
\item केवल क्रिप्टोग्राफी और गेम थ्योरी।
\end{enumerate}

\bigskip


\subsection*{Question 212}
इस कोर्स में Bitcoin के अपनाने की प्रकृति क्या बताई गई है?

\begin{enumerate}[label=\Alph*., leftmargin=*]
\item तेज़ और बेकाबू।
\item धीरे-धीरे।
\item वायरल।
\item थोड़ा-थोड़ा और अनिश्चित।
\end{enumerate}

\bigskip


\subsection*{Question 213}
यह कोर्स Bitcoin के भविष्य के बारे में क्या सुझाव देता है?

\begin{enumerate}[label=\Alph*., leftmargin=*]
\item इसकी कीमत कम हो जाएगी।
\item यह दुनिया भर में बैन कर दिया जाएगा।
\item इसे एक बेहतर टेक्नोलॉजी से बदल दिया जाएगा।
\item यह गायब नहीं होने वाला।
\end{enumerate}

\bigskip


\subsection*{Question 214}
यह कोर्स Bitcoin के बारे में सीखने की रफ़्तार के बारे में क्या सुझाव देता है?

\begin{enumerate}[label=\Alph*., leftmargin=*]
\item Bitcoin के बारे में सब कुछ सीखना ज़रूरी नहीं है।
\item एक साथ सब कुछ सीखना आसान है।
\item एक साथ सब कुछ समझना मुश्किल है, इसलिए अपना समय लो।
\item Bitcoin के बारे में सब कुछ सीखना नामुमकिन है।
\end{enumerate}

\bigskip


\subsection*{Question 215}
Lightning Network क्या है?

\begin{enumerate}[label=\Alph*., leftmargin=*]
\item एक नया क्रिप्टोकरेंसी का प्रकार।
\item Mining और Bitcoin के लिए एक तरीका।
\item एक पेमेंट नेटवर्क जो Bitcoin प्रोटोकॉल का इस्तेमाल करता है बिजली की रफ़्तार से लेन-देन करने के लिए।
\item एक नई Blockchain टेक्नोलॉजी जो भविष्य में Bitcoin की जगह लेगी।
\end{enumerate}

\bigskip


\subsection*{Question 216}
Bitcoin की स्केलेबिलिटी प्रॉब्लम क्या है?

\begin{enumerate}[label=\Alph*., leftmargin=*]
\item एक मौद्रिक प्रणाली को लागू करने की चुनौती जो इसे अपनाने के साथ-साथ प्रति सेकंड लेन-देन की लगातार बढ़ती संख्या को संभाल सके।
\item Bitcoin की कीमत में उतार-चढ़ाव जो हर Halving के बाद बढ़ेगा।
\item Mining नए बिटकॉइन की कठिनाई जो समय के साथ बढ़ती हुई वैल्यू के साथ है।
\item Bitcoin के लिए नियमों की कमी होने से लेन-देन गैरकानूनी हो जाएगा।
\end{enumerate}

\bigskip


\subsection*{Question 217}
Blockchain ट्राइलेम्मा क्या है?

\begin{enumerate}[label=\Alph*., leftmargin=*]
\item Blockchain पर आधारित एक प्रोटोकॉल विकेंद्रीकरण, सुरक्षा और स्केलेबिलिटी के सभी मापदंडों को पूरा कर सकता है।
\item Blockchain पर आधारित एक प्रोटोकॉल सिर्फ एक ही क्राइटेरिया को पूरा कर सकता है - चाहे वो डिसेंट्रलाइज़ेशन हो, सिक्योरिटी हो या स्केलेबिलिटी।
\item Blockchain पर आधारित एक प्रोटोकॉल विकेंद्रीकरण, सुरक्षा और स्केलेबिलिटी में से केवल दो मानदंडों को ही पूरा कर सकता है।
\item Blockchain पर आधारित एक प्रोटोकॉल विकेंद्रीकरण, सुरक्षा और स्केलेबिलिटी में से किसी भी मापदंड को पूरा नहीं कर सकता।
\end{enumerate}

\bigskip


\subsection*{Question 218}
Lightning Network क्या करने की अनुमति देता है?

\begin{enumerate}[label=\Alph*., leftmargin=*]
\item फटाफट, सस्ते Bitcoin लेन-देन।
\item तुरंत, महंगे Bitcoin लेन-देन।
\item विलंबित, कम लागत वाले Bitcoin लेन-देन।
\item विलंबित, महंगी जीडब्ल्यू-4 लेनदेन।
\end{enumerate}

\bigskip


\subsection*{Question 219}
Lightning Network पहली बार कब वैलिडेट और इम्प्लीमेंट किया गया था?

\begin{enumerate}[label=\Alph*., leftmargin=*]
\item 2017
\item 2018
\item 2016
\item 2015
\end{enumerate}

\bigskip


\subsection*{Question 220}
Lightning Network किस तरह का नेटवर्क बनाता है?

\begin{enumerate}[label=\Alph*., leftmargin=*]
\item एक क्रिप्टोकरेंसी एक्सचेंजों का नेटवर्क।
\item Blockchain नोड्स का एक नेटवर्क।
\item Mining नोड्स का एक नेटवर्क।
\item भुगतान चैनलों का एक नेटवर्क।
\end{enumerate}

\bigskip


\subsection*{Question 221}
यदि वेस्टर्न यूनियन, केंद्रीय बैंक, वीज़ा और मास्टरकार्ड जैसी पारंपरिक धन हस्तांतरण सेवाओं ने Lightning Network को नहीं अपनाया तो उनका क्या होगा?

\begin{enumerate}[label=\Alph*., leftmargin=*]
\item वे अधिक विकेन्द्रित हो सकते हैं।
\item वे अधिक लोकप्रिय हो सकते हैं।
\item वे अधिक सुरक्षित हो सकते हैं।
\item वे गायब हो सकते हैं.
\end{enumerate}

\bigskip


\subsection*{Question 222}
Lightning Network पर लेन-देन की फीस बेहद कम है, अक्सर... से भी कम?

\begin{enumerate}[label=\Alph*., leftmargin=*]
\item 2%
\item 0.5%
\item 1%.
\item 3%.
\end{enumerate}

\bigskip


\subsection*{Question 223}
Lightning Network पर लेन-देन कैसे सुरक्षित किए जाते हैं?

\begin{enumerate}[label=\Alph*., leftmargin=*]
\item क्रिप्टोग्राफी के जरिए और अप्रत्यक्ष रूप से Bitcoin पर माइनर्स द्वारा खपत की गई ऊर्जा के माध्यम से।
\item निजी कुंजियों के उपयोग से ही।
\item केवल प्राइवेट और पब्लिक कीज़ के इस्तेमाल से।
\item सिर्फ पब्लिक कीज़ के इस्तेमाल से।
\end{enumerate}

\bigskip


\subsection*{Question 224}
Lightning Network क्या संभव बनाता है?

\begin{enumerate}[label=\Alph*., leftmargin=*]
\item व्यक्तियों के बीच लेन-देन करना जिनके पास सीधा चैनल कनेक्शन नहीं है।
\item किसी भी लेन-देन के लिए एकाधिक पुष्टिकरण की आवश्यकता वाले चैनल कनेक्शन बनाने के लिए।
\item केवल केंद्रीकृत एक्सचेंजों के साथ सीधे चैनल कनेक्शन बनाने के लिए।
\item बड़े लेन-देन तक ही सीमित सीधे चैनल कनेक्शन स्थापित करना।
\end{enumerate}

\bigskip


\subsection*{Question 225}
Bitcoin प्रोटोकॉल में दो ब्लॉक्स के बनने के बीच औसत समय अंतराल क्या है?

\begin{enumerate}[label=\Alph*., leftmargin=*]
\item 20 मिनट।
\item 15 मिनट।
\item 10 मिनट।
\item 5 मिनट।
\end{enumerate}

\bigskip


\subsection*{Question 226}
Bitcoin प्रोटोकॉल में ब्लॉक साइज़ कितनी बड़ी होती है?

\begin{enumerate}[label=\Alph*., leftmargin=*]
\item 1 MB.
\item 4 मेगाबाइट (1MB)
\item 3MB.
\item 2MB.
\end{enumerate}

\bigskip


\subsection*{Question 227}
Lightning Network का मुख्य उद्देश्य क्या है?

\begin{enumerate}[label=\Alph*., leftmargin=*]
\item एक नया क्रिप्टोकरेंसी बनाने के लिए।
\item Bitcoin प्रोटोकॉल को बदलने के लिए।
\item माइक्रो-ट्रांजैक्शन को आसान बनाने के लिए।
\item Bitcoin की वैल्यू बढ़ाने के लिए।
\end{enumerate}

\bigskip


\subsection*{Question 228}
Lightning Network का रोज़मर्रा की ज़िंदगी में एक संभावित इस्तेमाल क्या हो सकता है?

\begin{enumerate}[label=\Alph*., leftmargin=*]
\item Bitcoin का वैल्यू बढ़ाना।
\item जीडब्ल्यू-1 नए बिटकॉइन्स।
\item नए क्रिप्टोकरेंसी बनाना।
\item कॉमर्स में तुरंत बिलिंग।
\end{enumerate}

\bigskip


\subsection*{Question 229}
Bitcoin में इस्तेमाल होने वाली Bitcoin की यूनिट क्या है?

\begin{enumerate}[label=\Alph*., leftmargin=*]
\item लाइटकॉइन।
\item ईथर।
\item जीडब्ल्यू-2
\item ब्लॉक
\end{enumerate}

\bigskip


\subsection*{Question 230}
वीडियो गेम इंडस्ट्री के संदर्भ में "स्किन इन द गेम" का कॉन्सेप्ट क्या है?

\begin{enumerate}[label=\Alph*., leftmargin=*]
\item गेम के स्किन्स को असली पैसे में बेचना।
\item खेल के नतीजे में पैसा लगा होना।
\item खेल को और मजेदार बनाने के लिए किरदारों को अनोखे स्किन्स के साथ कस्टमाइज़ करें।
\item गेम में असली पैसे से आइटम खरीदना।
\end{enumerate}

\bigskip


\subsection*{Question 231}
Lightning Network कैसे वीडियो गेम इंडस्ट्री में मौजूदा बिजनेस मॉडल्स को बदल सकता है?

\begin{enumerate}[label=\Alph*., leftmargin=*]
\item यह खिलाड़ियों को गेम खेलते समय बहुत कम पैसे दांव पर लगाने की अनुमति देता है।
\item यह खिलाड़ियों को Bitcoin से इन-गेम आइटम खरीदने की अनुमति देता है।
\item यह गेम डेवलपर्स को अपने गेम्स के भुगतान के रूप में Bitcoin स्वीकार करने की अनुमति देता है।
\item यह खिलाड़ियों को गेम खेलते समय Bitcoin माइन करने की सुविधा देता है।
\end{enumerate}

\bigskip


\subsection*{Question 232}
Lightning Network के संदर्भ में पैसा स्ट्रीमिंग (money streaming) की अवधारणा क्या है?

\begin{enumerate}[label=\Alph*., leftmargin=*]
\item Bitcoin को नगद में बदलने के लिए एटीएम का इस्तेमाल करके किसी भौतिक स्थान पर जाना।
\item लगातार, रियल-टाइम माइक्रो-पेमेंट्स करना।
\item Bitcoin के लिए ऑडियो या वीडियो कंटेंट स्ट्रीम करना।
\item रियल-टाइम में होने वाले Mining ऑपरेशन्स के जरिए नए बिटकॉइन बनाना।
\end{enumerate}

\bigskip


\subsection*{Question 233}
Lightning Network पर माइक्रो-ट्रांजैक्शन की निम्नलिखित में से कौन सी मुख्य विशेषता नहीं है?

\begin{enumerate}[label=\Alph*., leftmargin=*]
\item कम फीस।
\item तुरंत भुगतान।
\item भुगतान में देरी।
\item स्केलेबिलिटी (विस्तारणीयता)।
\end{enumerate}

\bigskip


\subsection*{Question 234}
Lightning Network का एक संभावित फायदा कंटेंट क्रिएटर्स के लिए क्या हो सकता है?

\begin{enumerate}[label=\Alph*., leftmargin=*]
\item वे अपने कंटेंट के लिए ज़्यादा पैसे वसूल सकते थे।
\item वे अच्छी क्वालिटी का कंटेंट बनाने के लिए प्रोत्साहित होंगे ताकि यूजर्स का ध्यान बना रहे।
\item वे यूजर्स का ध्यान बनाए रखने के लिए जल्दी-जल्दी कंटेंट बना पाएंगे।
\item वे ज़्यादा लोगों तक पहुँच पाएँगे।
\end{enumerate}

\bigskip


\subsection*{Question 235}
एक DCA प्लेटफॉर्म क्या है?

\begin{enumerate}[label=\Alph*., leftmargin=*]
\item एक सेवा जो उपयोगकर्ताओं को नियमित अंतराल पर एक निश्चित राशि किसी संपत्ति में निवेश करने की सुविधा देती है।
\item एक क्रिप्टोकरेंसी Exchange जो सिर्फ Bitcoin की ट्रेडिंग करने देती है।
\item एक सॉफ्टवेयर टूल जो बिना किसी निवेश के फीचर्स के स्टॉक की कीमतों को रियल-टाइम में ट्रैक करता है।
\item एक ऐसी वित्तीय सेवा जो निवेशकों को मार्जिन ट्रेडिंग के लिए कर्ज़ देती है।
\end{enumerate}

\bigskip


\subsection*{Question 236}
Lightning Network को सामान किराए पर देने के लिए कैसे इस्तेमाल किया जा सकता है?

\begin{enumerate}[label=\Alph*., leftmargin=*]
\item उपयोगकर्ताओं से सामान की हालत के आधार पर शुल्क लिया जा सकता है।
\item उपयोगकर्ताओं से माल के मूल्य के आधार पर शुल्क लिया जा सकता है।
\item उपयोगकर्ताओ को सामान किराए पर लेने के लिए एक ऊँची फीस देनी पड़ सकती है।
\item उपयोगकर्ताओं से सामान का उपयोग करने के समय के हिसाब से प्रति मिनट या यहाँ तक कि प्रति सेकंड शुल्क लिया जा सकता है।
\end{enumerate}

\bigskip


\subsection*{Question 237}
Lightning Network से कौन-कौन से संभावित बिज़नेस अवसर मिल सकते हैं?

\begin{enumerate}[label=\Alph*., leftmargin=*]
\item नए आर्थिक मॉडल जहां उपभोक्ता जो कंटेंट इस्तेमाल करते हैं, उसी के हिसाब से पैसे देते हैं।
\item नए आर्थिक मॉडल जहां उपभोक्ता अलग-अलग क्रिप्टोकरेंसी का इस्तेमाल करके भुगतान करते हैं।
\item नए आर्थिक मॉडल जहां उपभोक्ता अपने खुद के ब्लॉकचेन का इस्तेमाल करते हैं।
\item नए आर्थिक मॉडल जहां उपभोक्ता किसी भी कंटेंट के लिए पैसे नहीं देते।
\end{enumerate}

\bigskip


\subsection*{Question 238}
Lightning Network का स्वयं परीक्षण करने का एक तरीका क्या है?

\begin{enumerate}[label=\Alph*., leftmargin=*]
\item Mining Bitcoin द्वारा पर्सनल कंप्यूटर का उपयोग किया गया।
\item बैंक से Bitcoin खरीदकर।
\item फाउंटेन नामक पॉडकास्ट एप्लिकेशन को आज़माकर।
\item एक नई क्रिप्टोकरेंसी बनाकर।
\end{enumerate}

\bigskip


\subsection*{Question 239}
एआई टेक्नोलॉजी के विकास का एक संभावित असर क्या हो सकता है?

\begin{enumerate}[label=\Alph*., leftmargin=*]
\item AI टेक्नोलॉजीज़ से दुनिया की इकॉनमी बर्बाद हो जाएगी।
\item AI टेक्नोलॉजी इंसानी दिमाग को कमजोर कर देगी।
\item AI टेक्नोलॉजी सभी कंप्यूटरों को नष्ट कर देगी।
\item AI टेक्नोलॉजी की वजह से आज के 80% से ज़्यादा नौकरियाँ खत्म हो जाएँगी।
\end{enumerate}

\bigskip


\subsection*{Question 240}
Bitcoin का भविष्य की तकनीक और समाज से क्या संबंध है?

\begin{enumerate}[label=\Alph*., leftmargin=*]
\item जीडब्ल्यू-6 सरकारों के लिए अपने नागरिकों के वित्तीय लेन-देन पर नज़र रखने का एक तरीका है।
\item जीडब्ल्यू-5 दुनिया की आबादी को नियंत्रित करने का एक उपकरण है, जिससे गरीबी को खत्म किया जा सकता है।
\item Bitcoin एक AI टेक्नोलॉजी है जो लोगों को एक सेंट्रल सर्वर के जरिए Exchange इनफॉर्मेशन एक्सेस करने देती है।
\item Bitcoin एक तकनीकी क्रांति है जो बिना किसी भरोसेमंद तीसरे पक्ष के Exchange के मूल्य को सक्षम बनाती है।
\end{enumerate}

\bigskip


\subsection*{Question 241}
एक ऐसा मुख्य सवाल जो फाइनेंस के भविष्य को लेकर उठाया जाता है, वो है - "क्या क्रिप्टोकरेंसी और ब्लॉकचेन टेक्नोलॉजी पारंपरिक बैंकिंग सिस्टम की जगह ले लेंगे?"

\begin{enumerate}[label=\Alph*., leftmargin=*]
\item पैसे को कैसे बेकार बना सकते हैं?
\item सबके पास बराबर पैसा कैसे हो सकता है?
\item पैसा कौन रखे, मंजूरी दे और उसका पता लगाए?
\item पैसे को पूरी तरह से कैसे खत्म किया जा सकता है?
\end{enumerate}

\bigskip


\subsection*{Question 242}
Bitcoin की एक मुख्य खासियत क्या है?

\begin{enumerate}[label=\Alph*., leftmargin=*]
\item इसे सिर्फ चुने हुए लोग ही इस्तेमाल कर सकते हैं।
\item इसे सिर्फ कुछ खास देशों में ही इस्तेमाल किया जा सकता है।
\item इसे दुनिया में कहीं भी बिना किसी संस्था की अनुमति के बदला जा सकता है।
\item इसे इस्तेमाल करने के लिए केंद्रीय अधिकारी से मंजूरी चाहिए।
\end{enumerate}

\bigskip


\subsection*{Question 243}
नई टेक्नोलॉजीज जैसे Bitcoin को अपनाने का एक संभावित फायदा क्या है?

\begin{enumerate}[label=\Alph*., leftmargin=*]
\item यह generate पूरी दुनिया में बड़े पैमाने की अर्थव्यवस्थाएं ला सकता है।
\item इससे वैश्विक आर्थिक मंदी हो सकती है।
\item यह सभी नौकरियों को खत्म कर सकता है।
\item इससे इंसानी बुद्धिमत्ता में कमी आ सकती है।
\end{enumerate}

\bigskip


\subsection*{Question 244}
बैंकिंग सिस्टम को लेकर एक मुख्य चिंता क्या हो सकती है?

\begin{enumerate}[label=\Alph*., leftmargin=*]
\item बैंकिंग सिस्टम बहुत ज्यादा पारदर्शी है।
\item बैंकिंग सिस्टम समझने में बहुत आसान है।
\item बैंकिंग सिस्टम बहुत ज्यादा विकेंद्रित है।
\item बैंकिंग गेम के नियम सभी के लिए पारदर्शी और समझने योग्य नहीं हैं।
\end{enumerate}

\bigskip


\subsection*{Question 245}
सेंसरशिप से जुड़ा एक मुख्य मुद्दा क्या है?

\begin{enumerate}[label=\Alph*., leftmargin=*]
\item इससे सभा करने के अधिकार को और ज़्यादा समर्थन मिल सकता है।
\item इससे अभिव्यक्ति की स्वतंत्रता को लागू करने में मदद मिल सकती है।
\item यह अभिव्यक्ति की स्वतंत्रता और सभा के अधिकार को नष्ट कर सकता है।
\item इससे नवाचार में कमी आ सकती है।
\end{enumerate}

\bigskip


\subsection*{Question 246}
Bitcoin का एक मुख्य फायदा यह है कि यह बिना बैंक अकाउंट वाले लोगों को भी डिजिटल पेमेंट करने की सुविधा देता है।

\begin{enumerate}[label=\Alph*., leftmargin=*]
\item Bitcoin सिर्फ़ उन लोगों को फ़ायदा पहुँचाता है जिनकी सामाजिक हैसियत ऊँची है।
\item जीडब्ल्यू-3 का फायदा सिर्फ उन्हीं लोगों को मिलता है जिनके पास बैंक अकाउंट है और वामपंथी राजनीतिक राय रखते हैं।
\item जीडब्ल्यू-1 लेन-देन में समानता की अनुमति देता है, चाहे आपकी सामाजिक स्थिति या राजनीतिक पद कुछ भी हो।
\item जीडब्ल्यू-2 अमीर यूजर्स को ज्यादा फायदा देकर, लेन-देन में असमानता को बढ़ावा देता है।
\end{enumerate}

\bigskip


\subsection*{Question 247}
Bitcoin प्रोटोकॉल की एक मुख्य विशेषता क्या है?

\begin{enumerate}[label=\Alph*., leftmargin=*]
\item यह कुछ खास यूजर्स की तरफ झुका हुआ है।
\item यह सिर्फ़ कुछ ख़ास लोगों को फ़ायदा पहुँचाता है।
\item यह अपने सभी उपयोगकर्ताओं के प्रति निष्पक्ष है।
\item यह कुछ यूज़र्स को दूसरों से ज़्यादा फ़ेवर करता है।
\end{enumerate}

\bigskip


\subsection*{Question 248}
बिटकॉइन पाने का एक आसान तरीका क्या है?

\begin{enumerate}[label=\Alph*., leftmargin=*]
\item उन्हें असली दुनिया में ढूंढना।
\item उन्हें विनियमित या अविनियमित प्लेटफॉर्म के माध्यम से खरीदना।
\item जीडब्ल्यू-0 उन्हें लैपटॉप का इस्तेमाल करते हुए।
\item उन्हें एक खास सॉफ्टवेयर से बना रहे हैं।
\end{enumerate}

\bigskip


\subsection*{Question 249}
Bitcoin बनाने का एक मुख्य कारण क्या था?

\begin{enumerate}[label=\Alph*., leftmargin=*]
\item एक नया AI टेक्नोलॉजी फॉर्म बनाने का प्रस्ताव रखना।
\item पैसे को खत्म करने का प्रस्ताव देना।
\item दुनिया की आबादी को नियंत्रित करने का प्रस्ताव देना।
\item वित्तीय प्रणाली में बदलाव का प्रस्ताव रखना।
\end{enumerate}

\bigskip


\subsection*{Question 250}
बैंकिंग सिस्टम पर Bitcoin का एक मुख्य नतीजा क्या है?

\begin{enumerate}[label=\Alph*., leftmargin=*]
\item यह बैंकिंग सिस्टम में भ्रष्टाचार को बढ़ावा देता है।
\item यह बैंकिंग सिस्टम को खत्म करने की अनुमति देता है।
\item यह बैंकिंग सिस्टम को कंट्रोल करने की अनुमति देता है।
\item यह बैंकिंग सिस्टम से मुक्ति दिलाता है।
\end{enumerate}

\bigskip


\newpage
\section*{Answer Key}

\begin{enumerate}
\item \textbf{Question 1:} C

\item \textbf{Question 2:} C
\\ \textit{फीस ब्लॉक में लेन-देन को शामिल करने के लिए एक मुक्त बाजार बनाने के लिए जरूरी हैं। असल में, एक ब्लॉक का आकार 1 MB होता है (जिसे SegWit अपडेट के बाद 4MB तक बढ़ा दिया गया था), इसलिए हर ब्लॉक में "डाले" जा सकने वाले लेन-देन की संख्या कुछ हजार लेन-देन प्रति ब्लॉक तक सीमित होती है।}

\item \textbf{Question 3:} A
\\ \textit{Bitcoin एक कंप्यूटर प्रोटोकॉल है जो टेक्नोलॉजी का इस्तेमाल करता है... क्रिप्टोग्राफी और Blockchain की तरह, और एक मौद्रिक भी लेन-देन के लिए इस्तेमाल होने वाली इकाई।}

\item \textbf{Question 4:} C
\\ \textit{हायेक कहते हैं: "मैं नहीं मानता कि सरकार के हाथों से पैसा छीनने से पहले हमारे पास फिर से अच्छा पैसा होना चाहिए। अगर हम सरकार के हाथों से पैसे को हिंसक तरीके से नहीं छीन सकते, तो हम बस इतना कर सकते हैं कि किसी चालाकी या घुमावदार तरीके से कुछ ऐसा पेश करें जिसे वे रोक न सकें।"}

\item \textbf{Question 5:} B
\\ \textit{Bitcoin एक न्यूट्रल करेंसी है, यानी यह किसी के कंट्रोल में नहीं है। किसी भी संस्था या संगठन द्वारा नहीं, इसलिए हर यूज़र को एक जैसा अधिकार। यह इसके विकेंद्रीकरण का एक महत्वपूर्ण पहलू है।}

\item \textbf{Question 6:} C
\\ \textit{साइफरपंक्स लोगों का एक समूह है जो मानते थे कि क्रिप्टोग्राफी सुरक्षा के लिए एक उपकरण के रूप में काम कर सकती है सरकारों के हस्तक्षेप के खिलाफ व्यक्तिगत अधिकार और बड़ी कंपनियाँ।}

\item \textbf{Question 7:} A
\\ \textit{Cypherpunk मैनिफेस्टो 1993 में एरिक ह्यूजेस द्वारा लिखा गया था। और दावा करता है कि गोपनीयता एक मौलिक अधिकार है।}

\item \textbf{Question 8:} A
\\ \textit{डेविड चॉम ने गुमनाम इलेक्ट्रॉनिक की अवधारणा पेश की 1980 के दशक में उसने अपने प्रोजेक्ट डिजिकैश के साथ पैसा कमाया।}

\item \textbf{Question 9:} B
\\ \textit{साइबरस्पेस की स्वतंत्रता की घोषणा जॉन पेरी बार्लो ने 1996 में लिखी थी। इस दस्तावेज़ में कहा गया है कि साइबरस्पेस भौतिक क्षेत्र से अलग क्षेत्र है, इसलिए औद्योगिक सरकारें वहाँ अपना अधिकार लागू नहीं कर सकतीं।}

\item \textbf{Question 10:} D
\\ \textit{वेई दाई का बी-मनी एक गुमनाम डिजिटल करेंसी का आइडिया था। एक ऐसी मुद्रा जिसका कोई केंद्रीय नियंत्रण नहीं है, जिसने बहुत प्रेरित किया जीडब्ल्यू-2 का निर्माण। असल में, इसका जिक्र जीडब्ल्यू-2 में है व्हाइटपेपर।}

\item \textbf{Question 11:} B
\\ \textit{इसकी टेक्नोलॉजी से परे, Bitcoin एक क्रांति का प्रतीक है against traditional financial systems: the protocol}

\item \textbf{Question 12:} B
\\ \textit{पहले प्रकार की मुद्रा अक्सर ज़रूरी चीज़ों से जुड़ी होती थी माल जैसे अनाज, पशुधन, और अन्य वस्तुएं। हालांकि, इन सामानों में बड़ी खामियाँ थीं, जैसे कि नाशवानता, जिससे इन्हें एक के रूप में इस्तेमाल करना मुश्किल हो जाता है लंबे समय तक बचत का माध्यम।}

\item \textbf{Question 13:} A
\\ \textit{अरस्तू के अनुसार, पैसा तीन कार्य करता है: यह मूल्य का भंडार है, Exchange का माध्यम है, और खाते की इकाई है। आजकल, Exchange फ़ंक्शन का माध्यम अन्य दो से पहले मुख्य के रूप में व्यापक रूप से पहचाना जाता है।}

\item \textbf{Question 14:} C
\\ \textit{एक प्रभावी मुद्रा को फंजिबल (आपस में बदलने योग्य) होना चाहिए बिना मूल्य की हानि के), विभाज्य (छोटे हिस्सों में अलग किया जा सकने वाला भागों को अलग-अलग मात्रा के लेन-देन को आसान बनाने के लिए), और जीडब्ल्यू-3 (आसानी से माल या सेवाओं में परिवर्तनीय)।}

\item \textbf{Question 15:} A
\\ \textit{मुद्रा के रूप में सोने का एक मुख्य नुकसान यह है कि आवश्यकता पड़ने पर इसे छोटी इकाइयों में विभाजित करना मुश्किल होता है, जिससे इसका उपयोग कम व्यावहारिक हो जाता है।}

\item \textbf{Question 16:} B
\\ \textit{गोल्ड Exchange का स्टैंडर्ड मीडियम बन गया क्योंकि यह... rarity, durability, and divisibility: its rarity prevented}

\item \textbf{Question 17:} B
\\ \textit{पैसा एक संचार का साधन है जो भाषा और सीमाओं के पार संवाद करने में मदद करता है। जगह और वक़्त, जैसे हम अपना समय और ऊर्जा बदलते हैं... एक संपत्ति जिसे भविष्य में दोबारा इस्तेमाल किया जा सके बिना किसी जोखिम के मूल्यह्रास। यह एक सार्वभौमिक तरीके से संचार की भी अनुमति देता है। आम भाषा, जो दो अजनबियों को Exchange, व्यापार करने में सक्षम बनाती है, और चीज़ों की कीमत पर सहमत हो जाते हैं।}

\item \textbf{Question 18:} D
\\ \textit{फिएट करेंसीज का मुख्य नुकसान, जैसे कि... डॉलर और यूरो की बात यह है कि उनकी कीमत लगातार मुद्रास्फीति से खत्म हो गया। यह इस वजह से है। मूल्यह्रास जो उन्हें नियंत्रित करने वाली संस्थाओं (राजा, केंद्रीय बैंक, सम्राट, तानाशाह}

\item \textbf{Question 19:} C
\\ \textit{Bitcoin की सबसे बड़ी खासियत यह है कि यह दूसरे करेंसी फॉर्म्स से बेहतर है। यह इसलिए क्योंकि यह अपने उत्कृष्ट मूल्य संचय के कारण सख्ती से सीमित Supply और इसकी भरोसेमंद पार्टी की कमी। यह यह इसे लंबी अवधि की बचत के लिए एक अच्छा विकल्प बनाता है।}

\item \textbf{Question 20:} D
\\ \textit{सोना एक वैश्वीकृत और डिजिटल दुनिया में कदम नहीं मिला पा रहा क्योंकि इसे विभाज्य बनाने के लिए एक केंद्रीय संस्था की जरूरत होती है आसानी से बदला जा सकने वाला। इस वजह से इसे इस्तेमाल करना कम व्यावहारिक हो जाता है। एक ऐसी दुनिया जहां लेन-देन तेजी से डिजिटल होते जा रहे हैं और विकेंद्रित}

\item \textbf{Question 21:} B
\\ \textit{इसके गुणों के बावजूद, जैसे कि इसकी सख्त सीमा जीडब्ल्यू-3, जीडब्ल्यू-2 अभी भी व्यापार में व्यापक रूप से स्वीकार नहीं किया गया है आज, लोगों द्वारा अपनाने की कमी की वजह से, इसलिए... व्यापारियों को इसे स्वीकार करने के लिए प्रोत्साहित नहीं किया जाता है।}

\item \textbf{Question 22:} D
\\ \textit{मुद्रा एक ऐसी तकनीक है जो समय के साथ विकसित हुई है, कच्चे पत्थरों से दर्शाए जाने से शुरू करके, बदलते हुए... धातु के सिक्कों में, फिर नोटों में, बैंक कार्ड में, और आखिरकार पीयर-टू-पीयर नेटवर्क्स में इस्तेमाल के साथ जीडब्ल्यू-9, जीडब्ल्यू-8 और जीडब्ल्यू-7।}

\item \textbf{Question 23:} C
\\ \textit{सोना घना और भारी होता है, जिसका मतलब है कि इसे बड़ी मात्रा में ले जाना शारीरिक रूप से मुश्किल हो सकता है। हल्की चीज़ों की तुलना में इसका वज़न इसे ले जाने में कम सुविधाजनक बनाता है।}

\item \textbf{Question 24:} D
\\ \textit{फिएट करेंसीज, जैसे डॉलर और यूरो, बहुत Liquid और आसानी से ले जाने योग्य क्योंकि अब वे ज्यादातर डिजिटल। हालांकि, इनकी कीमत लगातार कम होती जा रही है। पैसों की मुद्रास्फीति।}

\item \textbf{Question 25:} D
\\ \textit{जीडब्ल्यू-1, अपने गुणों जैसे कि इसकी विभाज्यता की वजह से, सत्यापनीयता, अनुमति-रहित पहलू और परिवहनीयता, एक बेहतरीन मूल्य संचय (स्टोर ऑफ वैल्यू) प्रदान करता है। इसके अलावा, एक तटस्थ इंटरनेट करेंसी, यह Exchange के लिए एक अच्छा माध्यम है जो किसी सीमा को नहीं मानता।}

\item \textbf{Question 26:} C
\\ \textit{साउंड मनी एक प्राकृतिक सामाजिक घटना है जो अवश्य उभरनी चाहिए। बाज़ार से और स्वैच्छिक सहमति से। कोई इंसान या इंसानों का एक समूह पैसा बना सकता है।}

\item \textbf{Question 27:} A
\\ \textit{Bitcoin को पैसे का एक नया रूप माना जा सकता है, क्योंकि यह उभर रहा है एक पीयर-टू-पीयर नेटवर्क से जिसमें कोई केंद्रीय संस्था नहीं है।}

\item \textbf{Question 28:} B
\\ \textit{हालांकि यह आसानी से विभाज्य या परिवहन योग्य नहीं है, लंबी दूरियों के लिए, सोना एक बहुत ही टिकाऊ सामग्री है, जो यह एक अच्छा मूल्य संचय बनाता है। सोने के सिक्कों को एक... पहले करेंसी का इस्तेमाल होता था, लेकिन वे एक जैसा काम नहीं करते थे आज के डिजिटल युग में मकसद।}

\item \textbf{Question 29:} C
\\ \textit{फिएट मुद्रा को सरकारी आदेश द्वारा कानूनी निविदा के रूप में स्थापित किया जाता है, तथा इसका मूल्य किसी भी भौतिक वस्तु से प्राप्त नहीं होता है।}

\item \textbf{Question 30:} B
\\ \textit{फिड्यूशियरी करेंसी किसी भौतिक वस्तु (कमोडिटी) द्वारा समर्थित नहीं होती है। सोने या चाँदी की तरह। उनकी कीमत भरोसे से आती है। और लोगों का उन संस्थानों पर भरोसा जो जारी करते हैं और उन्हें नियंत्रित करो।}

\item \textbf{Question 31:} A
\\ \textit{केंद्रीय बैंक वो संस्थाएँ होती हैं जो एक मुद्रा जारी करने की ज़िम्मेदारी संभालती हैं। फिड्यूशियरी करेंसी। वे मौद्रिक नीति तय करते हैं और पता लगाएं कि कितना पैसा प्रचलन में डाला जाना चाहिए छपा हुआ।}

\item \textbf{Question 32:} D
\\ \textit{Bitcoin का कुल Supply, 21 मिलियन यूनिट्स तक सीमित है। यह Bitcoin की एक खास बात है जो इसे अलग बनाती है फ़िएट करेंसी से, जिसे बस छापा जा सकता है केंद्रीय बैंकों द्वारा असीमित मात्रा में।}

\item \textbf{Question 33:} B
\\ \textit{जीडब्ल्यू-1 को राज्य को अलग करने के उद्देश्य से बनाया गया था पैसे से। यह व्यवस्थागत चुनौतियों का एक जवाब है। वर्तमान वित्तीय प्रणाली, जहाँ राज्य का पूरा नियंत्रण है... पैसों पर कंट्रोल Supply का है और वो इसे अपने मुताबिक चला सकता है अपना फायदा।}

\item \textbf{Question 34:} D
\\ \textit{केंद्रीय संस्थाएं अक्सर मुद्रा का मूल्य घटा देती हैं उसकी आपूर्ति बढ़ाकर मुद्रा द्रव्यमान शुरुआती वाले की तुलना में कुछ प्रतिशत हर साल। इस चुपके से की गई मूल्यह्रास को अक्सर यह कहकर सही ठहराया जाता है कि... लोगों के फायदे में होना।}

\item \textbf{Question 35:} D
\\ \textit{पैसे छापना, या केंद्रीय बैंक द्वारा नया पैसा बनाना बैंकों से मुद्रास्फीति होती है। ऐसा इसलिए होता है क्योंकि बढ़ोतरी से... पैसा Supply पैसे की खरीदारी की ताकत को कम कर देता है, जिससे कीमतें बढ़ जाती हैं और मूल्य में कमी आ जाती है बचत।}

\item \textbf{Question 36:} A
\\ \textit{केंद्रीय बैंक डिजिटल मुद्राएँ (CBDCs) संभावित रूप से लोगों की आर्थिक आज़ादी में रुकावट डालना और सुविधा प्रदान करना अधिनायकवादी अत्याचार। ऐसा इसलिए क्योंकि वे पेशकश करेंगे अधिक केंद्रीय योजनाबद्ध अर्थव्यवस्था, जहां केंद्रीय बैंक पैसों पर और भी ज़्यादा कंट्रोल रखना Supply.}

\item \textbf{Question 37:} D
\\ \textit{इतिहास में, नेताओं ने सोने को केंद्रीकृत किया है और नए नियम लागू किए हैं। एक नई मुद्रा जिसकी कीमत सोने के बराबर हो। यह नई मुद्रा को धीरे-धीरे अवमूल्यन किया जाता है, अक्सर बहाने के तहत लोगों के हित में होने का।}

\item \textbf{Question 38:} A
\\ \textit{अचानक अवमूल्यन या किसी प्रत्ययी मुद्रा में विश्वास की कमी से हाइपरइन्फ्लेशन हो सकता है। यह एक ऐसी स्थिति है जिसमें वस्तुओं और सेवाओं की कीमतें अत्यधिक उच्च दर से बढ़ती हैं, जिससे अक्सर आर्थिक अस्थिरता पैदा होती है।}

\item \textbf{Question 39:} B
\\ \textit{इसके विकेंद्रीकृत और क्रिप्टोग्राफिक स्वभाव की वजह से, जीडब्ल्यू-1 को नकली नहीं बनाया जा सकता और यह सिर्फ 21 मिलियन तक ही सीमित है यूनिट्स। यह एक मुख्य विशेषता है जो इसे अलग बनाती है फिएट करेंसी, जो एक मौद्रिक नीति का पालन करती है जो... केंद्रीय बैंकों द्वारा बदला गया।}

\item \textbf{Question 40:} D
\\ \textit{जैसे कि राजनीतिक खिलाड़ी मौजूदा वित्तीय स्थिति से फायदा उठाते हैं सिस्टम में, उन्हें बड़े बदलाव करने का कोई प्रोत्साहन नहीं मिलता। यह सिस्टम को अपना काम जारी रखने देता है जब तक कि संभावित विस्फोट, क्योंकि सिस्टम नियंत्रित है और "अपने ही पतन से बचने के लिए प्रतिबंधित"}

\item \textbf{Question 41:} A
\\ \textit{हाइपरइन्फ्लेशन एक मौद्रिक घटना है जिसमें कीमतों में बेतहाशा वृद्धि होती है... मुद्रा छपाई के कारण मुद्रास्फीति में भारी बढ़ोतरी अधिकारियों द्वारा, जिससे पूरी तरह से नुकसान होता है मुद्रा में भरोसा।}

\item \textbf{Question 42:} B
\\ \textit{हाइपरइन्फ्लेशन की वजह से, लोग जल्दी-जल्दी मुद्रा में पूरी तरह से भरोसा खत्म हो जाना, जिसकी वजह से... इसकी वैल्यू में बहुत ज्यादा गिरावट आ गई कुछ ही समय में समय।}

\item \textbf{Question 43:} B
\\ \textit{हाइपरइन्फ्लेशन का पहला चरण मुद्रा में विश्वास खोना है।}

\item \textbf{Question 44:} A
\\ \textit{हाइपरइन्फ्लेशन का आखिरी चरण तब होता है जब एक नया नई करेंसी पुरानी करेंसी की जगह लेने के लिए लाई गई है।}

\item \textbf{Question 45:} D
\\ \textit{2\% मुद्रास्फीति के साथ, आप अपनी क्रय शक्ति का 2\% सालाना खो देते हैं, जो 5 वर्षों में 10\% हो जाता है।}

\item \textbf{Question 46:} A
\\ \textit{20\% मुद्रास्फीति के साथ, आप 3 वर्षों में अपनी क्रय शक्ति का लगभग आधा हिस्सा खो देते हैं।}

\item \textbf{Question 47:} B
\\ \textit{एक दिन में 100\% महंगाई दर एक यथार्थपरक परिदृश्य है जो ज़िम्बाब्वे में 2007 में हुआ था, जहाँ पैसे की कीमत लगभग... हर 24 घंटे में दोगुना हो जाता है।}

\item \textbf{Question 48:} D
\\ \textit{हाइपरइन्फ्लेशन का दूसरा चरण है मुद्रा का पतन और कीमतों में बढ़ोतरी, आत्मविश्वास खोने के चरण के ठीक बाद।}

\item \textbf{Question 49:} A
\\ \textit{जर्मनी में सबसे अधिक मुद्रास्फीति दर वाला महीना अक्टूबर 1923 था, जब यह 29,500\% के शिखर पर पहुंच गई थी।}

\item \textbf{Question 50:} C
\\ \textit{हंगरी में सबसे ज़्यादा महंगाई वाला महीना जुलाई था 1946, जब कीमतों में बेहिसाब बढ़ोतरी हुई थी 41,90,00,00,00,00,00,000\%।}

\item \textbf{Question 51:} B
\\ \textit{नवंबर 2008 में, ज़िम्बाब्वे के करेंसी की मासिक महंगाई दर 79,600,000,000\% तक पहुँच गई, जिसने ज़िम्बाब्वे डॉलर की कीमत को पूरी तरह से खत्म कर दिया (0 तक पहुँचा दिया)।}

\item \textbf{Question 52:} C
\\ \textit{अप्रैल 2009 में, वित्त मंत्री ने घोषणा की ज़िम्बाब्वे डॉलर का निलंबन और उपयोग की अनुमति दी व्यापार के लिए विभिन्न विदेशी मुद्राओं का।}

\item \textbf{Question 53:} C
\\ \textit{जी.डब्लू.-1 में एक कुशल और स्वस्थ मुद्रा होने के लिए सभी आवश्यक विशेषताएं हैं: विभाज्य, तुरंत परिवहन योग्य, बिना सेंसर, सत्यापन के लिए नगण्य लागत, और 21 मिलियन इकाइयों के साथ आने वाली शताब्दियों के लिए पहले से ही निर्धारित मौद्रिक नीति।}

\item \textbf{Question 54:} A
\\ \textit{Bitcoin पर लेन-देन को सत्यापित करने वाली प्रक्रिया नेटवर्क को Mining कहा जाता है। माइनर्स क्रिप्टोग्राफिक काम करो और नए बिटकॉइन पाने का इनाम पाओ।}

\item \textbf{Question 55:} B
\\ \textit{वह घटना जब Mining के लिए इनाम आधा कर दिया जाता है, उसे Halving कहा जाता है। यह घटना मौद्रिक निर्गम वक्र को सीढ़ी जैसा आकार देती है।}

\item \textbf{Question 56:} B
\\ \textit{वह तंत्र जो यह सुनिश्चित करता है कि एक नया ब्लॉक जोड़ा जाए जीडब्ल्यू-1 हर दस मिनट में डिफिकल्टी एडजस्टमेंट होता है। यह एडजस्टमेंट हर 2016 ब्लॉक्स के बाद होता है, यानी लगभग दो हफ़्ते।}

\item \textbf{Question 57:} B
\\ \textit{गेम थ्योरी एक गणितीय अवधारणा है जो आधारित है इंसानी तर्कशक्ति। Bitcoin में, प्रोटोकॉल में बदलाव होते हैं उपयोगकर्ताओं द्वारा लागू किया गया, Bitcoin में किसी भी बदलाव को करना प्रोटोकॉल बहुत जटिल है।}

\item \textbf{Question 58:} D
\\ \textit{बिटकॉइन की मात्रा को सत्यापित करने का कमांड Bitcoin नोड पर सर्कुलेशन bitcoin-cli है। gettxoutsetinfo. यह कमांड पारदर्शिता और सत्यापन योग्यता, Bitcoin प्रणाली में भरोसा मजबूत करना।}

\item \textbf{Question 59:} C
\\ \textit{जीडब्ल्यू-1 ऑस्ट्रियन इकोनॉमिक्स के सिद्धांतों के साथ मेल खाता है। इसकी सीमित मात्रा और केंद्रीय नियंत्रण की कमी इसे सुरक्षित रखती है मुद्रास्फीति के अंतर्निहित जोखिमों से जो देखने को मिलते हैं पारंपरिक मुद्राएँ।}

\item \textbf{Question 60:} A
\\ \textit{Halving मैकेनिज्म की वजह से, इसे गणितीय तरीके से यह अनुमान लगाया गया है कि बिटकॉइन का निर्माण 2140 में बंद हो जाएगा। जब बिटकॉइन की कुल संख्या अपनी सीमा 21 मिलियन तक पहुँच जाएगी दस लाख।}

\item \textbf{Question 61:} C
\\ \textit{कठिनाई समायोजन एक ऐसी प्रक्रिया है जो होती है हर 2016 ब्लॉक्स पर, या लगभग दो हफ्ते में, यह सुनिश्चित करने के लिए कि हर दस मिनट में Blockchain में एक नया ब्लॉक जोड़ा जाता है औसत, चाहे दुनिया भर में Hashrate कुछ भी हो।}

\item \textbf{Question 62:} A
\\ \textit{Mining का इनाम हर 210,000 पर आधा होने के लिए प्रोग्राम किया गया है। blocks, which is approximately every four years: an event}

\item \textbf{Question 63:} D
\\ \textit{पहले Halving के बाद, जो ब्लॉक हाइट पर हुआ था 2,10,000, अनुमानित संख्या में चलन में मौजूद बिटकॉइन की 10,500,000 BTC था, जो कुल संख्या का आधा है सर्कुलेशन में आने वाले बिटकॉइन्स।}

\item \textbf{Question 64:} B
\\ \textit{चौथे Halving के बाद, जो ब्लॉक हाइट पर होता है 8,40,000, Mining के लिए BTC इनाम है 3.125 BTC।}

\item \textbf{Question 65:} A
\\ \textit{7वां Halving 2036 में होना चाहिए, और अनुमानित... Bitcoin की प्रचलन में संख्या 20,835,937.5 होगी। जो कुल का लगभग 99.22\% होगा जीडब्ल्यू-3}

\item \textbf{Question 66:} B
\\ \textit{Bitcoin Wallet का इस्तेमाल Bitcoin के साथ इंटरैक्ट करने के लिए किया जाता है। नेटवर्क। इसके तीन मुख्य कार्य हैं: डेटा प्राप्त करने की अनुमति देना, बिटकॉइन्स, बिटकॉइन्स भेजने की अनुमति देते हैं, और उन्हें सुरक्षित रखते हैं हैकिंग और चोरी की कोशिशें।}

\item \textbf{Question 67:} D

\item \textbf{Question 68:} C
\\ \textit{हालांकि आपकी प्राइवेट कुंजियाँ आपके Wallet में स्टोर हैं, लेकिन बिटकॉइन असल में Bitcoin में स्टोर होते हैं। जीडब्ल्यू-7, जो कि जनता के बीच वितरित जीडब्ल्यू-8 है जीडब्ल्यू-9 पीयर-टू-पीयर नेटवर्क।}

\item \textbf{Question 69:} B
\\ \textit{2017 से, प्राइवेट की को एक साधारण लिस्ट में कोड किया जा सकता है। 12 या 24 शब्दों का, जिसे Mnemonic वाक्यांश कहा जाता है। यह वाक्यांश क्या आपके Bitcoin और Wallet का बैकअप है, तो यह आपको अनुमति देता है अपने बिटकॉइन्स को एक्सेस करने के लिए किसी भी Bitcoin या Wallet का उपयोग करें सॉफ़्टवेयर/ऐप}

\item \textbf{Question 70:} D
\\ \textit{पब्लिक की प्राइवेट की से बनती है और इसका इस्तेमाल generate और Bitcoin पते पर पैसे प्राप्त करने के लिए।}

\item \textbf{Question 71:} A
\\ \textit{सार्वजनिक कुंजी साझा करने से गोपनीयता को खतरा हो सकता है, लेकिन... सुरक्षा के लिए। ऐसा इसलिए है क्योंकि पब्लिक की का इस्तेमाल जीडब्ल्यू-6 जीडब्ल्यू-7 पते और इस तरह पैसा प्राप्त करते हैं, लेकिन यह बिटकॉइन के Ownership को प्रतिनिधित्व नहीं करता है।}

\item \textbf{Question 72:} B
\\ \textit{गवाह-8 वाला डिवाइस खो जाने पर ज़रूरी नहीं कि इसका मतलब आपके बिटकॉइन का नुकसान हो। क्या चीज़ आपको अनुमति देती है अपने Wallet को दोबारा बनाने और अपने Bitcoin को खर्च करने का मतलब है जीडब्ल्यू-11 वाक्यांश।}

\item \textbf{Question 73:} C
\\ \textit{अगर आप अपने डिफ़ॉल्ट सुरक्षा से संतुष्ट नहीं हैं तो Bitcoin Wallet, आप इसे हमेशा मजबूत बना सकते हैं... जीडब्ल्यू-14 से आपके जीडब्ल्यू-15 जीडब्ल्यू-16।}

\item \textbf{Question 74:} A
\\ \textit{Wallet बनाने के लिए इस्तेमाल की गई क्रिप्टोग्राफी की बदौलत, किसी के द्वारा गलती से आपके 12 या 24 शब्दों की सूची का अनुमान लगाने की संभावना खगोलीय रूप से कम है। यह 1 और $2^256$ के बीच सही संख्या खोजने जैसा है, जो ब्रह्मांड में एक परिभाषित परमाणु को खोजने के लगभग बराबर है।}

\item \textbf{Question 75:} B

\item \textbf{Question 76:} D
\\ \textit{बिटकॉइन खर्च करने के लिए मालिक को एलिप्टिक कर्व का इस्तेमाल करना पड़ता है। डिजिटल सिग्नेचर अल्गोरिदम (ECDSA) का इस्तेमाल करके सिग्नेचर बनाना Bitcoin लेन-देन के प्रसार के लिए।}

\item \textbf{Question 77:} D
\\ \textit{बिटकॉइन प्राप्त करने के तरीके को समझाने के लिए जो उदाहरण दिया गया है भेजने वाले को आपके फंड चुराने की अनुमति दिए बिना डाकघर। लोग इसमें पैसे जमा करते हैं, लेकिन सिर्फ तुम ही जो इसे खोल सके।}

\item \textbf{Question 78:} B
\\ \textit{जब आपके पास बिटकॉइन होते हैं, तो सबसे बड़ी चिंता सुरक्षा की होती है। आपके फंड्स का। ऐसा इसलिए है क्योंकि Bitcoin, किसी भी अन्य तरह की तरह पैसे की तरह, अगर ठीक से सुरक्षित नहीं किया गया तो चोरी हो सकता है या खो सकता है।}

\item \textbf{Question 79:} D
\\ \textit{Hot Wallet एक शब्द है जिसका उपयोग Bitcoin Wallet का वर्णन करने के लिए किया जाता है जिसे इंटरनेट एक्सेस वाले डिवाइस जैसे फ़ोन या कंप्यूटर पर संग्रहीत किया जाता है। इस तरह, लेन-देन करना आसान है, लेकिन इसका मतलब यह भी है कि Wallet ऑनलाइन खतरों के लिए संभावित रूप से असुरक्षित है।}

\item \textbf{Question 80:} A
\\ \textit{Cold Wallet एक ऐसा शब्द है जिसका इस्तेमाल Bitcoin Wallet को बताने के लिए किया जाता है। जो एक ऐसे डिवाइस पर स्टोर है जो कनेक्टेड नहीं है इंटरनेट। इससे यह ऑनलाइन खतरों के प्रति कम संवेदनशील हो जाता है, लेकिन इसका मतलब यह भी है कि Wallet तक पहुँचना थोड़ा मुश्किल हो जाता है रोज़ाना लेन-देन।}

\item \textbf{Question 81:} A
\\ \textit{Multisig Wallet एक प्रकार का Bitcoin Wallet है जिसमें लेनदेन करने के लिए कई हस्ताक्षरों की आवश्यकता होती है। यह सुरक्षा का एक अतिरिक्त Layer जोड़ता है, क्योंकि इसका मतलब है कि एक भी व्यक्ति दूसरों की मंजूरी के बिना लेनदेन नहीं कर सकता है।}

\item \textbf{Question 82:} B

\item \textbf{Question 83:} B
\\ \textit{अपनी सुरक्षा और पहुंच को ज़्यादा पेचीदा बना लेना बिटकॉइन आपको नुकसान पहुंचा सकते हैं अगर, उदाहरण के लिए, आप उन्हें गलत तरीके से हैंडल करें अपने वॉलेट्स के बैकअप। ऐसा इसलिए क्योंकि साधारण बैकअप प्रक्रियाएं आपके धन तक फिर से पहुंच पाने के लिए बेहद ज़रूरी हैं अगर आपका फोन या कंप्यूटर खो जाए तो।}

\item \textbf{Question 84:} C
\\ \textit{एक मोबाइल जीडब्ल्यू-1 रोज़मर्रा के खर्चों के लिए सुझाया जाता है क्योंकि इसके साथ लेन-देन करना आसान है। हालांकि, यह बड़ी मात्रा में स्टोर करने के लिए या के लिए सुझावित नहीं है लंबे समय तक बचत क्योंकि यह संभावित रूप से कमजोर है ऑनलाइन धमकियाँ।}

\item \textbf{Question 85:} C
\\ \textit{ऑफलाइन, या Cold, फिजिकल Wallet की सिफारिश की जाती है बड़ी मात्रा में स्टोर करना क्योंकि यह कम असुरक्षित है ऑनलाइन खतरों से Hot वॉलेट्स की तुलना में ज्यादा सुरक्षित है। हालांकि, यह पूरी तरह से... रोज़मर्रा के खर्चों के लिए सुझाया गया है क्योंकि यह कम है रोज़मर्रा के लेन-देन के लिए सुलभ।}

\item \textbf{Question 86:} B
\\ \textit{एक जीडब्ल्यू-2 जीडब्ल्यू-3 बड़ी मात्रा में स्टोर करने के लिए रेकमेंडेड है और कॉर्पोरेट इस्तेमाल के लिए क्योंकि इसमें कई चीज़ों की ज़रूरत पड़ती है लेन-देन करने के लिए हस्ताक्षर। यह प्रक्रिया जोड़ती है एक्स्ट्रा जीडब्ल्यू-4 सिक्योरिटी की, क्योंकि एक आदमी अकेला नहीं कर सकता दूसरों की मंजूरी के बिना लेन-देन करना।}

\item \textbf{Question 87:} A
\\ \textit{जब आप Wallet को passphrase के साथ इस्तेमाल करते हैं, तो आपको बैकअप लेना होगा। 12 या 24 शब्दों की सूची और आपका passphrase, दोनों क्योंकि शब्दों की सूची और passphrase दोनों अपने फंड तक फिर से पहुंच पाने के लिए जरूरी अपना फोन या कंप्यूटर खो देना। एक के बिना या दूसरे के बिना, फंड्स एक्सेस नहीं हो सकते।}

\item \textbf{Question 88:} A
\\ \textit{जब आप Multisig या Wallet का इस्तेमाल करते हैं, तो हर एक प्राइवेट की का जीडब्ल्यू-4 को अलग-अलग जगहों पर स्टोर करना चाहिए। यह क्योंकि Multisig को चलाने के लिए कई सिग्नेचर चाहिए होते हैं एक लेन-देन करने के लिए, तो एक निश्चित संख्या में वॉलेट्स होते हैं ज़रूरत है। इसलिए, हर हिस्से को अलग-अलग जगह पर स्टोर करना। एक्स्ट्रा Layer सुरक्षा जोड़ देता है।}

\item \textbf{Question 89:} D
\\ \textit{HODL शब्द की उत्पत्ति 2013 में Bitcoin फ़ोरम पर एक गलत वर्तनी वाली ऑनलाइन पोस्ट से हुई थी। समय के साथ, HODL एक व्यापक निवेश रणनीति का प्रतिनिधित्व करने के लिए विकसित हुआ है। यह निवेशकों को बाज़ार में उतार-चढ़ाव की परवाह किए बिना, लंबे समय तक अपनी क्रिप्टोकरेंसी को बनाए रखने के लिए प्रोत्साहित करता है।}

\item \textbf{Question 90:} B
\\ \textit{Bitcoin Wallet में, निजी कुंजी को अक्सर 12 या 24 शब्दों की सूची द्वारा दर्शाया जाता है, जिसे seed वाक्यांश या Mnemonic वाक्यांश के रूप में भी जाना जाता है।}

\item \textbf{Question 91:} A
\\ \textit{अगर आपका प्राइवेट की किसी तीसरे पक्ष के हाथ लग जाए, तो जुड़े हुए फंड अब सुरक्षित नहीं हैं, क्योंकि प्राइवेट... कुंजी उन्हें आपके फंड तक पहुंचने की अनुमति देती है।}

\item \textbf{Question 92:} C
\\ \textit{कागज़ के बदले स्टोर करने के लिए एक सुझाया गया विकल्प है तुम्हारा Mnemonic वाक्यांश एक धातु की प्लेट पर उकेरा जा रहा है। समाधान ज़्यादा टिकाऊ और लंबे समय तक चलने वाला है।}

\item \textbf{Question 93:} D
\\ \textit{एक बार आपके Mnemonic वाक्यांश के शब्द लिखे जाने के बाद, यह दूसरी कॉपी बनाने और इसे स्टोर करने की सलाह दी जाती है पहले वाले स्थान से अलग जगह। इससे बैकअप मिलता है। पहली कॉपी के खो जाने या दुर्घटना होने की स्थिति में।}

\item \textbf{Question 94:} C
\\ \textit{12 या 24 शब्दों की सूची का न होना, यही है जीडब्ल्यू-7 वाक्यांश, जीडब्ल्यू-8 में आपको सचेत करना चाहिए कि जीडब्ल्यू-9 सुरक्षित नहीं है क्योंकि यह स्टैंडर्ड बैकअप प्रोटोकॉल का पालन नहीं करता प्रक्रियाएँ।}

\item \textbf{Question 95:} A
\\ \textit{प्राइवेट की को बैकअप करने का स्टैंडर्ड तरीका है Bitcoin Wallet आपके Mnemonic वाक्यांश को किसी में डालकर है जीडब्ल्यू-9 सॉफ्टवेयर। यह आपको अपने जीडब्ल्यू-9 को रिस्टोर करने की सुविधा देता है।}

\item \textbf{Question 96:} C
\\ \textit{अपने Mnemonic वाक्यांश को एक मजबूत सामग्री जैसे कि पर उकेरना स्टील आपकी चाबियों का एक भौतिक बैकअप बनाता है जो पानी के नुकसान और आग दोनों से सुरक्षित, इस तरह आपका बचाव करता है लॉन्ग टर्म में बिटकॉइन।}

\item \textbf{Question 97:} D
\\ \textit{एक वारिस योजना यह सुनिश्चित करती है कि आपके बिटकॉइन्स आपकी मृत्यु के बाद सही तरीके से संभाला जाए। इसमें शामिल हो सकता है हाथ से लिखा हुआ पत्र जिसमें आपकी संपत्तियों और उन तक पहुंच का विवरण हो तरीके, और भरोसेमंद लोगों का संपर्क विवरण संपर्क किया जाना।}

\item \textbf{Question 98:} C
\\ \textit{Mnemonic वाक्यांश Bitcoin के संदर्भ में 12 की एक सूची है या 24 शब्द जो आपको किसी भी डिवाइस पर अपना Wallet रिस्टोर करने की अनुमति देते हैं जीडब्ल्यू-8 जीडब्ल्यू-9 एप्लीकेशन। इस लिस्ट तक पहुंच रखने वाला कोई भी तुम अपने बिटकॉइन्स को भी एक्सेस कर सकते हो।}

\item \textbf{Question 99:} B
\\ \textit{Bitcoin के संदर्भ में निजता महत्वपूर्ण है ताकि इसे कम से कम किया जा सके। आपके फंड्स के ट्रेस होने के खतरे, जिससे आसानी हो सकती है तुम्हारी निगरानी।}

\item \textbf{Question 100:} A
\\ \textit{Exchange पर अपने Bitcoins को छोड़ने से बचना ही बेहतर है। प्लेटफॉर्म क्योंकि ये हैकर्स के लिए आसान शिकार हो सकते हैं हमलों।}

\item \textbf{Question 101:} B
\\ \textit{Bitcoin वंशानुक्रम के संदर्भ में, यह महत्वपूर्ण है कि... बिटकॉइन के वारिस होने के बारे में एक नोटरी से बात करें ताकि यह सुनिश्चित हो सके टैक्स अनुपालन (Tax Compliance)}

\item \textbf{Question 102:} D
\\ \textit{जीडब्ल्यू-2 वॉलेट्स एक तरह का सॉफ्टवेयर है जिसका इस्तेमाल किया जाता है वो प्राइवेट कीज़ स्टोर करो जो तुम्हें तुम्हारे लेन-देन करने की अनुमति देती हैं बिटकॉइन्स।}

\item \textbf{Question 103:} C
\\ \textit{Bitcoin में, वित्तीय संप्रभुता साथ-साथ चलती है खुद की जिम्मेदारी, क्योंकि तुम्हें अपना ख्याल खुद रखना होगा अपने फंड्स को सुरक्षित रखने के लिए सुझाए गए बैकअप प्रक्रियाओं का पालन करें।}

\item \textbf{Question 104:} C
\\ \textit{ब्लॉकमिट एक कम खर्चीला तरीका है जिससे आप अपनी चीज़ों पर नक्काशी कर सकते हैं। Mnemonic वाक्यांश एक प्रतिरोधी सामग्री जैसे स्टील पर, इस तरह लंबे समय तक अपने बिटकॉइन को सुरक्षित रखना।}

\item \textbf{Question 105:} D
\\ \textit{KYC यानी "अपने ग्राहक को जानें" एक ऐसी प्रक्रिया है जो व्यवसाय अपने ग्राहकों की पहचान सत्यापित करने के लिए अपनाते हैं। यह आमतौर पर अनुपालन सुनिश्चित करने के लिए किया जाता है नियम और गैरकानूनी गतिविधियों जैसे कि कालेधन को रोकने के लिए धोखाधड़ी या फ्रॉड।}

\item \textbf{Question 106:} A
\\ \textit{नॉन-केवाईसी बिटकॉइन खरीदने का मतलब है बिना स्टैंडर्ड "नो योर कस्टमर" प्रक्रिया से गुजरना कई एक्सचेंज और प्लेटफॉर्म के वेरिफिकेशन प्रोसेस एक्सक्लूसिवली एक मुख्य वजह जो लोगों को यह चुनने के लिए प्रेरित कर सकती है रास्ता यह है कि वे अपनी प्राइवेसी बनाए रखें, यह सुनिश्चित करते हुए कि उनका... व्यक्तिगत विवरण और वित्तीय गतिविधियाँ सुरक्षित रहती हैं अनिर्दिष्ट।}

\item \textbf{Question 107:} A
\\ \textit{जहाँ एक व्यक्ति सुरक्षा को प्राथमिकता दे सकता है, वहीं दूसरा... उपयोगकर्ता-मित्रता या विशेष सुविधाओं को महत्व दें। इसलिए, एक एकल Bitcoin Wallet सबके लिए पर्याप्त नहीं हो सकता आवश्यकताएं, जिससे कई Wallet की जरूरत पड़ती है मौजूद होने के विकल्प।}

\item \textbf{Question 108:} B
\\ \textit{Bitcoin को पहली बार दुनिया के सामने 31 अक्टूबर को पेश किया गया था। 2008, जब Satoshi नाकामोटो ने Cypherpunk को एक ईमेल भेजा था जीडब्ल्यू-4 व्हाइट पेपर वाली मेलिंग लिस्ट।}

\item \textbf{Question 109:} C
\\ \textit{हाल फिन्नी Bitcoin नेटवर्क में शामिल होने वाले पहले व्यक्ति थे। Satoshi नाकामोतो खुद के बाद। उसने ट्वीट भी किया "Running" जीडब्ल्यू-4 इस अवसर को चिह्नित करने के लिए।}

\item \textbf{Question 110:} D
\\ \textit{पहला Bitcoin लेन-देन जो रिकॉर्ड किया गया था, वह एक लेन-देन था 10 बीटीसी Satoshi नाकामोटो और हाल फिन्नी के बीच।}

\item \textbf{Question 111:} D
\\ \textit{जीडब्ल्यू-1 नाकामोटो 2010 में गायब हो गए, अपनी रचना छोड़कर समुदाय के हाथों में।}

\item \textbf{Question 112:} C
\\ \textit{Bitcoin और Blockchain का पहला ब्लॉक, जिसे Genesis ब्लॉक में, मैसेज है 03/jan/2009 चांसलर बैंकों के लिए दूसरे बेलआउट के कगार पर, जो कि उस दिन का समय जीडब्ल्यू-6 से जीडब्ल्यू-7 ब्लॉक तक।}

\item \textbf{Question 113:} B
\\ \textit{बिटकॉइनटॉक फोरम Bitcoin पर चर्चा करने की एक जगह है। यूजर्स। इसे Satoshi नाकामोटो ने 22 नवंबर को बनाया था। 2009.}

\item \textbf{Question 114:} A
\\ \textit{पहली बार जब Bitcoin का इस्तेमाल भौतिक सामान खरीदने के लिए किया गया माल तब था जब लास्ज़लो हैनयेक्ज़ ने 10,000 बिटकॉइन में 2 पिज़्ज़ा खरीदे थे। बीटीसी, 2010/05/22 को जिसे अब पिज़्ज़ा डे के नाम से जाना जाता है}

\item \textbf{Question 115:} B
\\ \textit{2009 से Bitcoin की ऑनलाइन उपलब्धता 99.988\% है।}

\item \textbf{Question 116:} D
\\ \textit{2009 से Bitcoin की ऑनलाइन उपलब्धता 99.988\% है।}

\item \textbf{Question 117:} D
\\ \textit{पहला Bitcoin लेन-देन, जो एक लेन-देन था 10 BTC Satoshi नाकामोटो और हाल फिन्नी के बीच, रिकॉर्ड किया गया है 170वें ब्लॉक में।}

\item \textbf{Question 118:} B

\item \textbf{Question 119:} D
\\ \textit{Satoshi का आखिरी ज्ञात निजी संदेश की तारीख नाकामोटो ने ईमेल के जरिए 23 अप्रैल, 2011 को भेजा था। यह भेजा गया था डेवलपर माइक हर्न।}

\item \textbf{Question 120:} D
\\ \textit{Bitcoin लेन-देन डिजिटल तरीके से पैसे ट्रांसफर करने का एक तरीका है जीडब्ल्यू-8 बिटकॉइन्स का, जीडब्ल्यू-10 जीडब्ल्यू-9 के इस्तेमाल से यह एलिप्टिक कर्व क्रिप्टोग्राफी की वजह से संभव हुआ है। जीडब्ल्यू-11 उस व्यक्ति के जीडब्ल्यू-12 के लिए अनोखा है जो इसे प्राप्त कर रहा है बिटकॉइन्स।}

\item \textbf{Question 121:} C
\\ \textit{Bitcoin लेन-देन में लेन-देन शुल्क एक माइनर्स को अगले ब्लॉक में लेन-देन को शामिल करने के लिए प्रोत्साहन ब्लॉक. फीसों से इसमें शामिल होने के लिए एक मुक्त बाजार बनता है। ब्लॉक में लेन-देन।}

\item \textbf{Question 122:} A
\\ \textit{बिटकॉइन का सही तरीके से प्रबंधन करने के लिए, अपनी मृत्यु के बाद इन्हें कैसे हैंडल किया जाए, इसकी योजना बनाना बेहद जरूरी है। इस योजना में आपके संपत्तियों की सूची वाला एक हाथ से लिखा पत्र शामिल हो सकता है, उनके पहुंच तरीके, और विश्वसनीय व्यक्तियों का संपर्क विवरण जिनसे संपर्क किया जा सके}

\item \textbf{Question 123:} C

\item \textbf{Question 124:} C

\item \textbf{Question 125:} B
\\ \textit{Bitcoin नेटवर्क में एक Bitcoin नोड संभालता है नए ब्लॉक्स और शामिल लेनदेन की जाँच (वैलिडेशन)। यह है यह Bitcoin नोड्स का यह नेटवर्क ही है जो इसे सच में डिसेंट्रलाइज़्ड (विकेंद्रित)}

\item \textbf{Question 126:} D
\\ \textit{Bitcoin नोड्स का भौगोलिक वितरण इसे बनाता है सभी कॉपीज़ को नष्ट करना लगभग नामुमकिन है जीडब्ल्यू-5. चूंकि नोड्स फैले हुए हैं, इसलिए यह मुश्किल है सबको ज़बरदस्ती पकड़ लो।}

\item \textbf{Question 127:} C

\item \textbf{Question 128:} C
\\ \textit{Bitcoin नेटवर्क में ब्लॉक Mining के संदर्भ में, एक वैध Hash एक अनोखी फिंगरप्रिंट है जो इससे जुड़ी हुई है। ब्लॉक, जिसमें 256 अक्षर होते हैं। इसकी वैधता Hash, Bitcoin नेटवर्क की मुश्किलों पर निर्भर करता है।}

\item \textbf{Question 129:} D

\item \textbf{Question 130:} D

\item \textbf{Question 131:} B
\\ \textit{जीडब्ल्यू-5 नेटवर्क के यूजर्स अपने जीडब्ल्यू-4 को ट्रांसफर करते हैं पैसे अपने प्राइवेट की से डिजिटली ट्रांजैक्शन्स पर साइन करके चाबियाँ। इससे पुष्टि होती है कि वे इसके मालिक हैं। वो बिटकॉइन जो वो ट्रांसफर करना चाहते हैं।}

\item \textbf{Question 132:} D

\item \textbf{Question 133:} D

\item \textbf{Question 134:} B
\\ \textit{जीडब्ल्यू-3 जीडब्ल्यू-2 का आकार 2023 तक लगभग 500GB है।}

\item \textbf{Question 135:} C
\\ \textit{रास्पबेरी पाई 4 पर Bitcoin नोड चलाने में थोड़ा खर्च आता है सालाना बिजली खपत के मामले में 10€ से कम २०२३ तक।}

\item \textbf{Question 136:} D

\item \textbf{Question 137:} B
\\ \textit{अगर एक नोड यूज़र ने दूसरे नियमों को फॉलो करने का फैसला किया तो वर्तमान सहमति नियमों को अमान्य कर दें, फिर यह नोड अब Bitcoin नेटवर्क का हिस्सा नहीं रहेंगे।}

\item \textbf{Question 138:} D

\item \textbf{Question 139:} A
\\ \textit{2017 का ब्लॉक वॉर एक झगड़ा था जिसमें सवाल यह था कि क्या... Bitcoin को ब्लॉक साइज़ बढ़ाकर विस्तारित करें लेन-देन की क्षमता।}

\item \textbf{Question 140:} B
\\ \textit{Bitcoin नोड की सस्ती कीमत और आसान उपलब्धता है महत्वपूर्ण है क्योंकि यह विकेंद्रीकरण को आसान बनाता है नेटवर्क।}

\item \textbf{Question 141:} C
\\ \textit{अगर ब्लॉक्स 100 गुना ज़्यादा भारी होते, तो Bitcoin चलाना... नोड को 50TB Hard डिस्क की ज़रूरत होगी, साथ ही 50TB से ज़्यादा की बैंडविड्थ चाहिए होगी। 500GB/महीना, और हार्डवेयर जो सैकड़ों को वैलिडेट करने में सक्षम हो 10 मिनट से भी कम समय में हजारों लेन-देन। यह Bitcoin नोड चलाना मुश्किल बना देगा आम आदमी, विकेंद्रीकरण को कमजोर करते हुए प्रोटोकॉल और लेन-देन की अपरिवर्तनीयता तथा सहमति नियम}

\item \textbf{Question 142:} D

\item \textbf{Question 143:} C
\\ \textit{जीडब्ल्यू-7 एक तुरंत जीडब्ल्यू-8 पेमेंट नेटवर्क है यह Bitcoin Blockchain पर आधारित है।}

\item \textbf{Question 144:} B

\item \textbf{Question 145:} B

\item \textbf{Question 146:} C

\item \textbf{Question 147:} D

\item \textbf{Question 148:} C
\\ \textit{ASICs (एप्लिकेशन-स्पेसिफिक इंटीग्रेटेड सर्किट्स) होते हैं Bitcoin के लिए विशेष रूप से डिज़ाइन किया गया विशेष हार्डवेयर जीडब्ल्यू-6. ये SHA256 हैशिंग करने के लिए ऑप्टिमाइज़्ड हैं फंक्शन, जो माइनर्स के लिए एक गणितीय समस्या है Proof of Work प्रक्रिया में हल करना होगा।}

\item \textbf{Question 149:} A

\item \textbf{Question 150:} A

\item \textbf{Question 151:} B
\\ \textit{बाइज़ेंटाइन जनरल्स प्रॉब्लम कंप्यूटर साइंस में एक मशहूर समस्या है जो डिस्ट्रिब्यूटेड कंप्यूटिंग, जिसमें सहमति बनाना शामिल है एक नेटवर्क में जहां कुछ नोड्स खराब हो सकते हैं या दुर्भावनापूर्ण। Bitcoin इस समस्या का समाधान एक के उपयोग से करता है जीडब्ल्यू-2, जीडब्ल्यू-1 पर आधारित है, जो यह सुनिश्चित करता है कि सभी लेन-देन के इतिहास की स्थिति पर नोड्स सहमत होते हैं।}

\item \textbf{Question 152:} B

\item \textbf{Question 153:} B

\item \textbf{Question 154:} A
\\ \textit{51\% हमले में, एक एजेंट जिसके पास आधे से ज़्यादा नेटवर्क Hashrate संभावित रूप से हेरफेर कर सकता है जीडब्ल्यू-7 को ब्लॉकों की एक वैकल्पिक श्रृंखला बनाकर। हालांकि, ऐसा हमला बेहद महंगा होगा और फायदेमंद होने की संभावना नहीं है, इसलिए यह स्थिति होने की उम्मीद कम है।}

\item \textbf{Question 155:} A

\item \textbf{Question 156:} D
\\ \textit{Bitcoin प्रोटोकॉल में माइनर्स वर्तमान और भविष्य दोनों को सुरक्षित करते हैं। ऐतिहासिक नेटवर्क। वे कीमत को नियंत्रित नहीं करते जीडब्ल्यू-5, अपने जीडब्ल्यू-6 को कंट्रोल नहीं कर सकता, न ही लेन-देन मैनेज कर सकता है उपयोगकर्ताओं के बीच}

\item \textbf{Question 157:} B

\item \textbf{Question 158:} A

\item \textbf{Question 159:} D
\\ \textit{Bitcoin प्रोटोकॉल यह सुनिश्चित करता है कि औसत समय के बीच दो ब्लॉक्स की संख्या को स्थिर रखा जाता है, डिफिकल्टी को एडजस्ट करके। हर 2016 ब्लॉक्स पर एक वैध Hash ढूंढना। यह घटना होती है चाहे जितने भी माइनर्स हों और उनकी कंप्यूटिंग क्षमता कैसी भी हो ताक़त।}

\item \textbf{Question 160:} A
\\ \textit{माइनर्स बिजलीघरों के लिए आखिरी उपाय के खरीदार के रूप में काम कर सकते हैं, जो... अभी तक ग्रिड से नहीं जुड़े हैं। इससे पावर प्लांट्स को वित्तीय सहायता हासिल करने के लिए, नेटवर्क से जुड़ने से पहले भी इलेक्ट्रिकल नेटवर्क।}

\item \textbf{Question 161:} D
\\ \textit{वर्तमान वित्तीय प्रणाली की आलोचना की जाती है क्योंकि यह प्रोत्साहित करती है ज़रूरत से ज़्यादा खर्च और कर्ज़, जो गंभीर ख़तरा पैदा करता है... पर्यावरणीय समस्याएं। यह आसान पहुंच के कारण होता है। क्रेडिट, बैंकों द्वारा मुद्रा जारी करना, और उपयोग आंशिक रिज़र्व बैंकिंग।}

\item \textbf{Question 162:} C

\item \textbf{Question 163:} D
\\ \textit{Bitcoin प्रोटोकॉल Miner को इनाम देता है जो एक वैध खोजता है ग्वे-7 अगले ब्लॉक के लिए, जिसका इनाम राशि है सहमति नियमों द्वारा परिभाषित, साथ ही फीस से वैध ब्लॉक में शामिल सभी लेन-देन।}

\item \textbf{Question 164:} B

\item \textbf{Question 165:} B
\\ \textit{Bitcoin प्रोटोकॉल यह सुनिश्चित करता है कि इसमें सेंसरशिप का विरोध हो। लचीलापन क्योंकि प्रोटोकॉल का हर हिस्सा दुनिया भर में भौगोलिक रूप से फैला हुआ। उदाहरण के लिए, लगभग 40,000 Bitcoin नोड्स सभी जगह मौजूद हैं महाद्वीप।}

\item \textbf{Question 166:} C

\item \textbf{Question 167:} C

\item \textbf{Question 168:} B

\item \textbf{Question 169:} C

\item \textbf{Question 170:} A

\item \textbf{Question 171:} B

\item \textbf{Question 172:} B

\item \textbf{Question 173:} D

\item \textbf{Question 174:} C

\item \textbf{Question 175:} C
\\ \textit{Bitcoin अभी भी जिंदा है और और भी ज्यादा बढ़ रहा है पारंपरिक में खुद को और अधिक एकीकृत करते हुए बाजार। Bitcoin ETFs का आगमन, स्पष्ट नियमन, और स्टोरेज टूल्स का बेहतर अधिग्रहण सिर्फ इसे और बढ़ावा देता है ट्रेंड।}

\item \textbf{Question 176:} A

\item \textbf{Question 177:} C

\item \textbf{Question 178:} D

\item \textbf{Question 179:} D

\item \textbf{Question 180:} B

\item \textbf{Question 181:} D

\item \textbf{Question 182:} A

\item \textbf{Question 183:} D

\item \textbf{Question 184:} C
\\ \textit{AMF, यानी Autorité des Marchés Financiers, एक फ्रेंच बिटकॉइन को रेगुलेट करने वाला संगठन खरीदा।}

\item \textbf{Question 185:} D

\item \textbf{Question 186:} A

\item \textbf{Question 187:} C

\item \textbf{Question 188:} D
\\ \textit{BTCpay Server एक टूल है जो बड़े पैमाने पर Bitcoin स्वीकार करने के लिए है। संरचनाएँ या जोशीले बिटकॉइन समर्थक। अन्य समाधान ओपननोड को सरल ऑनलाइन व्यवसायों और स्विस के लिए शामिल करें जीडब्ल्यू-7 शौकिया व्यापारियों के लिए भुगतान।}

\item \textbf{Question 189:} B

\item \textbf{Question 190:} B

\item \textbf{Question 191:} A
\\ \textit{कोर्स टेक्स्ट में बताया गया है कि Bitcoin बेहद अस्थिर है फाइनेंशियल एसेट, और इसकी कीमत गिरकर 0 हो सकती है, जो एक Bitcoin खरीदने से जुड़े जोखिमों के बारे में}

\item \textbf{Question 192:} C
\\ \textit{कोर्स के टेक्स्ट में यह बताया गया है कि Bitcoin खरीदने से पहले, आपको एक सुरक्षित Wallet होना चाहिए जहाँ तुम अपनी चाबियाँ रखोगे।}

\item \textbf{Question 193:} D
\\ \textit{पाठ्यक्रम पाठ में बताया गया है कि डॉलर लागत औसत तकनीक में नियमित अंतराल पर Bitcoin की छोटी मात्रा खरीदना शामिल है। यह विधि समय के साथ मूल्य में उतार-चढ़ाव को संतुलित करती है और स्वामित्व वाली Bitcoin की मात्रा में निरंतर वृद्धि सुनिश्चित करती है।}

\item \textbf{Question 194:} D
\\ \textit{पाठ्यक्रम पाठ में बताया गया है कि एक सहज खरीद का उपयोग तुरंत Bitcoin के संपर्क में आने के लिए किया जाता है। उदाहरण के लिए, इसका मतलब क्रैश के दौरान खरीदारी करना या बोनस का लाभ उठाना हो सकता है।}

\item \textbf{Question 195:} A
\\ \textit{कोर्स टेक्स्ट में बताया गया है कि "नो योर कस्टमर" (KYC) नियमों के मुताबिक, यूज़र्स को पहचान दिखानी पड़ती है आतंकवाद को वित्त पोषित करने, टैक्स चोरी और काले धन से लड़ना धोखाधड़ी (money laundering के लिए colloquial term)}

\item \textbf{Question 196:} B
\\ \textit{कोर्स के टेक्स्ट में बताया गया है कि कई तरीके हैं बिना KYC के बिटकॉइन हासिल करने के तरीकों में से एक है Bitcoin एटीएम।}

\item \textbf{Question 197:} C
\\ \textit{कोर्स के टेक्स्ट में सलाह दी गई है कि Bitcoin खरीदने के बाद, एक... Exchange प्लेटफॉर्म से UTXOs को तुरंत निकालना होगा। हैकिंग और फंड ब्लॉकिंग के खतरों को कम करने के लिए।}

\item \textbf{Question 198:} B
\\ \textit{इस कोर्स टेक्स्ट में बताया गया है कि जब DCA चुनते समय प्लेटफॉर्म, सही चुनाव इस पर ज़्यादा निर्भर करेगा आपके देश में उपलब्धता।}

\item \textbf{Question 199:} A

\item \textbf{Question 200:} A

\item \textbf{Question 201:} B

\item \textbf{Question 202:} A
\\ \textit{कोर्स के टेक्स्ट में यह बताया गया है कि नियमों का पालन न करने से आपके देश में Bitcoin खरीदते समय आपको जोखिम हो सकता है।}

\item \textbf{Question 203:} D

\item \textbf{Question 204:} C

\item \textbf{Question 205:} C

\item \textbf{Question 206:} A
\\ \textit{Bitcoin को एक बहु-विषयक विषय के रूप में वर्णित किया गया है क्योंकि यह विभिन्न क्षेत्रों से जुड़े सवाल खड़े करता है, जिनमें शामिल हैं... क्रिप्टोग्राफ़ी, गेम थ्योरी, अर्थशास्त्र और मौद्रिक नीति, कंप्यूटर साइंस, फिलॉसफी, एनर्जी, कानून, और रेगुलेशन।}

\item \textbf{Question 207:} A
\\ \textit{इस कोर्स टेक्स्ट में कहा गया है कि Bitcoin का मुख्य लक्ष्य... क्रांति का मकसद एक ऐसी दुनिया बनाना है जहां इंसानी हुकूमत हो। एक सही, प्राइवेसी डिफ़ॉल्ट रूप से सम्मानित होती है, और पैसा नहीं छेड़ा गया।}

\item \textbf{Question 208:} D
\\ \textit{मिल्टन फ्राइडमैन, एक मशहूर अर्थशास्त्री, ने 1999 में भविष्यवाणी की थी कि जल्द ही एक भरोसेमंद ई-कैश विकसित हो जाएगा, एक तरीका जहां इंटरनेट पर आप पैसे A से B में ट्रांसफर कर सकते हैं, बिना A को B के बारे में पता हो या B को A के बारे में।}

\item \textbf{Question 209:} C
\\ \textit{Bitcoin में सामान्य ऊपर की ओर प्रवृत्ति एक निश्चित अवधि में इसकी कीमत में लगातार वृद्धि को संदर्भित करती है। यह ध्यान रखना महत्वपूर्ण है कि जबकि ऊपर की ओर प्रवृत्ति एक सकारात्मक प्रक्षेपवक्र का सुझाव देती है, यह गारंटी नहीं देती है कि कीमत बिना किसी उतार-चढ़ाव के बढ़ती रहेगी।}

\item \textbf{Question 210:} A
\\ \textit{कोर्स के टेक्स्ट में उल्लेख है कि कुछ देशों ने प्रतिबंध लगा दिया है और Bitcoin के इस्तेमाल को अपराध घोषित कर दिया, जिससे और बढ़ गया... Bitcoin अपनाने की जटिलता संस्कृतियों, युगों और... के आधार पर देशों ने। हालांकि, अल सल्वाडोर ने इसे अपनाने की हिम्मत दिखाई। एक कानूनी मुद्रा के रूप में।}

\item \textbf{Question 211:} A
\\ \textit{Bitcoin एक बहु-विषयक विषय है जो उठाता है... विभिन्न क्षेत्रों से संबंधित प्रश्न, जिनमें शामिल हैं क्रिप्टोग्राफी, गेम थ्योरी, अर्थशास्त्र और मौद्रिक नीति, कंप्यूटर साइंस, फिलॉसफी, एनर्जी, कानून, और रेगुलेशन।}

\item \textbf{Question 212:} C
\\ \textit{कोर्स टेक्स्ट में Bitcoin के अपनाने को वायरल बताया गया है, यह सुझाव देता है कि यह एक से तेजी से और व्यापक रूप से फैलता है एक व्यक्ति से दूसरे व्यक्ति में फैलता है, बिल्कुल वायरस की तरह, लेकिन एक अच्छा वाला।}

\item \textbf{Question 213:} D
\\ \textit{कोर्स के टेक्स्ट के मुताबिक, Bitcoin नहीं होने वाला है गायब मत हो। बल्कि, क्रांति अभी शुरू हुई है।}

\item \textbf{Question 214:} C
\\ \textit{कोर्स बताता है कि इतना कुछ है जो एक्सप्लोर करने लायक है Bitcoin कि सब कुछ समझना मुश्किल है एक बार, इसलिए अपना समय लेना चाहिए।}

\item \textbf{Question 215:} C
\\ \textit{Lightning Network एक भुगतान नेटवर्क है जो Bitcoin प्रोटोकॉल तेज़ लेन-देन को सक्षम करने के लिए। यह नहीं है अलग क्रिप्टोकरेंसी नहीं, न ही Mining Bitcoin के लिए कोई तरीका।}

\item \textbf{Question 216:} A

\item \textbf{Question 217:} C
\\ \textit{Blockchain ट्राइलेमा उस चुनौती को संदर्भित करता है जिसमें Blockchain पर आधारित प्रोटोकॉल सिर्फ दो ही मापदंडों को पूरा कर सकता है विकेंद्रीकरण, सुरक्षा, और स्केलेबिलिटी।}

\item \textbf{Question 218:} A

\item \textbf{Question 219:} A
\\ \textit{जीडब्ल्यू-1 को पहली बार मान्य किया गया और लागू किया गया 2017 में।}

\item \textbf{Question 220:} D

\item \textbf{Question 221:} D
\\ \textit{वेस्टर्न यूनियन, केंद्रीय बैंक, वीज़ा और मास्टरकार्ड जैसी पारंपरिक धन हस्तांतरण सेवाएँ गायब हो सकती हैं यदि वे Lightning Network तकनीक को नहीं अपनाते हैं।}

\item \textbf{Question 222:} B

\item \textbf{Question 223:} A

\item \textbf{Question 224:} A
\\ \textit{Lightning Network से आपके बीच लेन-देन संभव बनाता है वे लोग जिनका कोई सीधा चैनल कनेक्शन नहीं है।}

\item \textbf{Question 225:} C
\\ \textit{दो के निर्माण के बीच औसत समय अंतराल Bitcoin प्रोटोकॉल में ब्लॉक्स का समय 10 मिनट है।}

\item \textbf{Question 226:} B

\item \textbf{Question 227:} C

\item \textbf{Question 228:} D

\item \textbf{Question 229:} C
\\ \textit{Lightning Network सतोशी (जिसे Sats भी कहा जाता है) का इस्तेमाल करता है, जीडब्ल्यू-5 का दशमलव, काम करने के लिए।}

\item \textbf{Question 230:} B
\\ \textit{खेल में शामिल होने का मतलब यह है कि हाल ही में वीडियो गेम इंडस्ट्री में काफी लोकप्रिय हुआ है। यह असल में यह डेवलपर्स और हितधारकों का व्यक्तिगत निवेश या दाँव लगा होता है एक गेम की सफलता, जो उनके हितों को संरेखित करती है एक अच्छा प्रोडक्ट पहुंचाना।}

\item \textbf{Question 231:} A

\item \textbf{Question 232:} B

\item \textbf{Question 233:} C
\\ \textit{Lightning Network माइक्रो-ट्रांजैक्शन के साथ, यूजर्स कर सकते हैं... इसमें लेन-देन शुल्क कम होते हैं, जिससे पैसे भेजना सस्ता पड़ता है बहुत ही कम मात्रा में Bitcoin। वे तुरंत भी कर सकते हैं भुगतान, जिससे रियल-टाइम लेनदेन संभव होता है। आखिरकार वे स्केलेबिलिटी मिल सकती है, जिससे बड़ी मात्रा को हैंडल करने में मदद मिलती है एक साथ होने वाले लेन-देन की संख्या। इसलिए, गलत जवाब है लंबा भुगतान समय।}

\item \textbf{Question 234:} B
\\ \textit{Lightning Network के साथ, कंटेंट क्रिएटर्स को अच्छी क्वालिटी का कंटेंट बनाने के लिए प्रोत्साहित किया जाता है ताकि लोग बने रहें उपयोगकर्ताओं का ध्यान। उपयोगकर्ता, बदले में, केवल भुगतान करेंगे वे जो कंटेंट देखते हैं, उसके लिए सीधे सब्सक्रिप्शन की जरूरत नहीं होगी फीस।}

\item \textbf{Question 235:} A
\\ \textit{एक डीसीए प्लेटफॉर्म वह प्लेटफॉर्म है जो सुविधा प्रदान करता है डॉलर-कॉस्ट एवरेजिंग (डीसीए), एक निवेश रणनीति जहां एक व्यक्ति एक निश्चित रकम को निवेश करता है नियमित अंतराल पर विशेष संपत्ति, चाहे जो भी हो एसेट की कीमत।}

\item \textbf{Question 236:} D

\item \textbf{Question 237:} A

\item \textbf{Question 238:} C

\item \textbf{Question 239:} D

\item \textbf{Question 240:} D
\\ \textit{टेक्स्ट बताता है कि Bitcoin, मतलब के लिए इंटरनेट की तरह है संचार का, एक तकनीकी क्रांति है जो नए बड़े पैमाने के तरीकों के लिए मौका बनाता है संगठन, जिसने हमें आदान-प्रदान की संभावना दी बिना किसी भरोसेमंद तीसरे पक्ष के मूल्य।}

\item \textbf{Question 241:} C

\item \textbf{Question 242:} C
\\ \textit{पाठ में उल्लेख है कि Bitcoin छद्मनामी है और हो सकता है दुनिया में कहीं भी बिना अनुमति के विनिमय किया जा सकता है किसी भी संस्था का।}

\item \textbf{Question 243:} A

\item \textbf{Question 244:} D

\item \textbf{Question 245:} C

\item \textbf{Question 246:} C
\\ \textit{पाठ में उल्लेख है कि दुनिया में 2.4 अरब लोग हैं बैंक खाते के बिना दुनिया, और Bitcoin की अनुमति देता है लेन-देन में समानता, चाहे उनका सामाजिक स्तर कुछ भी हो या राजनीतिक पद।}

\item \textbf{Question 247:} C

\item \textbf{Question 248:} B
\\ \textit{टेक्स्ट में उल्लेख है कि यूज़र्स बिटकॉइन्स प्राप्त कर सकते हैं उन्हें व्यापार के लिए स्वीकार करना या खरीदकर नियमित या अनियमित प्लेटफॉर्म।}

\item \textbf{Question 249:} D
\\ \textit{पाठ में उल्लेख है कि Satoshi नाकामोटो ने Bitcoin को बनाया था 2008 में वित्तीय प्रणाली में बदलाव का प्रस्ताव रखा गया।}

\item \textbf{Question 250:} D


\end{enumerate}

\end{document}
